\documentclass[11pt,a4paper]{article}
\usepackage[T1]{fontenc}
\usepackage[utf8]{inputenc}
\usepackage{lmodern,microtype}
\usepackage[margin=.85in]{geometry}
\usepackage{amsmath,amssymb,amsthm,mathtools,aliascnt}
\usepackage{booktabs,tabularx,array,enumitem,needspace,float}
\usepackage[section]{placeins}
\usepackage{xcolor,tikz,mdframed,graphicx}
\usetikzlibrary{arrows.meta,positioning,calc,decorations.pathreplacing}
\usepackage[hidelinks]{hyperref}
\usepackage[nameinlink,noabbrev]{cleveref}
\usepackage{fancyhdr}
\hypersetup{pdftitle={Sync HotStuff RR, Part II: Sequencer Rotation},pdfsubject={Regional rotation with Espresso ordering and explicit L1 multisignature control}}
\setlength{\headheight}{14pt}
\pagestyle{fancy}\fancyhf{}
\lhead{\small Sync HotStuff RR, Part II}\rhead{\small Sequencer Rotation}
\cfoot{\thepage}
\setlength{\parindent}{0pt}\setlength{\parskip}{4pt}
\setlength{\emergencystretch}{2em}
\renewcommand{\topfraction}{.92}
\renewcommand{\textfraction}{.06}
\renewcommand{\floatpagefraction}{.85}
\makeatletter
\setlength{\@fptop}{0pt}\setlength{\@fpbot}{0pt plus 1fil}
\makeatother
\setlist{nosep,leftmargin=*}
\newcolumntype{Y}{>{\raggedright\arraybackslash}X}
\definecolor{rrRed}{HTML}{A51C30}
\definecolor{rrBlue}{HTML}{2563EB}
\definecolor{rrTeal}{HTML}{0F766E}
\newcommand{\RR}[1]{{\color{rrRed}#1}}
\newcommand{\QC}{\mathcal C}
\newcommand{\FC}{\mathcal F}
\newcommand{\Hash}{\mathsf H}
\newcommand{\rank}{\operatorname{rank}}
\newcommand{\TB}{\mathit{TB}}
\newcommand{\tagg}[1]{\ensuremath{\mathsf{#1}}}
\newcommand{\defterm}[1]{\textcolor{rrBlue}{\textit{#1}}}
\newtheorem{definition}{Definition}[section]
\newaliascnt{assumption}{definition}
\newtheorem{assumption}[assumption]{Assumption}
\aliascntresetthe{assumption}
\crefname{assumption}{assumption}{assumptions}
\Crefname{assumption}{Assumption}{Assumptions}
\newaliascnt{proposition}{definition}
\newtheorem{proposition}[proposition]{Proposition}
\aliascntresetthe{proposition}
\crefname{proposition}{proposition}{propositions}
\Crefname{proposition}{Proposition}{Propositions}
\newcounter{protocol}
% Protocol boxes float to the top of a page, so long boxes leave no gaps.
\newenvironment{protocol}[3][htbp]{\begin{figure}[#1]\refstepcounter{protocol}\label{#2}
\begin{mdframed}[linewidth=.6pt,skipabove=0pt,skipbelow=0pt,
innertopmargin=7pt,innerbottommargin=7pt]\small
\textbf{Protocol \theprotocol. #3}\par\smallskip}{\end{mdframed}\end{figure}}
\crefname{protocol}{Protocol}{Protocols}
\Crefname{protocol}{Protocol}{Protocols}
\title{Sync HotStuff RR, Part II: Sequencer Rotation}
\author{}\date{19 October 2026}
\begin{document}
\maketitle
\begin{abstract}
Part~I runs one synchronous BFT engine inside one region. This document
lets several regions take turns producing one chain. Each regional turn
is a \emph{session} with a fixed DATA quota. When a session commits its
last DATA, the next eligible region takes over at once, without waiting
for global confirmation. A Publish Layer sends complete regional results
to Espresso. Every verifier reads the same finalized L1 controls and
Espresso history, and so accepts the same continuous prefix of DATA,
called the global prefix~$G$.

A multisignature committee can force a handoff. It does so with an
atomic L1 branch update that takes effect at a future Espresso time and
carries the complete region map for the new branch.
A Force closes the current production branch and restarts from $G$.
It never removes DATA that is already in $G$. Missing output alone never
causes a Force. On each node, a Session Manager prepares, enters and
retires the local programs of different sessions. The document states
which guarantees hold unconditionally and which depend on regional
service, publication service and committee behavior.
\end{abstract}

\section{Overview}
\label{sec:overview}
Part~I~\cite{part1} gives one region a synchronous BFT engine.
This document connects several regional engines into one chain.

\paragraph{A running example.}
We use one example throughout. The regions are A, B, C and D, in this
cyclic order.
\begin{itemize}
\item A's turn produces DATA 101--110. A commits all of them inside its
region and hands off to B.
\item B continues with DATA 111--115.
\item Meanwhile, verifiers have globally accepted only DATA up to 108.
\end{itemize}
Now the multisignature committee orders a Force. When the order takes
effect:
\begin{itemize}
\item everything up to 108 stays in the chain;
\item the unconfirmed DATA 109--115 of the old branch can no longer be
accepted;
\item if B is selected again, B starts a \emph{new} program that produces
a new DATA 109 after 108;
\item B's old program, which continued after 110, shuts down separately.
The two programs never share votes, locks or timers.
\end{itemize}
The rest of the document explains when a handoff happens, who chooses
the next region, why the new program continues after 108, and how old
and new programs run side by side.

\subsection{How the protocol works}
\paragraph{Two kinds of handoff.}
A \defterm{regular handoff} happens when a session finishes its quota.
A commits its last DATA, collects a regional commitment receipt, and
gives the result to B. B checks the data, its own selection and the
execution state, and then continues, possibly before A's last
DATA is globally accepted.

A \defterm{Force} is a handoff ordered by the multisignature committee.
The committee submits an L1 instruction that names a future Espresso
time, the target branch version and one complete new status map. When
verifiers reach the first Espresso block at or after that time, they all
apply the same instruction if its target is current: they close the old
branch, install the new map and continue from the current global prefix.
If the target is stale, the whole instruction is a no-op. A Force happens only
through such an instruction. Silence of a region, a monitoring alarm or
a local timeout does not cause one. The committee may use slow delivery
as a reason to act, but that judgment belongs to its own policy.

\paragraph{Three stages of a block.}
A DATA block passes through three separate stages.
\begin{enumerate}
\item \emph{Regional commitment.} The producing region commits the block
under its BFT rules and collects enough FINAL signatures.
\item \emph{Publication.} The Publish Layer sends the block and its
evidence to Espresso, which fixes its position in a finalized log.
\item \emph{Global acceptance.} Verifiers check the regional evidence,
the L1 controls and the rotation rules. If the block extends the
current global prefix, they append it.
\end{enumerate}
Only the third stage adds a block to the \defterm{global prefix} $G$,
and $G_q$ is its value at Espresso height $q$. We write ``$G=108$'' for ``the global prefix ends at
DATA 108''.

The \defterm{production branch} $R$ is what regions have produced so
far. It can run ahead of $G$: in the example, $R$ reaches 115 while
$G$ is still 108. The \defterm{export session} is the session that owns
the next DATA height after $G$. In the example it is still A's
session, because 109 is A's block, even though B is the region that is
currently producing. A Force closes the part of $R$ beyond $G$; it
never touches $G$.

A later stage, the rollup DA publication, has its own accepted prefix.
It may lag behind $G$, but it may not skip any DATA
in $G$. The Publish Layer of this document, which delivers regional
results to Espresso, is a different service.

A regional receipt remains valid evidence of regional commitment even
if the branch containing the block is later closed. Global acceptance is
stronger: no Force ever removes DATA from $G$. These statements concern
input content, order and provenance. Correct execution and native state
continuity belong to the application and publication interfaces.

\begin{table}[htbp]\centering\small
\caption{Participants and responsibilities.}\label{tab:roles}
\begin{tabularx}{\linewidth}{@{}>{\raggedright\arraybackslash}p{.19\linewidth}Y@{}}\toprule
Role & Responsibility\\\midrule
L1 & Publishes immutable authorities and the ordered control log.\\
Regional replica & Runs the DATA/DRAIN engine of its sessions and
serves complete commitment evidence.\\
Multisignature service & Announces future atomic branch updates on L1.
Each update supplies the complete map for the new branch and causes a
Force if its target is current. It cannot choose a retained prefix or
declare data global.\\
Publish Layer & Delivers complete regional results to Espresso. The
leader, replicas and external verifiers may act as senders.\\
Espresso & Supplies a finalized ordered log, authenticated timestamps,
complete namespace contents, and bounded finalization and visibility.\\
Observer / certifier & Verifies both sources, replays control and
exports, and derives global output and forced boundaries.\\
\bottomrule
\end{tabularx}
\end{table}

\subsection{The four view indices}
The full view is
\[
 (\beta,r,g,t).
\]
\begin{itemize}
\item $\beta$ is the \defterm{branch version}. Each effective Force
increases it by one. It is not an Espresso epoch or a time interval; it
only separates the old production branch from the new one.
\item $r$ is the \defterm{turn index}. It selects the immutable L1
authority (voting snapshot and DATA quota) used by a regional turn. It
is not a physical region identifier.
\item $g$ is the \defterm{generation}: how many times the session has
run regional recovery.
\item $t$ is the \defterm{term}: which leader turn within the
generation.
\end{itemize}
A \defterm{session} is the pair $(\beta,r)$: one region's production
turn on one branch version. Its voting snapshot and DATA quota stay
fixed through all of its recoveries. Every node and certifier derives
these numbers by replaying the same authenticated history, so no other
session identifier is needed.

\begin{table}[htbp]\centering\small
\caption{How the view changes.}\label{tab:views}
\begin{tabularx}{\linewidth}{@{}>{\raggedright\arraybackslash}p{.19\linewidth}>{\raggedright\arraybackslash}p{.25\linewidth}Y@{}}\toprule
Change & New view & Trigger and effect on timers\\\midrule
Leader rotation & $(\beta,r,g,t+1)$ & A terminal QC below $H_{\beta,r}$;
running commit timers continue.\\
Regional recovery & $(\beta,r,g+1,0)$ & A quit proof, the $\Delta_i$ wait
and finalized authority; the quit cancels generation-$g$ timers.\\
Regular handoff & $(\beta,r+1,0,0)$ & Ready on a REGULAR witness
with the terminal receipt; the old engine stops through Completion.\\
Force & $(\beta+1,r',0,0)$ & $r'=r_c$ if the export session has zero
accepted DATA, otherwise $r_c+1$; an unassigned boundary keeps its
target index. Old programs close separately.\\
\bottomrule\end{tabularx}
\end{table}

For example, $(7,20,0,0)\to(7,20,0,1)$ is a leader rotation inside A,
$(7,20,0,1)\to(7,20,1,0)$ is a recovery inside A, and
$(7,20,1,0)\to(7,21,0,0)$ is a regular handoff to B. None of these
changes $\beta$.

A Force can reuse a turn index. Suppose export session $(7,20)$ has
already added DATA to $G$. A Force then opens $(8,21,0,0)$, even if
a regular successor already used $(7,21)$. If $(7,20)$ had added no
DATA to $G$, the Force would open $(8,20,0,0)$ instead. In both cases
the new pair differs from every old pair, so votes, locks and evidence
stay separate. Where the new session starts is fixed by the DATA anchor
it keeps, not by its turn index. An L1 control that is finalized but
not yet effective does not stop production or global acceptance.

\subsection{What the protocol leaves to others}
\label{sec:scope}
The protocol fixes the following boundaries.
\begin{itemize}
\item Only a finalized, authorized MSIG instruction can cause a
cross-region Force. A Force needs no waiting period, no Ready and no
prior output.
\item Region eligibility changes only when an atomic L1 branch update
successfully installs the complete map of a new branch. The protocol
never disables or restores a region by itself.
\item Deciding \emph{when} to intervene is the multisignature
committee's policy. \Cref{app:policy} sketches a reference policy;
replay never depends on it.
\item The protocol does not enforce fair service among regions or
eventual inclusion of a client request. Request persistence and
resubmission belong to the execution layer.
\end{itemize}
Regional BFT timers, BLAME and Quit remain in place. Inside a region,
the only changes from Part~I are those listed in \cref{sec:protocols}.

\section{Model and assumptions}
\label{sec:model}
\subsection{L1 authorities and atomic branch updates}
\label{sec:l1}
\label{sec:runtime-controls}
Nodes authenticate finalized L1 headers, authorities and a complete,
ordered control history. L1 timestamps only label L1 blocks. They are
neither regional commit times nor Espresso activation times.

\begin{definition}[\defterm{Atomic branch update}]
\label{def:control-update}
An atomic branch update is
\[
 U=(\mathit{chain},c,\theta_a,\beta_{\rm target},S_{\rm new})
\]
together with a signature set $\Sigma_U$ on its hash. Here $c$ is its
sequence number, $\theta_a$ its Espresso activation threshold,
$\beta_{\rm target}$ the branch version it asks to close, and
\[
 S_{\rm new}:\mathcal L_R\longrightarrow
 \{\tagg{ON},\tagg{OFF}\}
\]
is one explicit, complete status map; it may set every region OFF, but
the map field itself may not be omitted. The update names no retained
DATA height. There are no status-only updates and no separate Force tag.
\end{definition}

\begin{definition}[\defterm{Valid atomic branch update}]
\label{def:control-valid}
An update $U$ in finalized L1 history is valid when:
\begin{enumerate}[label=(V\arabic*)]
\item $\mathit{chain}$ is this chain;
\item $\Sigma_U$ satisfies the control authority $\mathcal K$ fixed by
the genesis checkpoint;
\item $c$ is larger than the sequence number of every earlier valid
update in L1 order;
\item $\beta_{\rm target}\in\mathbb N$ and $S_{\rm new}$ assigns exactly
one flag to every region of $\mathcal L_R$; and
\item $\theta_a\ge\theta_a(U')+T_{\rm update}$, where $U'$ is the
previous valid update, or the last update of the genesis prefix.
\end{enumerate}
\end{definition}

The \defterm{control log} $\Lambda=\langle U_1,U_2,\ldots\rangle$ is the
sequence of valid updates in L1 order. A verifier processes every entry
exactly once, including stale entries, and keeps the last processed
sequence and control-prefix commitment as its control cursor $\kappa$.

\begin{table}[htbp]\centering\small
\caption{L1 authority and control updates.}\label{tab:l1}
\begin{tabularx}{\linewidth}{@{}p{.27\linewidth}Y@{}}\toprule
Field & Meaning\\\midrule
$\tagg{AUTH}_{r,0}$ & Fixed regional snapshots, leader schedule and
DATA quota $\TB_r$ for index $r$.\\
$\tagg{AUTH}_{r,g}$, $g>0$ & Recovery leader schedule; it does not
change the initial snapshot or quota.\\
Branch update & One L1 record
$(\mathit{chain},c,\theta_a,\beta_{\rm target},S_{\rm new})$
(\cref{def:control-update}).\\
Control sequence $c$ & Unique increasing index binding the chain and the
update.\\
Activation time $\theta_a$ & Espresso timestamp threshold, announced
before its first crossing can become final. One per atomic update;
strictly increasing along the log, at least
$T_{\rm update}$ apart.\\
Target $\beta_{\rm target}$ & The branch version that the update asks to
close. A stale target makes the entire update a no-op.\\
Status map $S_{\rm new}$ & A complete ON/OFF map for every region. If
the update executes, this is the immutable map of the new branch
version.\\\bottomrule
\end{tabularx}
\end{table}

\paragraph{Authority records.}
For each $(r,g)$, L1 supplies at most one immutable authority.
$\tagg{AUTH}_{r,0}$ fixes $\mathit{cfg}_r$ and $\TB_r$. The configuration
$\mathit{cfg}_r$ contains the members, weights, $F_i$, $\Delta_i$ and $K$
of each region. For $g>0$, only the leader schedule changes. Nodes wait
until the evidence they need arrives. Later controls and slashing never
reweight an old certificate or reset a session quota.

Authority is indexed by $r$ alone, not by branch. A new branch version
that reuses $r$ reads the same authority. Its pair $(\beta,r)$ still
names a different engine. An authority, once assigned, is never
rebound.

\paragraph{Ingesting controls.}
The L1 control history is one log of atomic branch updates, with sequence
numbers and a binding history commitment. A verifier keeps both pending
and applied updates. Semantically, however, replay carries forward only
the not-yet-activated queue, the cursor $\kappa$ and the current
$S_\beta$; Protocol~3 never rescans applied entries to select a region.
Applied bytes remain available as audit and closure evidence. Finalization
on L1 only announces an update; it takes effect at its activation time.
Activation times strictly increase
along the log, at least $T_{\rm update}$ apart (\cref{def:control-valid}),
so entries are processed in sequence order (\cref{def:control-effect}). Malformed, unauthorized or duplicate updates
grant nothing (\cref{def:control-valid}). To change an instruction, the
committee submits a later update; finalized history is never edited.

An executing update atomically closes its target branch, increments the
branch version and installs $S_{\rm new}$ as the new version's immutable
status map.  If its target is stale, the entire update is a no-op: its
map does not apply and its operational intent is not carried forward.
The committee must resubmit that intent against the current version.

\paragraph{Espresso positions and when controls apply.}
Let $J_q$ be Espresso block $q$, with authenticated timestamp $\theta_q$.
Within a block, namespace-transaction and record indices give a complete
order, also when timestamps are equal. For a threshold $d$, let
\[
 q(d)=\min\{q:\theta_q\ge d\}
\]
be the first block whose timestamp reaches $d$. For example, if blocks
500, 501, 502 have timestamps 119, 122, 125, then $q(120)=501$.

A control with threshold $d$ takes effect at the start of $J_{q(d)}$,
before any EXPORT in that block is examined. We call this start-of-block
step the \defterm{prelude}.  It processes all due entries in sequence
order but executes at most one branch update: the first due entry whose
target equals the branch version $\beta_0$ at the start of the prelude.
A version created in this prelude is never targeted again in the same
prelude.  Every other due entry is stale and has no state effect.  There
is no other path to a cross-region Force. We write ``$G$ at the start of
$J_q$'' for the global prefix just before the block's EXPORTs.

\begin{definition}[\defterm{Atomic update effect}]
\label{def:control-effect}
Let $\Lambda_q$ be the valid updates with $q(\theta_a)=q$, in sequence
order, and let $\beta_0$ be the branch version at the start of $J_q$.
If no entry of $\Lambda_q$ targets $\beta_0$, the prelude changes no
protocol state. Otherwise let
\[
 U^*=(\mathit{chain},c^*,\theta_a^*,\beta_0,S^*)
\]
be the first such entry.  The prelude atomically closes $\beta_0$, sets
$\beta:=\beta_0+1$, fixes
\[
 S_{\beta}:=S^*,
\]
and executes the Force transition of \cref{alg:global-rule}.  All other
entries of $\Lambda_q$ are stale whole-update no-ops.  Every entry,
effective or stale, advances the processed control cursor exactly once.
\end{definition}

Genesis fixes $S_0$.  The map $S_\beta$ is immutable for the lifetime of
branch version $\beta$; regular handoffs change $r$ but continue to use
the same map.  A region is eligible for turn index $r$ exactly when
$S_\beta$ sets it ON and $\tagg{AUTH}_{r,0}$ gives it positive total
weight. A monitoring report, missing output or local timeout changes
nothing.  If every region is ineligible, the next session stays pending
until another atomic branch update creates a later version with a new
complete map.  No update removes accepted DATA.

\paragraph{Choosing the next region.}
$\tagg{NextLive}(i,S_\beta;r)$ searches the other eligible regions after
$i$ in the fixed cyclic order $\mathcal L_R$, using
$\tagg{AUTH}_{r,0}$ for the target turn.  It returns $i$ itself only if
no other region is eligible and $i$ still is; otherwise it returns
$\bot$.  A synthetic source $\bot$ searches from the beginning.  Because
$S_\beta$ is fixed, every regular handoff in version $\beta$ obtains the
same result regardless of when its terminal receipt is delivered.

\paragraph{Genesis.}
The system starts from an initial checkpoint. It fixes the chain
identifier, the cyclic region order $\mathcal L_R$, the genesis DATA
anchor, the first replay position, the Espresso quorum trust anchor,
the control authority $\mathcal K$, the minimum update gap
$T_{\rm update}$, the control prefix, the complete initial map $S_0$,
and the initial branch version $\beta=0$. The initial
authority and state are either trusted or verified from their stated
sources; starting this replay does not make Espresso authenticate
earlier application history.

The genesis anchor carries synthetic commitment evidence and fixed,
authenticated L1 and Espresso references. At replay position
$q_{\rm init}$, a synthetic boundary with no source region and target
index 0 applies $\tagg{NextLive}(\bot,S_0;0)$. If it returns a region,
session $(0,0)$ opens with $q^{\rm open}_{0,0}=q_{\rm init}$; otherwise
the same boundary stays pending. A later successful atomic update closes
the pending version, installs a new complete map and tries the same target
index on the new version. In either case,
$\lambda^{\rm open}_{\beta,0}=\max\{\lambda(\text{genesis}),
\lambda(\tagg{AUTH}_{0,0})\}$, where the second term is the L1 height
at which the authority was published. A later checkpoint used as a
starting point must instead include its current boundary, its reference
bounds and the complete replay state.

\subsection{Data delivery and the Publish Layer}
\label{sec:delivery}
The \defterm{Publish Layer} is a delivery service run by existing
participants. It runs no consensus of its own and does not decide which
DATA is valid. A regional leader, another replica or an external
verifier can submit. The rotation rules rely only on the service
guarantee in \cref{ass:services}(i); \cref{sec:publish-implementation}
describes one implementation.

A \defterm{complete regional result} is a DATA block together with its
direct or covering FINAL receipt and the ancestry, boundaries and source
evidence needed to validate it. A proposal QC, a block hash, or a receipt
without its payload is not complete. Correct replicas offer complete
results to the Publish Layer as soon as they have them. The service may
batch results and builds EXPORT records as specified in
\cref{sec:exports}.

Regions may also send results directly to other regions. This lets a
regular successor become Ready before the terminal result is finalized
in Espresso. Without direct delivery, the successor reads the result
from Espresso after finalization, which takes longer. There is no
separate cross-region transport service.

Espresso carries finalized records, including DATA and commitment
evidence, to online readers. Readers still check the application and
rotation rules. L1 supplies its own authority and control proofs. An
execution-state snapshot may be transferred directly or rebuilt from
authenticated input history. Rebuilding state is part of Ready,
not of Espresso inclusion.

\subsection{Stake and faults}
\label{sec:stake}
Member $v$ of region $i$ has voting weight $\omega_{i,v}>0$. Let
$W_i=\sum_{v\in V_i}\omega_{i,v}$. The stake configuration fixes a bound
$F_i<W_i/2$, and Byzantine members of region $i$ hold total weight at
most $F_i$. Every regional certificate needs matching signatures from
distinct members with total weight greater than $F_i$. So every
certificate has a correct signer, and the correct members alone can form
one. This replaces the $f+1$ of $n=2f+1$ in Part~I.

Membership and weights stay fixed for a whole session, including its
recoveries. Slashing and later atomic branch updates do not change them. A later
session checks old certificates against the old snapshot. The bound
$F_i$ counts every key of the snapshot that is compromised while its
certificates are still accepted. Compromise of old keys after that is
outside this model.

Correct replicas keep their protocol state and committed blocks.
Signatures cannot be forged, and hashes are binding.

\subsection{Assumptions}
\label{sec:assumptions}
\begin{assumption}[Finalized L1 records]
\label{ass:authority}
At every finalized L1 height, the authority records and the ordered
control history through that height are common and authenticated.
Finalized records never change. A verifier obtains a complete control
prefix, not an arbitrary subset of valid events.
\end{assumption}

\begin{assumption}[Authenticated synchronous communication and global clock]
\label{ass:network}
\leavevmode\par
\begin{itemize}
\item Messages use authenticated communication channels.
\item Within region $i$, the actual delay of a correct-to-correct
message is $\delta\le\Delta_i$; $\delta$ may differ across messages.
\item All nodes share a global clock.
\end{itemize}
\end{assumption}
The delay bound covers delivery and processing of complete validation
dependencies. Each replica records its own entry and voting events on
the global clock.

\begin{assumption}[Honest multisignature committee]
\label{ass:committee}
The committee acts as one honest, coordinated control service. It
submits only authorized, immutable, future-dated atomic branch updates, with
activation times that satisfy the advance-announcement condition below. It keeps no private protected prefix and cannot select
an older retained anchor. Honesty here imposes no response deadline; the
minimum gap $T_{\rm update}$ between updates is a validity rule
(\cref{def:control-valid}), not part of honesty.  The reference policy
normally keeps one unactivated update in flight and resubmits any stale
operational intent against the current branch version.
\end{assumption}

\Needspace{8\baselineskip}
\begin{assumption}[Espresso history and time]
\label{ass:espresso}
Finalized Espresso history is consistent, persistent and publicly
verifiable, including the complete namespace content of each block, so a
verifier can also confirm that a block holds no record for this chain.
Each block carries a non-decreasing timestamp that runs at most
$\epsilon_E\ge0$ ahead of real time: if $J_q$ first becomes final at real
time $t_E(q)$, then $\theta_q\le t_E(q)+\epsilon_E$.
\end{assumption}
Long-term retrieval is a separate service, for example archival readers;
a finality certificate alone does not provide it. A returning node
catches up, or starts from an authorized checkpoint, before it confirms
new output; catching up has no fixed time bound.

\begin{assumption}[Bounded external services]
\label{ass:services}
There are fixed bounds $T_{\rm pub},T_{\rm L1},T_E>0$, valid under the
declared input-size, admission and load limits, such that:
\begin{enumerate}[label=(\roman*)]
\item \emph{Publish Layer.} If a correct participant has a complete
regional result available at real time $s$, an Espresso-admissible record
carrying that result, with its DATA and commitment evidence intact, is
submitted to Espresso by $s+T_{\rm pub}$. Submission does not rely on the
original leader alone: a working submission path and the required fee
resources exist.
\item \emph{L1.} A valid control operation actually submitted at real
time $s$ is finalized by $s+T_{\rm L1}$, and by then every correct online
participant can retrieve and verify it.
\item \emph{Espresso.} An admissible complete record submitted at real
time $s_E$ is in a finalized block by $s_E+T_E$, and by then every correct
participant that stays online and keeps up with the history can retrieve
the complete record and verify its inclusion and finality. Every finalized
block, with its complete relevant namespace, becomes retrievable and
verifiable to those participants within $T_E$ of finalization, also when
it holds no record for this chain.
\end{enumerate}
\end{assumption}
Participants are regional replicas, external observers and certifiers.
Each bound covers the whole path: inclusion, finalization, dissemination,
retrieval and validation. Each clock starts at an actual event: in (i),
when the complete result is available, not when a leader proposes the
block or one replica commits it; in (ii) and (iii), at actual submission,
which excludes any detection or decision time before it. ``Submitted''
means the complete record reaches the submission interface under its
admission rules; a local send attempt is not enough. A result already
recorded in Espresso satisfies (i). Forming the FINAL receipt is the
regional protocol's job.

These are delivery bounds only. Espresso delivers full payloads and
evidence, but does not promise that an EXPORT passes this chain's
acceptance rules: a stale base, a closed branch or a missing predecessor
can still stop it from advancing $G$, and retries needed for acceptance
are a separate task. The bounds never trigger a Force, and $T_E$ does not
replace the clock-lead bound of \cref{ass:espresso} or the
advance-announcement rule below.

For a newly available complete result, (i) and (iii) give
\[
 s_{\rm submit}\le s+T_{\rm pub},\qquad
 s_{\rm visible}\le s+T_{\rm pub}+T_E,
\]
where $s_{\rm visible}$ is the time by which the result and its Espresso
evidence are available to the readers in (iii). This bounds neither when
the regional result first forms, nor when a new region restores its
execution state.

\paragraph{Advance announcement.}
A control must be visible to everyone before it can take effect.
For a control submitted at real time $s$, the committee chooses
\[
 \theta_a>s+T_{\rm L1}+\epsilon_E.
\]
Then the control is available to correct participants before the first
block crossing $\theta_a$ can become final. This is a model condition,
not a conclusion from average L1 or Espresso block times. The L1 bound
already includes dissemination, so no second announcement delay is
added. $T_{\rm L1}$ starts at actual submission; detection, deliberation
and signature collection before submission are not covered by any
bound. Offline or unsynchronized participants catch up before they
confirm new output. Missing control history pauses verification.

\paragraph{Verification interfaces.}
We assume that regional receipt validity and the application-validity
checks below are sound under the regional model. A deployed certifier
must authenticate both source histories and bind the control version it
used. A signature or a TEE alone does not show that an old control
snapshot is complete. The publication interface additionally requires
sound execution and DA correspondence checks.

The bounded preparation assumption for a newly selected region is
stated with the Ready rules, in \cref{ass:ready-service}. The conditions
needed for progress are collected in \cref{sec:progress}.

\section{Regional engine: definitions}
\label{sec:regional}
This section defines the objects that one regional engine uses: blocks
and their positions, the session quota, certificates and receipts, and
the rules for admitting views. The engine is Part~I's engine with
sessions, a quota, DRAIN blocks and FINAL receipts added.
\Cref{tab:notation} lists the main symbols. Anchors and entry times
carry the same four indices as the view.
\begin{table}[htbp]\centering\small
\caption{Notation.}\label{tab:notation}
\begin{tabularx}{\linewidth}{@{}p{.32\linewidth}Y@{}}\toprule
Symbol & Meaning\\\midrule
$(\beta,r,g,t)$; $j$ & Branch version, turn index, generation, term; fresh position.\\
$(\beta,r)$ & Session: one region\textquotesingle s turn on one branch version.\\
$i_{\beta,r}$; $V_i$ & Session region; its regional members.\\
$\omega_{i,v},W_i,F_i$ & Member weight, total weight, and Byzantine bound $F_i<W_i/2$.\\
$\Delta_i$; $K$ & Regional synchronous delay bound; maximum DATA per term.\\
$\tagg{AUTH}_{r,g}$; $\mathit{cfg}_r$ & Immutable authority; fixed session snapshot.\\
$A_{\beta,r,g,t}$ & Anchor: the block just before the first fresh position of that view.\\
$h(B)$; $k_{\beta,r}(B)$ & DATA height; physical position in the session, counting DRAIN.\\
$\TB_r,H_{\beta,r},T_{\beta,r}$ & Session DATA quota, last DATA height, and terminal DATA.\\
$m_{\beta,r,g,t}$; $L(\beta,r,g,t)$ & Number of fresh positions in the term; designated leader.\\
$\QC_{\beta,r,g,t}(B)$; ANCHOR QC & Proposal certificate; re-certification of the opening.\\
$\FC_{\beta,r}(B)$ & Regional FINAL receipt, or its same-session coverage.\\
$Q_{\beta,r}^0,Q_v,Q^*$ & Inherited commitment evidence; entry lock; selected justification.\\
$e_v(\beta,r,g,t),\tau_v(P),a_v$ & Real-time entry, vote and quit events.\\
$\rho,\lambda(B)$ & Finalized L1 reference and its L1 height.\\
$\xi(B)$ & Signed Espresso header reference.\\
$J_q,\theta_q$ & Espresso block and authenticated timestamp.\\
$S_\beta$; $\mathcal L_R$ & Immutable status map of branch version $\beta$; cyclic region order.\\
$q_{\beta,r}^{\rm open}$ & The session's lower bound for Espresso references; not a successor-selection time.\\
$\lambda_{\beta,r}^{\rm open}$ & The session's lower bound for L1 references.\\
$G$, $G_q$; $R$ & Global DATA prefix, and its value at Espresso height $q$; production branch.\\
$T_{\rm L1},\epsilon_E$ & L1 announcement-service bound; maximum Espresso clock lead.\\
$T_{\rm pub},T_E$ & Complete-result submission bound; Espresso finalization and visibility bound.\\
$T_{\rm ready}$ & Preparation bound, measured from the public opening position.\\
$\nu=(\beta,r,i,A,H,q^{\rm open},\mathit{cause})$ & Complete descriptor of the session responsible for the next DATA.\\
$\sigma_G(B)$ & Session record $\nu$ bound to an accepted non-genesis DATA block $B\in G$.\\
$\mathsf{Tail}$; $\omega$ & Tail $\tagg{NEXT}(\nu)$ or $\tagg{PENDING}(\omega)$; boundary with no selected region.\\
$a(G,\nu)$ & Number of DATA of session $\operatorname{sid}(\nu)$ already accepted into $G$.\\
$U=(\mathit{chain},c,\theta_a,\beta_{\rm target},S_{\rm new})$; $\mathcal K$ & Atomic branch update; control authority that signs it.\\
$S_\beta$; $\Lambda$; $\Lambda_q$; $T_{\rm update}$ & Immutable map of branch $\beta$; control log; updates due at $J_q$; minimum activation gap.\\
$\mathcal X_q$; $\kappa_q$ & Maintained replay state; processed control cursor.\\
$\mathit{st}(B)$ & Executable state after block $B$.\\
$\mathcal V_R,\mathcal V_G,\mathcal V_L$ & The Session Manager's regional, global and control views.\\
$E$; Ready & Regional entry proof; the Ready condition of \cref{def:ready}.\\
\bottomrule\end{tabularx}
\end{table}

\subsection{Blocks, positions, and L1 references}
\label{sec:blocks}
The \defterm{anchor} $A_{\beta,r,g,t}$ is the block just before the first
fresh position of view $(\beta,r,g,t)$. In short: the anchor is the block
the view continues from. We always write all four indices.
\begin{itemize}
\item $A_{\beta,r,0,0}$ is the DATA block inherited from the preceding
session; the first session inherits genesis. It keeps the regional
commitment evidence of the session that produced it: a direct or
covering receipt (\cref{sec:coverage}). Genesis has synthetic evidence.
\item For $g>0$, NEWVIEW chooses the \defterm{recovery anchor}
$A_{\beta,r,g,0}$ (\cref{alg:recovery}). It is certified but may be
uncommitted.
\item For $t>0$, $A_{\beta,r,g,t}$ is the terminal DATA block of term
$(\beta,r,g,t-1)$, certified by that term's terminal QC.
\end{itemize}
The anchor of a view never changes. A generation starts at
$A_{\beta,r,g,0}$, and its later terms use their own anchors. ANCHOR
voting re-certifies only $A_{\beta,r,g,0}$, once per generation.

In the running example, A's session starts after anchor 100, and B's
regular session starts after anchor 110. After the Force, B's new session
starts after anchor 108. The inherited anchor keeps its original
producing session, height and hash: DATA 108 produced by $(7,20)$ stays
a $(7,20)$ block even when $(8,21)$ continues from it.

Each block has two positions.
\begin{itemize}
\item The \defterm{DATA height} $h(B)$ counts the DATA blocks on the chain
from genesis to $B$.
\item The \defterm{physical position} $k_{\beta,r}(B)$ counts DATA and DRAIN
blocks within session $(\beta,r)$, with $k_{\beta,r}(A_{\beta,r,0,0})=0$.
This local count does not change the signed fields of $A_{\beta,r,0,0}$.
\end{itemize}
A new block of session $(\beta,r)$ has the form
\[
 B=(\beta,r,k,h,\mathit{type},x,\Hash(B_{\rm parent}),\rho,\xi),
 \qquad \mathit{type}\in\{\tagg{DATA},\tagg{DRAIN}\},
\]
with $k=k_{\beta,r}(B_{\rm parent})+1$. A \defterm{DATA block} carries a
batch $x$ of requests and has $h=h(B_{\rm parent})+1$. A \defterm{DRAIN
block} carries no batch and has $h=h(B_{\rm parent})$. So DATA advances
both positions, and DRAIN advances only the physical position.

\paragraph{L1 references.}
The \defterm{reference} $\rho$ names a finalized L1 block. It holds that
block's height $\lambda(B)$ and hash, and a registry proof for the L1
control records up to that height. It tells validators which fixed
authorities and rules apply. These rules hold:
\begin{itemize}
\item along DATA ancestry, $\lambda$ never decreases;
\item the first DATA block of session $(\beta,r)$ has $\lambda\ge
\lambda^{\rm open}_{\beta,r}$, the L1 lower bound fixed by its boundary
(\cref{sec:witness}; genesis uses \cref{sec:l1});
\item a DRAIN block copies the reference of $T_{\beta,r}$.
\end{itemize}
The timestamp of the L1 block named by $\rho$ is readable from its
header, but it is not a commit time for $B$: recovery may commit $B$
much later (\cref{sec:late-receipt}). The Vote step applies no
receive-time freshness filter to $\rho$.

The L1 reference serves regional validation only. It does not choose
the time of a Force or the Espresso position used for selection.
Regional validity uses the fixed session authority, so a later status
update does not change an old receipt's voting weights.

\paragraph{Espresso references.}
The outer field $\xi(B)$ names a finalized Espresso header $J_q$ by
height and hash. Its height does not decrease along DATA ancestry, and
the first DATA of session $(\beta,r)$ has
$q\ge q_{\beta,r}^{\rm open}$. Terminal DATA uses the same rule as other
DATA. DRAIN copies the Espresso reference of $T_{\beta,r}$, just as it
copies its L1 reference. The proposal hash and signatures bind these
references and the session context. They do not change the native
mini-block encoding. Eligibility is not read from $J_q$: every session
in branch version $\beta$ uses the fixed map $S_\beta$.

Regional full validation authenticates the referenced history, directly
or from its validated cache. Public \tagg{CheckFINAL} establishes the
intrinsic regional validity of the reference. Global acceptance
separately authenticates the required Espresso context and requires the
reference to precede the Espresso block that contains the EXPORT. So no
reference can justify itself through that same EXPORT. Historical
references make no claim about freshness at receipt time. The reference
of $T_{\beta,r}$ contributes only to the regular successor's reference
floor in \cref{sec:witness}.

\paragraph{Why the Espresso reference is needed.}
The terminal reference does not choose the next region. That choice is
already determined by the immutable map $S_\beta$, so two nodes that
receive the terminal receipt at different times compute the same
successor. Instead, the terminal reference gives the successor a common
lower bound for the Espresso headers that its first DATA may cite. It is
not proof that the terminal DATA is globally final.

For example, let the region order be A, B, C and let $S_7$ set all three
regions ON. If A's terminal DATA 110 binds
$\xi(110)=(500,\Hash(J_{500}))$, every verifier selects B under $S_7$,
whether it receives the terminal receipt before or after $J_{505}$.
Changing B to OFF requires an atomic update that closes version 7 and
creates a new version with a new map; it cannot retroactively change the
regular boundary inside version 7.

The next regular session uses the greater of two heights as its reference
floor:
\[
 q^{\rm open}_{\beta,r+1}
 =\max\{\operatorname{height}(\xi(T_{\beta,r})),
         q^{\rm open}_{\beta,r}+1\}.
\]
This advances successive floors without delaying successor selection.
The number is a lower bound, not necessarily the height of an already
available header and not the time at which the successor is selected.
The successor may be assigned and enter from the direct terminal receipt;
its first DATA must wait until it can cite a finalized header at or above
the floor. DRAIN keeps the terminal DATA's reference.

\subsection{Session quota, term positions, and DRAIN}
\label{sec:quota}
\paragraph{1. Fix the session endpoint.}
Session $(\beta,r)$ may add $\TB_r\ge1$ DATA blocks after
$A_{\beta,r,0,0}$. Its \defterm{last DATA height} is
\[
 H_{\beta,r}=h(A_{\beta,r,0,0})+\TB_r.
\]
The DATA block at this height is the \defterm{terminal DATA}
$T_{\beta,r}$. Recovery never changes $H_{\beta,r}$. The quota counts the
DATA blocks after $A_{\beta,r,0,0}$ in a proposal's ancestry, committed
or not. DRAIN, resends and ANCHOR votes use no quota.

\paragraph{2. Produce DATA while space remains.}
If $h(A_{\beta,r,g,0})<H_{\beta,r}$, generation $(\beta,r,g)$ is a
\defterm{DATA generation}. Each term begins at its anchor
$A_{\beta,r,g,t}$ and has
\[
 m_{\beta,r,g,t}=\min\{K,\;H_{\beta,r}-h(A_{\beta,r,g,t})\}
\]
new positions, numbered $j=0,\ldots,m_{\beta,r,g,t}-1$: at most $K$ per
term, and never more than the remaining quota. At position $j$, the DATA
block has
\[
 h=h(A_{\beta,r,g,t})+j+1,\qquad
 k=k_{\beta,r}(A_{\beta,r,g,t})+j+1.
\]
A new position is a \defterm{fresh position}, and its block is a
\defterm{fresh block}. The anchor itself is not fresh. The last position
$j=m_{\beta,r,g,t}-1$ is the \defterm{terminal position}, and its QC is
the \defterm{terminal QC}.

If the terminal DATA height of a term is below $H_{\beta,r}$, its QC
admits term $(\beta,r,g,t+1)$, with that terminal block as
$A_{\beta,r,g,t+1}$. If it equals $H_{\beta,r}$, there is no next term.
Only the last DATA term can have fewer than $K$ positions. The initial
generation is always a DATA generation, since $\TB_r\ge1$.

\paragraph{3. Use DRAIN when recovery starts at the endpoint.}
If $g>0$ and $h(A_{\beta,r,g,0})=H_{\beta,r}$, generation $(\beta,r,g)$ is
a \defterm{drain generation}. It has only term $t=0$ and position $j=0$,
with $m_{\beta,r,g,0}=1$. Its one fresh block is a DRAIN, whose parent is
$A_{\beta,r,g,0}$ and whose justification is the generation's ANCHOR QC.
It carries no requests, keeps DATA height $H_{\beta,r}$, and has no next
term.

DRAIN exists to supply a commit timer. ANCHOR voting starts no timer, so
re-certifying $A_{\beta,r,g,0}$ alone does not commit it. At height
$H_{\beta,r}$, a DATA child would exceed the quota. A DRAIN child is
voted on like any proposal, and its timer commits its ancestors,
including $T_{\beta,r}$.
\begin{itemize}
\item A DATA generation that reaches $H_{\beta,r}$ does not append DRAIN.
It waits for $T_{\beta,r}$'s timer or recovers after a quit.
\item A later drain generation may start from an earlier DRAIN.
Each interrupted recovery may add one DRAIN, so their number is not
bounded.
\item The choice depends on the recovery anchor's DATA height, not on
whether a replica has already committed $T_{\beta,r}$.
\end{itemize}
DRAIN stays in the regional log. The external chain contains only DATA,
and the next session inherits $T_{\beta,r}$.

\paragraph{Example.}
Let $h(A_{\beta,r,0,0})=100$, $\TB_r=10$, and $K=4$, so $H_{\beta,r}=110$.
\begin{center}\small
\begin{tabular}{@{}lll@{}}\toprule
View & Anchor height & Fresh DATA heights\\\midrule
$(\beta,r,0,0)$ & 100 & 101--104\\
$(\beta,r,0,1)$ & 104 & 105--108\\
$(\beta,r,0,2)$ & 108 & 109--110\\\bottomrule
\end{tabular}
\end{center}
Recovery from DATA 107 starts a new generation whose first term produces
108--110. Suppose instead that 110 has a QC but a Quit cancels its timer,
and recovery selects 110 again. DATA 111 would exceed the quota, and
ANCHOR voting on 110 starts no timer. So the recovery leader proposes a
DRAIN on top of 110; its timer commits 110. In both cases the endpoint
stays 110.

\subsection{Certificates and signed messages}
\label{sec:certificates}
\Cref{tab:certificates} lists the certificates of a session. Each needs
matching signatures from distinct members of $V_{i_{\beta,r}}$ with total
weight $>F_{i_{\beta,r}}$.

\begin{table}[htbp]\centering\small
\caption{Certificates of a session.}
\label{tab:certificates}
\begin{tabularx}{\linewidth}{@{}>{\raggedright\arraybackslash}p{.25\linewidth}>{\raggedright\arraybackslash}p{.24\linewidth}Y@{}}
\toprule
Certificate & Signed message & Purpose\\
\midrule
\defterm{Proposal QC} $\QC_{\beta,r,g,t}(B)$ & DATA-VOTE or DRAIN-VOTE & Justifies a
child. A terminal QC below $H_{\beta,r}$ also admits the next term.\\
\defterm{ANCHOR QC} & ANCHOR-VOTE & Re-certifies the opening $A_{\beta,r,g,0}$ in
generation $g$, before any fresh proposal.\\
\defterm{BLAME proof} & BLAME & Freezes a generation (\cref{sec:quit}).\\
\defterm{Regional receipt} $\FC_{\beta,r}(B)$ & FINAL & Proves that a correct replica
committed DATA block $B$.\\
\bottomrule
\end{tabularx}
\end{table}

A correct replica casts at most one fresh vote per $(\beta,r,g,t,j)$ and
one ANCHOR vote per $(\beta,r,g)$. DATA and DRAIN votes follow the same
rule but are different signed types. No QC and no ANCHOR vote starts a
commit timer. A QC means only that enough votes exist; it does not mean
that any timer has expired.

\paragraph{Signed context.}
Every signature binds the chain identifier, protocol domain, session
$(\beta,r)$, voting configuration, assigned region and message type.
View messages also bind $(g,t)$ and the L1 authority reference
$\tagg{AUTH}_{r,g}$; their full view is $(\beta,r,g,t)$. The other
signed fields are:
\begin{itemize}
\item \textbf{PROPOSE and fresh votes:} the generation's opening,
position $j$, block type, physical position, DATA height and block hash;
\item \textbf{NEWVIEW and ANCHOR-VOTE:} the chosen opening. NEWVIEW
also binds its justification and entry proof;
\item \textbf{STATUS:} the sender's entry lock and entry proof;
\item \textbf{BLAME:} the view only, with no opening, block hash or
progress count. All complaints for a view sign the same statement;
\item \textbf{FINAL:} the session and DATA fields below, with no $g$ or $t$.
\end{itemize}
Displayed messages include $\beta$ and $r$ but omit the common chain,
protocol, region and configuration fields where they are clear from
context. A verifier also checks the boundary that assigned the pair
$(\beta,r)$; the pair alone is not evidence of authority.

\paragraph{Rank and the base justification.}
Within session $(\beta,r)$, a QC formed in term $(\beta,r,g,t)$ for block
$B$ has \defterm{rank} $(g,t,k_{\beta,r}(B))$, compared lexicographically.
The ANCHOR QC of generation $g$ has rank $(g,0,k_{\beta,r}(A_{\beta,r,g,0}))$,
just below the generation's first fresh QC. A fixed total order on
certified content breaks ties; signer sets and encodings do not matter.
The \defterm{opening identity} of generation $g$ is
$(g,k_{\beta,r}(A_{\beta,r,g,0}),h(A_{\beta,r,g,0}),\Hash(A_{\beta,r,g,0}))$.
As in Part~I, ``opening'' always means the generation's opening block
$A_{\beta,r,g,0}$. QCs are ranked only within one session $(\beta,r)$.

The anchor $A_{\beta,r,0,0}$ also has a \defterm{base justification}
$Q_{\beta,r}^0$: the anchor together with its direct or covering
commitment evidence in the boundary witness
(\cref{sec:coverage,sec:witness}). Its justified block is exactly
$A_{\beta,r,0,0}$, even if the evidence uses a receipt for a higher
block of the old session; that higher block is evidence only.
$Q_{\beta,r}^0$ ranks below every QC of session $(\beta,r)$. Recovery
uses it when no QC of the session is known. A QC of another session is
never a justification in session $(\beta,r)$, so it cannot replace
$A_{\beta,r,0,0}$.

\paragraph{Regional receipts.}
After a replica commits a DATA block $B$ of session $(\beta,r)$, by $B$'s
own timer or as an ancestor, it sends
\[
 \langle\tagg{FINAL},\beta,r,\mathit{cfg}_r,h(B),\Hash(B)\rangle_v
\]
to its region. Matching FINAL messages of weight $>F_{i_{\beta,r}}$ form
the \defterm{regional receipt} $\FC_{\beta,r}(B)$. The receipt carries
$B$ and the ancestry needed to validate it, or refers to a validated
local cache.

On first obtaining a complete result, each correct holder offers it to
the Publish Layer and multicasts the full validated receipt, the DATA
and their dependencies to the fixed regional membership. Duplicates are
removed by session and DATA hash; the outgoing action is kept until it
is dispatched. This rule also covers a receipt used for Completion and
a holder that is Completed or archival. The holder may also send the
result directly to other regions for an early handoff.

\begin{definition}[\defterm{Valid regional receipt}]
\label{def:receipt}
A regional receipt $\FC_{\beta,r}(B)$ is valid when all of the following hold:
\begin{enumerate}[label=(\roman*)]
\item $B$ is a DATA block of session $(\beta,r)$, extends
$A_{\beta,r,0,0}$, and has $h(A_{\beta,r,0,0})<h(B)\le H_{\beta,r}$;
\item distinct members in session $(\beta,r)$'s fixed voting snapshot have
valid signatures on the same
$\langle\tagg{FINAL},\beta,r,\mathit{cfg}_r,h(B),\Hash(B)\rangle$,
with total weight $>F_{i_{\beta,r}}$ and the signed context specified above; and
\item $B$ and its required ancestry pass the block-validation rules of
\cref{sec:admission}.
\end{enumerate}
Regional engines check the full evidence under \cref{sec:admission};
they do not replay other signers' timers. The global finality rule uses
\cref{sec:receipt-kernel} to establish the same intrinsic validity from
the exact DATA and FINAL signatures. The weight threshold guarantees a
correct signer, and that signer's own validation and commit ground the
conclusion.
\end{definition}
Here $\mathit{cfg}_r$ identifies session $(\beta,r)$'s \defterm{fixed
voting snapshot}. $\FC(B)$ without a session subscript means the receipt
from the session that produced $B$. An inherited anchor keeps that
receipt; it does not acquire a new one. The public checker of
\cref{sec:receipt-kernel} establishes condition~(iii) without requiring
any optional admission sidecar from a publisher. The regional engines
still check the full evidence.
\begin{itemize}
\item FINAL binds no generation or term, so FINAL messages for one block
combine across recoveries. The block hash keeps the block's signed
context.
\item If a DRAIN commits $T_{\beta,r}$, replicas send FINAL for
$T_{\beta,r}$, never for the DRAIN.
\item Replicas of session $(\beta,r)$ send no FINAL for $A_{\beta,r,0,0}$.
It keeps its original commitment evidence, which may come from an earlier
session if intermediate sessions added no DATA.
\item A replica may also commit $B$ and its uncommitted ancestors when it
obtains a valid $\FC_{\beta,r}(B)$ of its own session. Completion does
this for $T_{\beta,r}$ (\cref{alg:steady}). Such a commit sends no FINAL,
because the receipt already exists. Only Completion also cancels timers.
Receipts of other sessions serve only as history or anchor evidence;
they never change this session's locks or log.
\end{itemize}

\paragraph{Receipt coverage.}
\label{sec:coverage}
A valid $\FC_{\beta,r}(D)$ also proves commitment of every DATA ancestor
$B$ of $D$ in the same session: at least one correct signer committed
$D$, and committing $D$ commits its ancestors. The \defterm{covering
evidence} for $B$ is the valid receipt for $D$ plus the fully verified
ancestry from $D$ to $B$. It does not claim that the signatures name $B$
directly. For example, a receipt for 110 with the exact path back to 109
proves that 109 was committed.

Coverage never crosses sessions: a receipt of B's session for 111 does
not replace A's evidence for 110. Coverage also does not show that $B$
belongs to the authorized external branch. A later Force may discard $D$
while keeping $B$; the regional evidence for $B$ still holds. External
observers, retained-anchor witnesses and base justifications may use
direct or covering evidence, with no extra FINAL messages or certificate
round. We write ``\defterm{commitment evidence} for $B$'' for either form.

\subsection{Admission, validation, and engine states}
\label{sec:admission}
\paragraph{\defterm{Admission.}}
Replica $v$ admits views as follows.
\begin{itemize}
\item $(\beta,r,0,0)$ is admitted at $v$'s initial Enter for the boundary
that opens session $(\beta,r)$ (\cref{def:ready}). The boundary witness is
the \defterm{entry proof}. The NEWVIEW must choose exactly
$A_{\beta,r,0,0}$.
\item $(\beta,r,g,0)$ with $g>0$ is admitted once $v$ holds
$\tagg{AUTH}_{r,g}$ and a quit proof for generation $g-1$. That proof is
the \defterm{entry proof}. The NEWVIEW is justified by a QC of an earlier
generation of session $(\beta,r)$, or by $Q_{\beta,r}^0$.
\item $(\beta,r,g,t)$ with $t>0$ is admitted once $v$ also holds the
terminal QCs of terms $0,\ldots,t-1$, each for DATA below $H_{\beta,r}$.
\end{itemize}
Term zero is admitted before its NEWVIEW arrives, so its progress checks
run even if the initial leader is silent. Fresh proposals need the
generation's ANCHOR QC. This holds in every generation, including
generation $0$ of the first session from genesis; Part~I's genesis
shortcut does not apply. A message is \defterm{admissible} if its view is
admitted and its signatures are valid.

A verifier that checks a remote session's full regional history has no
local admission event for it. There, a view counts as admitted if the
evidence contains its \defterm{admission witness}: the boundary witness
for $(\beta,r,0,0)$, the chain of entry proofs for $(\beta,r,g,0)$, and
the earlier terminal QCs for $t>0$.

\paragraph{\defterm{Complete regional validation.}}
Complete regional validation checks signatures, authority, admission,
block type, quota, both positions, the L1 and Espresso reference rules,
the parent rule (\cref{tab:parent}), parent QCs, hash links and
application validity throughout the required ancestry. Every relevant
ancestor comes with its signed proposal and block data; a chain of
certificate hashes alone is not enough. Dependencies come from the
message or from a validated local cache. Correct senders send
\defterm{complete validation bundles}. A reference to a dependency is
allowed only if the receiver can resolve it within the delay bound of
\cref{ass:network}. Bundles list each dependency once.

Regional Vote, recovery, and processing of own-session evidence and
receipts always use this complete procedure. Ready also checks the
target's admission, exact state and boundary inputs in full; only the
receipt of the inherited anchor, which another session produced, may be
checked with the public kernel of \cref{sec:receipt-kernel}.

\paragraph{Evidence takes priority.}
Before validating a message or discarding a duplicate, inspect all signed
material in it for a quit proof (\cref{sec:quit}). A proof needs its own
valid signatures and admission witnesses. A valid proof counts even if
the message that carries it is invalid or incomplete. Compare incoming
material with all stored material, not only with objects inside the
message.

\paragraph{Processing proposals.}
Unless retirement has closed the signing gate, every admitted fresh
proposal of the current generation that $v$ has not voted for goes
through the Vote step. Closing programs still inspect and forward signed
evidence. This includes proposals inside attached histories, proposals
of earlier terms, and proposals whose blocks are already committed as
ancestors. A proposal's identity is $(\beta,r,g,t,j,\Hash(B))$; the signer
set or encoding of its parent QC does not change it. Duplicates are
removed by this identity after evidence inspection. Old-generation
proposals are stored without new votes or timers. Future-generation
messages wait for entry.

\paragraph{Engine states.}
A replica runs one engine for each session pair $(\beta,r)$ assigned to
its region. Engines of different sessions stay separate, even if they
share a turn index. The engine of session $(\beta,r)$ has three core
phases. The program that hosts the engine has its own four states
(\cref{sec:session-lifecycle}): it is Running while the engine is Active
or Frozen, and Closing once the engine is Completed or a Close has
arrived.
\begin{itemize}
\item \textbf{\defterm{Active}:} the current generation is admitted and
unfrozen. The engine proposes, votes and changes terms.
\item \textbf{\defterm{Frozen}:} the current generation has a quit proof.
The engine stops proposals, votes and term changes. It still processes
evidence, old votes, QCs, receipts and recovery.
\item \textbf{\defterm{Completed}:} the replica holds
$\FC_{\beta,r}(T_{\beta,r})$ satisfying \cref{def:receipt}. Production in
session $(\beta,r)$ is over, even if the external handoff is still
pending. The engine only serves data.
\end{itemize}
Committing $T_{\beta,r}$ locally does not complete the engine: other
members may still need this replica's recovery votes before a receipt
forms. Each message goes to the engine named by its signed pair
$(\beta,r)$, so evidence from one session or branch never freezes
another.

These states describe the local engine, not external finality. The
external chain may move to another branch while an old engine still
runs. L1 atomic branch updates do not rewrite an engine's state or log. An
effective Force closes the external branch version; the Session Manager
then delivers Close to the program (\cref{sec:retirement}). The engine
of a closing program keeps its Active or Frozen phase until a genuine
Completion or Exit.

\begin{table}[htbp]\centering\small
\caption{Parent selection for the next fresh position $j$ of term $(\beta,r,g,t)$.}
\label{tab:parent}
\begin{tabularx}{\linewidth}{@{}l Y Y@{}}
\toprule
Condition & $B_{\mathrm{parent}}$ & $Q_{\mathrm{parent}}$\\
\midrule
$j=0,t=0$ (DATA or DRAIN) & Opening $A_{\beta,r,g,0}$ & Its ANCHOR QC in $(\beta,r,g,0)$\\
$j=0,t>0$ & Anchor $A_{\beta,r,g,t}$, the terminal block of $(\beta,r,g,t-1)$ & Its proposal QC in $(\beta,r,g,t-1)$\\
$1\le j<m_{\beta,r,g,t}$ & Position $j-1$ of $(\beta,r,g,t)$ & Its proposal QC in $(\beta,r,g,t)$\\
\bottomrule
\end{tabularx}
\end{table}

\section{Regional protocols}
\label{sec:protocols}
\Cref{alg:steady,alg:recovery} run inside the engine of one session
$(\beta,r)$. They are Protocols 1 and 2 of Part~I~\cite{part1}, with the
changes listed below. \RR{Red text} marks where they differ from Part~I.
The Espresso reference $\xi$ is defined in \cref{sec:blocks}.

Two rules connect these protocols to the cross-region machinery. First,
external replay never changes the expiry of a running regional timer;
only a genuine Quit or Completion cancels one. Second, retirement can
block fresh signing, new timers and further entry, through the gate
below.

\paragraph{Notation.}
The full view is $(\beta,r,g,t)$ throughout. Every anchor, receipt and
entry time below names its branch and session explicitly. In these two
protocols, $\Delta$ and $F$ mean $\Delta_{i_{\beta,r}}$ and
$F_{i_{\beta,r}}$. The authority fields $\mathit{cfg}_r$, $\TB_r$ and
$\tagg{AUTH}_{r,g}$ are indexed by $r$ alone: a later branch that reuses
$r$ reads the same immutable authority. Its anchor, endpoint and engine
state can still differ.

\paragraph{Changes from Part~I.}
\begin{enumerate}[label=\arabic*.]
\item \emph{Stake.} Certificates need weight $>F$, not $f+1$ signatures.
\item \emph{Blocks.} A block also carries its session, physical
position, type, L1 reference $\rho$, and Espresso reference $\xi$
(\cref{sec:blocks}). Validity checks all of them.
\item \emph{Quota.} A term has $m_{\beta,r,g,t}\le K$ positions. No term
follows $H_{\beta,r}$.
\item \emph{DRAIN.} A drain generation proposes one DRAIN. Its vote
counts as progress.
\item \emph{Receipts.} Replicas send FINAL after committing DATA. A
receipt for $T_{\beta,r}$ completes the engine (step~5 of \cref{alg:steady}).
\item \emph{Timers.} Completion also cancels timers. A Force never does.
\item \emph{Recovery entry.} After the $\Delta$ wait, a replica enters
when $\tagg{AUTH}_{r,g+1}$ is finalized. $Q_{\beta,r}^0$ can serve as the
entry lock and as the justification.
\item \emph{Genesis.} The initial session $(\beta,0)$ also certifies its
anchor with an ANCHOR QC before its first fresh proposal.
\item \emph{Sessions.} Each message goes to the engine of its session.
The initial session $(\beta,0)$ opens by bootstrap (\cref{alg:bootstrap}).
\end{enumerate}

\paragraph{Retirement gate and core events.}
The following actions require that $\beta$ is not marked closed at the
node (\cref{sec:session-manager}), that the program has not
received Close, and that the program is not
Terminated: any fresh proposal; any DATA, DRAIN or ANCHOR vote; STATUS
and NEWVIEW; BLAME from a progress check; term entry; and recovery entry.
The same check applies to buffered-message and authority-arrival
callbacks.

The gate never prevents checking, comparing, storing and relaying signed
evidence; processing a genuine Quit; executing a surviving commit timer
and sending its FINAL; or applying a valid receipt and Completion. It
does not overwrite a lock or cancel a timer. \Cref{sec:retirement}
defines when the gate closes and how existing work drains. Entry, Quit,
recovery, Completion and timer-drain changes emit a \texttt{CoreEvent} to
the local Session Manager; that event does not change public $G$.

\paragraph{Message-processing order.}\leavevmode
\RR{Replica $v$ first passes each message to the engine of the session
pair $(\beta,r)$ that the message binds. L1 updates go to \cref{alg:force}. Boundary witnesses and receipts of
other sessions go to \cref{alg:rotation,alg:external}.} The engine handles each received
message in this order:
\begin{enumerate}[label=(\arabic*)]
\item check all signed material for a quit proof;
\item check completeness and validity;
\item process attached DATA \RR{and DRAIN} proposals under Vote, parents
first;
\item process attached terminal QCs under step~4 of \cref{alg:steady},
in increasing term order;
\item process the message itself.
\end{enumerate}
When a message and a timer fall at the same time, $v$ handles the
message first.

\subsection{Steady state}
\label{sec:steady}
\Cref{alg:steady} gives the steady state of a session: proposals, votes,
timers, term rotation, and completion.

\begin{protocol}{alg:steady}{Propose, vote, commit, normal rotation\RR{, and completion}}
Replica $v$ runs these steps \RR{in the engine of session $(\beta,r)$}. Its
current term is $(\beta,r,g,t)$.
\begin{enumerate}[label=\textbf{\arabic*)},itemsep=3pt]
\item \textbf{Propose.} Run this step only if $v=L(\beta,r,g,t)$, $g$ is
unfrozen\RR{, and session $(\beta,r)$ is not complete}.
\RR{In a DATA generation,} for each position
$j=0,\ldots,\RR{m_{\beta,r,g,t}}-1$, in order:
 \begin{itemize}
 \item Wait until $v$ holds $B_{\mathrm{parent}}$ and
       $Q_{\mathrm{parent}}$ for $j$ (\cref{tab:parent}).
 \item Build \RR{$B=(\beta,r,k,h,\tagg{DATA},x,\Hash(B_{\mathrm{parent}}),\rho,\xi)$,
       with $k$ and $h$ from \cref{sec:quota} and the references
$\rho,\xi$ from \cref{sec:blocks}}.
       The batch $x$ follows \cref{sec:requests}.
 \item Broadcast
       $P=\langle\tagg{PROPOSE},\beta,r,g,t,j,B,Q_{\mathrm{parent}}\rangle_v$
       with its validation bundle \RR{to regional replicas. }
 \item Process $P$ under Vote.
 \end{itemize}
 \RR{In a drain generation, propose once, at $(t,j)=(0,0)$. The block is
 a DRAIN with parent $A_{\beta,r,g,0}$, no batch, DATA height $H_{\beta,r}$, and the L1
 and Espresso references of $T_{\beta,r}$. Broadcast and process it as above.}
 Do not wait for the parent to commit. Never sign two different
 blocks for one position. A resend repeats the same $P$.

\item \textbf{Vote.} Run this step on every DATA \RR{or DRAIN} proposal
$P$ of term $(\beta,r,g_P,t_P)$ at position $j_P$. Here $t_P$ may be earlier
than $t$. The step also applies to proposals inside attached histories.
 \begin{enumerate}[label=(\alph*)]
 \item If $P$ or its attached material contains a quit proof for $g$,
       run Quit and stop.
 \item If $g_P<g$, keep $P$ and stop. If $g_P>g$, buffer $P$
       until $v$ enters generation $g_P$.
 \item If $g$ is frozen \RR{or session $(\beta,r)$ is complete}, keep $P$ and
       stop.
 \item If a dependency is missing, keep $P$, request the dependency,
       and stop. Rerun this step when it arrives.
 \item If $P$ is invalid, discard it and stop. \RR{Validity includes the
       block type, the quota, both positions, and the L1 and Espresso references.}
 \item If $v$ has already voted for $P$, stop.
       A duplicate never restarts a timer.
 \item Otherwise, let $\tau_v(P)$ be the current time. Forward $P$ and
       its bundle to all \RR{regional replicas. }
       Broadcast
       $\langle\tagg{DATA\mbox{-}VOTE},\beta,r,g,t_P,j_P,\Hash(B)\rangle_v$
       \RR{(DRAIN-VOTE if $B$ is a DRAIN)}.
       Start the commit timer of $P$. It expires at $\tau_v(P)+2\Delta$.
 \end{enumerate}

\item \textbf{Commit.} When the timer of $P$ expires, commit $B$ and
its uncommitted ancestors. \RR{For each DATA block $D$ of session $(\beta,r)$
that this commits, send
$\langle\tagg{FINAL},\beta,r,\mathit{cfg}_r,h(D),\Hash(D)\rangle_v$ to the regional fixed membership and the Publish Layer collection interface. A committed DRAIN yields FINAL for $T_{\beta,r}$, never for
itself. } Only Quit \RR{and Completion} cancel a timer. Entering a
new term does not cancel it. A QC never starts a timer. \RR{A force
instruction never cancels a timer.}

\item \textbf{Terminal QC and rotation.} Run this step the first time
$v$ obtains a valid terminal QC for a term $(\beta,r,g,t')$.
``Obtains'' means that $v$ assembles the QC from votes or finds it in
any message, including as another proposal's $Q_{\mathrm{parent}}$.
 \begin{enumerate}[label=(\alph*)]
 \item Process each unseen proposal in the QC's history under Vote,
       parents first.
 \item If $g$ is now frozen, stop. \RR{If the QC certifies DATA at
       $H_{\beta,r}$, or a DRAIN, stop: no next term exists.}
 \item Send the QC and its validation bundle to $L(\beta,r,g,t'+1)$.
 \item If $t\le t'$, enter $(\beta,r,g,t'+1)$ and set $e_v(\beta,r,g,t'+1)$ to the
       current time. Cancel the progress checks of all earlier
       terms.
 \end{enumerate}
 A QC never moves $v$ to an earlier term. It never changes a recorded
 entry time.

\item \RR{\textbf{Completion.} Run this step the first time $v$ obtains a
regional receipt $\FC_{\beta,r}(T_{\beta,r})$ satisfying \cref{def:receipt}.}
 \begin{itemize}\color{rrRed}
 \item Multicast the complete validated terminal receipt and its dependencies
       to the fixed regional membership, once, and persist this outgoing action.
 \item Commit $T_{\beta,r}$ and its uncommitted ancestors.
 \item Mark session $(\beta,r)$ Completed. Stop its proposals, votes, progress
       checks, term changes, and recovery. Cancel its unexpired timers.
 \item Offer the complete terminal result to the Publish Layer.
       It may also be sent directly to the next region.
 \item Keep serving data. \Cref{alg:rotation} runs the handoff.
 \end{itemize}
 \RR{A force instruction never runs this step.}
\end{enumerate}
\end{protocol}

\Needspace{9\baselineskip}
The \defterm{acceptance time} $\tau_v(P)$ is when complete validation and the first
vote occur, not an earlier arrival of an incomplete packet. Each replica
records its own entry time for each term. In session $(\beta,r)$, $e_v(\beta,r,0,0)$ is
set at the atomic Enter call, in step~1 of \cref{alg:bootstrap}. For $g>0$, $e_v(\beta,r,g,0)$ is
set in step~3 of \cref{alg:recovery}. For $t>0$, $e_v(\beta,r,g,t)$ is set in
step~4 of \cref{alg:steady}. A term's progress checks therefore start
when the term is admitted locally, not when its leader's first proposal
arrives. A silent new leader cannot delay them.

A proposal or QC at height $H_{\beta,r}$ does not end the session. Without a
quit, the timer of $T_{\beta,r}$ commits it; no later block is needed. Only a
receipt for $T_{\beta,r}$ completes the engine.

\subsection{Quit proofs and recovery}
\label{sec:quit}
A quit proof freezes a generation. \Cref{alg:recovery} then moves the
session to the next generation.

\begin{protocol}{alg:recovery}{Blame, quit, and recovery}
Replica $v$ runs these steps \RR{in the engine of session $(\beta,r)$}. Its
current term is $(\beta,r,g,t)$.
\begin{enumerate}[label=\textbf{\arabic*)},itemsep=3pt]
\item \textbf{Blame.} For each $p=1,\ldots,\RR{m_{\beta,r,g,t}}$ (one
\RR{in a drain generation}), run a
progress check at time
\[
 e_v(\beta,r,g,t)+(2p+4)\Delta.
\]
\RR{Schedule the check for $p=1$ at entry, even if the opening is not yet
known. Schedule the others once $m_{\beta,r,g,t}$ is known. A check whose time has
already passed runs at once.}
At that time, broadcast $\langle\tagg{BLAME},\beta,r,g,t\rangle_v$ if all of
these hold:
 \begin{itemize}
 \item $v$'s current term is still $(\beta,r,g,t)$, $g$ is unfrozen\RR{, and
       session $(\beta,r)$ is not complete};
 \item $v$ has cast DATA \RR{or DRAIN} votes for fewer than $p$ positions
       of term $(\beta,r,g,t)$ (ANCHOR votes do not count);
 \item $v$ has not yet sent BLAME for $(\beta,r,g,t)$.
 \end{itemize}

\item \textbf{Quit.} Run this step the first time $v$ holds a quit proof
for its current generation $g$.
 \begin{itemize}
 \item Freeze $g$. Stop all progress checks, proposals, votes, and
       term changes in $g$. Cancel every unexpired commit timer of $g$.
 \item Set $a_v$ to the current time. Broadcast the proof.
 \item Keep all commitments. Keep collecting votes and QCs of $g$.
       They can raise the entry lock chosen in step~3. \RR{Keep
       processing regional receipts (step~5 of \cref{alg:steady}).}
 \item If the proof is of type (b) or (c), submit it to L1. Do
       not wait for L1 to process it.
 \end{itemize}

\item \textbf{Status and entry.} \RR{Skip this step if session $(\beta,r)$ is
complete.} \RR{Wait until at least $a_v+\Delta$ and until
$\tagg{AUTH}_{r,g+1}$ is finalized. Set $e_v(\beta,r,g+1,0)$ to the time both
conditions hold, and read that authority from L1.} Then:
 \begin{itemize}
 \item Set the entry lock $Q_v$ to the highest-ranked QC
       (DATA\RR{, DRAIN,} or ANCHOR) that $v$ knows from generations
       $0,\ldots,g$\RR{, or to $Q_{\beta,r}^0$ if it knows none}.
       $Q_v$ never changes after this. QCs that arrive later do not
       update it.
 \item Enter $(\beta,r,g+1,0)$ as unfrozen and start its progress checks.
       \RR{Keep $A_{\beta,r,0,0}$, $H_{\beta,r}$, and the voting snapshot of session $(\beta,r)$.}
 \item Send $\langle\tagg{STATUS},\beta,r,g+1,0,Q_v,E\rangle_v$ and its
       validation bundle to $L(\beta,r,g+1,0)$. Here $E$ is the quit proof for
       $g$, which is also the entry proof of $g+1$.
 \item Process any buffered proofs for $g+1$.
 \end{itemize}

\item \textbf{New-view.} Run this step only at $L(\beta,r,g+1,0)$. Wait $2\Delta$
after its own entry. If $g+1$ is still current and unfrozen:
 \begin{itemize}
 \item Let $Q^*$ be the highest-ranked QC from generations $0,\ldots,g$
       that the leader knows, including the QCs in STATUS
       messages\RR{; if it knows none, let $Q^*=Q_{\beta,r}^0$}.
 \item Let $A_{\beta,r,g+1,0}$ be the block that $Q^*$ certifies. \RR{If
       $h(A_{\beta,r,g+1,0})<H_{\beta,r}$, generation $g+1$ is a DATA generation.
       Otherwise it is a drain generation.}
 \item Broadcast
       $\langle\tagg{NEWVIEW},\beta,r,g+1,0,A_{\beta,r,g+1,0},Q^*,E\rangle_{L(\beta,r,g+1,0)}$
       with its validation bundle. Choose only one opening block for
       $g+1$.
 \end{itemize}

\item \textbf{Anchor vote.} On receiving a NEWVIEW for $g+1$:
 \begin{enumerate}[label=(\alph*)]
 \item If it contains a quit proof for $g+1$, run Quit and stop.
 \item Validate its leader, entry proof $E$, opening block $A_{\beta,r,g+1,0}$,
       $Q^*$, and dependencies. If data are missing, request them and
       wait without voting. If the message is invalid, discard it.
 \item Vote only if all of these hold: $g+1$ is current and unfrozen; $v$
       has cast no ANCHOR vote in $g+1$; and $\rank(Q^*)\ge\rank(Q_v)$.
 \item To vote, forward the NEWVIEW to all replicas. Broadcast
       \[\langle\tagg{ANCHOR\mbox{-}VOTE},\beta,r,g+1,0,\Hash(A_{\beta,r,g+1,0})\rangle_v.\]
       Start no timer. Use no fresh position.
 \end{enumerate}
 Matching ANCHOR votes \RR{of total weight $>F$} form the ANCHOR QC.
 When $L(\beta,r,g+1,0)$ obtains it, the leader starts Propose in
 \cref{alg:steady} at $j=0$\RR{; in a drain generation, this is the
 DRAIN}. Every replica, including one that cast no ANCHOR vote,
 processes DATA \RR{and DRAIN} proposals backed by a valid ANCHOR QC
 under Vote.
\end{enumerate}
\end{protocol}

\paragraph{\defterm{Complete quit proofs.}}
For an admitted generation $g$ \RR{of session $(\beta,r)$}, a proof consists of
either:
\begin{enumerate}[label=(\alph*)]
\item distinct signatures \RR{of total weight $>F$} on the same
$\langle\tagg{BLAME},\beta,r,g,t\rangle$ for an admitted term, with that term's
admission witness;
\item two different admissible DATA \RR{or DRAIN} blocks signed by
$L(\beta,r,g,t)$ at the same $(\beta,r,g,t,j)$, with their admission \RR{and quota}
witnesses; or
\item two valid NEWVIEWs signed by $L(\beta,r,g,0)$ selecting different opening
blocks for generation $g$.
\end{enumerate}
\RR{A DATA block beyond the quota, or a DRAIN outside a drain generation,
is invalid. It is rejected and never counts for (b).}
Different openings count even when their blocks are comparable, because
they fix different starting heights $h(A_{\beta,r,g,0})$. The height formula then
maps the same position to different heights.

Different QCs for the same opening block are not equivocation. Evidence
from any admitted term can freeze its generation, including a term not
yet locally reached. Old-generation proofs are archival and may be
submitted to L1. \RR{Evidence from another session never freezes this
engine.} Future-generation proofs are buffered until entry. Neither
resets an unrelated recovery.

The Publish Layer uses these DATA and FINAL receipts. No publication
acknowledgement is required for regional commitment or Completion.
Participants may forward regional bundles directly; no external send
or Espresso confirmation delays a regional vote or timer.

The entry-lock test governs ANCHOR voting only. A replica that has
committed $T_{\beta,r}$ locally still takes part in recovery until it holds
$\FC_{\beta,r}(T_{\beta,r})$, because other members may need its votes. This continued-participation
rule applies until a verified Close closes the fresh-action gate.

\section{Regional and global finality}
\label{sec:external}
Regional commitment and global acceptance are separate events. A node
follows a production branch $R$ and a continuous global DATA prefix
$G$. The branch may extend several sessions beyond $G$. A Force
closes that extension, never $G$ itself.

This section first states what a regional receipt proves
(\cref{sec:regional-finality}) and gives the public check of a receipt
(\cref{sec:receipt-kernel}). It then defines the global chain $G$ and
global finality (\cref{sec:global-chain}), and the rule by which every
verifier computes $G$ from the Espresso log and L1
(\cref{sec:exports}).

\subsection{Regional finality}
\label{sec:regional-finality}
\begin{definition}[\defterm{Regional finality}]
\label{def:regional-finality}
A DATA block $B$ has regional finality on a specified production branch
when its boundary history assigns it to its claimed region and session,
it has valid direct or same-session covering commitment evidence, and
the DATA path to it is valid and continuous. Each session on that path
supplies its own authority and commitment evidence.
\end{definition}
The region is derived from the boundary before its receipt is checked.
A valid receipt therefore cannot authorize another region's DATA or a
different session. Current live/off state does not replace historical
authority. When a Force closes an unconfirmed suffix, its receipts remain
evidence about the old regional engine, but they no longer let that
suffix extend the global prefix.

The progress statements in the remark below concern sessions that run
normally, with no Close during the interval considered. Closing
sessions follow \cref{sec:retirement}; there a replica may keep a local
commitment for which no complete receipt ever forms.

\paragraph{Remark (Commitment by one replica and by all replicas).}
Here ``all replicas'' means all correct voters of the session that
produced $B$, with recovery authorities set in advance.
\begin{enumerate}[label=(\arabic*),itemsep=4pt]
\item \textbf{What a receipt proves.} A valid direct or covering receipt
proves that at least one correct replica has committed $B$. Other correct
replicas may still have running timers, or may have canceled their timers
on Quit. Regional finality does not claim that they have all committed.

\item \textbf{When all replicas commit.} Regional safety preserves $B$
through recovery: the selected anchor is $B$ or a descendant of $B$.
ANCHOR voting re-certifies that block but starts no commit timer.
Suppose a correct replica votes for and forwards a fresh DATA or DRAIN
proposal for a descendant $D$ at time $s$. If no correct replica freezes
that generation through $s+3\Delta_{i_{\beta,r}}$, all correct replicas
commit $B$ by that time. Replicas that have not completed the session
vote by $s+\Delta_{i_{\beta,r}}$; their timers commit $D$ and its
ancestors at most $2\Delta_{i_{\beta,r}}$ later. Completion may cancel a
timer, but it first commits $T_{\beta,r}$ and its ancestors, including
$B$. In particular, a DRAIN need not itself commit at a replica that
completes through the receipt for $T_{\beta,r}$.

\item \textbf{Why one honest leader is not enough.} An honest leader may
finish its term before its timers expire. A Byzantine successor can then
equivocate and freeze the generation. What is needed is the uninterrupted
commit window of (2). Under regional safety and the progress conditions
of \cref{sec:progress}, including permanent equivocation exclusion and
fair leader schedules, the engine eventually obtains such a window while
it keeps running, or completes through the terminal receipt. Earlier
recoveries do not revoke $B$. These conditions give no fixed bound on
when the window begins.

\item \textbf{When a replica can know.} A correct regional replica that
directly commits $D$ by its timer at time $c$ can conclude that all
correct regional replicas have committed $B$ if it stays Active in the
same generation through $c+2\Delta_{i_{\beta,r}}$, after processing the
messages due at that time. Other replicas' timers are due by
$c+\Delta_{i_{\beta,r}}$. A Quit that canceled one of those timers would
reach this replica by $c+2\Delta_{i_{\beta,r}}$ and freeze its
generation. A replica whose timer was canceled by Completion has already
committed $B$. Normal term rotation does not affect this test.

\item \textbf{What an external observer can know.} A receipt alone does
not show that the uninterrupted window occurred. The delivery bounds
cover complete submitted results, not every regional BLAME or quit event,
so the absence of a published quit proof does not prove that no replica
quit. Espresso evidence shows that the result is publicly available. A
bound on when every replica runs its local commit step would also need
the relevant online and receipt-processing conditions. This protocol
protects the external prefix against Force through Espresso acceptance;
it uses no external all-replicas-committed test.
\end{enumerate}

\subsection{Public receipt verification}
\label{sec:receipt-kernel}
This subsection answers one question: given a DATA block $B$, can anyone
check from public data that $B$ is regionally final
(\cref{def:regional-finality})? The answer is the check
$\tagg{CheckFINAL}$. It needs three things: the context of the session
that produced $B$, which the verifier derives itself; FINAL signatures
on a root block $D$; and the exact bytes that link $D$ back to $B$. It
does not need the proposal or recovery history that a publisher may
attach, which can name recovery authorities the verifier does not have.

\paragraph{Who uses it.}
The global finality rule (\cref{sec:exports}) and handoff verifiers use
$\tagg{CheckFINAL}$ for receipts that another session produced, such as
a historical anchor or a REGULAR terminal receipt. A handoff verifier
need not fetch the producer's optional recovery-admission path, but it
still checks the exact inherited DATA and path, the quota and the
canonical boundary and selection. Everything else keeps full
validation: own-session Vote, recovery and receipt-driven Completion;
Ready, which validates the target boundary, authority and exact
executable state; and admission in the target session.

\paragraph{The session context $C$.}
Every check is made against the \defterm{session context} $C$ of the
producing session $(\beta,r)$. $C$ consists of two authenticated
objects: the session's incoming boundary (the handoff record that opened
it, \cref{sec:witness}) and its initial authority $\tagg{AUTH}_{r,0}$.
They fix the following fields.
\begin{center}\small
\begin{tabularx}{\linewidth}{@{}>{\raggedright\arraybackslash}p{.36\linewidth}>{\raggedright\arraybackslash}p{.2\linewidth}Y@{}}\toprule
Field & Source & Example: A's $(7,20)$\\\midrule
chain identifier and protocol domain & boundary & this chain\\
producing pair $(\beta,r)$ & boundary & $(7,20)$\\
assigned region $i$ & boundary & A\\
anchor $A_{\beta,r,0,0}$: height, hash, original producer & boundary & DATA 100\\
fixed voting snapshot $\mathit{cfg}_r$: members, weights, $F_i$ & $\tagg{AUTH}_{r,0}$ & $\mathit{cfg}_{20}$\\
DATA quota $\TB_r$, hence $H_{\beta,r}$ & $\tagg{AUTH}_{r,0}$ & 10, so $H=110$\\\bottomrule
\end{tabularx}
\end{center}
The verifier derives $C$ from public history, never from the sender of a
record. The global finality rule (\cref{alg:global-rule}) installs a
regular boundary when it accepts the predecessor's terminal DATA, a
forced boundary at a Force, and fixes the region by applying
$\tagg{NextLive}$ to the branch's immutable map; the
initial checkpoint supplies the genesis boundary; a handoff verifier
derives the same boundary with $\tagg{VerifyHandoffHistory}$
(\cref{sec:verify-handoff}). $C$ contains no recovery authority
$\tagg{AUTH}_{r,g}$ with $g>0$ and no generation or term: FINAL does not
bind them, and the voting snapshot is the same in every generation of
the session.

\paragraph{The check.}
Write
\[
 \tagg{CheckFINAL}(C,B,D,\Sigma,\pi).
\]
Its inputs are the following.
\begin{center}\small
\begin{tabularx}{\linewidth}{@{}>{\raggedright\arraybackslash}p{.05\linewidth}Y>{\raggedright\arraybackslash}p{.15\linewidth}>{\raggedright\arraybackslash}p{.22\linewidth}@{}}\toprule
Input & What it is & Supplied by & Example: A's $(7,20)$\\\midrule
$C$ & the session context above & derived by the verifier & anchor DATA 100, quota 10\\
$B$ & the candidate DATA to check & publisher & DATA 103\\
$D$ & the root DATA that carries the FINAL certificate: $B$ itself, or a same-session descendant of $B$ & publisher & DATA 105\\
$\Sigma$ & the FINAL signatures on $D$ & publisher & FINAL on DATA 105\\
$\pi$ & the exact hash-linked path from $D$ back to $B$, with full bodies & publisher & $105\to104\to103$\\\bottomrule
\end{tabularx}
\end{center}
The triple $(D,\Sigma,\pi)$, read against $C$, forms the \defterm{kernel} for $B$.
The check has two parts.
\begin{enumerate}[label=(\arabic*),itemsep=2pt]
\item \textbf{Signatures.} Write $\tagg{RawFINAL}_C(D,\Sigma)$ when the
bytes $D$ have the named hash and height, and $\Sigma$ contains valid
FINAL signatures from distinct members of $\mathit{cfg}_r$, with total
weight $>F_i$. The signed context must match $C$ in full: chain,
protocol domain, message type, producing pair, assigned region and
configuration, even where the displayed FINAL notation omits those
fields. ``Raw'' means a cryptographic and type check only. It does not
validate the block, include \cref{def:receipt}(iii), read a generation
or term, or check any attached admission witness.
\item \textbf{Path.} Each body on $\pi$ is an actual DATA or DRAIN body
with the correct hash, parent link and position; $\pi$ ends at $B$; $D$
and $B$ belong to the same session $(\beta,r)$; and
$h(B)\le H_{\beta,r}$. For a direct receipt, $D=B$ and $\pi$ is just $B$.
A hash without its payload is not enough, and a later session's receipt
never covers an earlier session's DATA.
\end{enumerate}
The check returns \emph{true} if both parts pass, \emph{false} if one
fails, and \emph{incomplete} if a required object is absent. The caller
handles \emph{incomplete} under the missing-object rule; it is not an
instruction to fetch optional sidecars. Required bytes must be supplied
or resolved from established public objects.

In the example, part (1) checks that the FINAL signatures on DATA 105
have weight $>F_A$ under $\mathit{cfg}_{20}$. Part (2) checks that
$105\to104\to103$ is correctly linked inside $(7,20)$ and that
$103\le110$. The check returns \emph{true}: DATA 103 is regionally final.

\paragraph{Why it suffices.}
\begin{proposition}[Soundness of the public FINAL kernel]
\label{prop:final-kernel-soundness}
Let $C$ be the authentic incoming context of session $(\beta,r)$, and
suppose $\tagg{RawFINAL}_C(D,\Sigma)$ holds. Then some correct replica of
that session committed exactly this DATA $D$ through its own commit
timer, and before doing so it fully validated $D$ and the inputs $D$
depends on. Hence $D$ is a validly admitted regional input under its
anchor, including the $\rho$ and $\xi$ conditions, whatever witness a
publisher happens to attach. The same holds for every same-session DATA
ancestor of $D$ on its hash-linked path.
\end{proposition}
This relies on the standing assumptions: the fault bound of the fixed
snapshot, the key-compromise bound, the cryptographic assumptions, sound
validation interfaces, and the regional rules for full validation and
for sending FINAL. The intuition is short. Weight above $F_i$ contains a
correct signer; a correct replica signs FINAL only after its own timer
commits the block; and it starts that timer only after full validation.
Commits caused by a receipt send no FINAL, so they cannot be the source
of the first correct signature. The precise statement and proof are in
Appendix~\ref{app:final-kernel-proof}.

So if $\tagg{CheckFINAL}$ returns \emph{true}, part (2) places $B$ on
$D$'s validated path and the proposition makes $B$ regionally final.
The proposition shows that a valid witness \emph{exists}; it does not
certify the one a publisher attaches. Because FINAL binds no generation
or term, a publisher cannot force a fetch of $\tagg{AUTH}_{r,g}$ by
naming a different $g$, and optional archival validation never gates
acceptance of a complete kernel.

\paragraph{What it does not check.}
A \emph{true} result says that $B$ is regionally final, and nothing more.
\begin{itemize}[itemsep=2pt]
\item \textbf{Global acceptance.} Whether $B$ exactly extends the
current $G$, the current branch permission and export duty, the strict
precedence of Espresso references, and the cutoff after a
Force are checked by the caller (\cref{sec:exports}). The caller must
still obtain the complete L1 control history; FINAL cannot vouch
for it.
\item \textbf{Uniqueness.} The threshold $>F_i$ alone is not a
quorum-intersection argument. Regional safety, proved separately under
its premises, makes competing valid commitments compatible.
\item \textbf{Execution.} The kernel certifies only the hash-bound
input history. Unbound native context, state objects, execution output
and later DA publications keep their own execution and DA
checks; a receipt is not a state object, so Ready still requires the
exact native state. A native claim bound by the DATA hash and checked by
the regional validity interface is covered, and only that claim.
\item \textbf{Delivery.} The kernel is a subset of a complete result.
It does not relax the delivery and retention rules, adds no deadline,
and proves nothing about request inclusion.
\end{itemize}

\subsection{The global chain}
\label{sec:global-finality}
\label{sec:global-chain}
The \defterm{global chain} $G$ is the sequence of DATA blocks that every
verifier accepts. It starts at the genesis anchor $A_0$ and passes
through a sequence of sessions.

\paragraph{Session descriptors and block-bound records.}
A \defterm{session descriptor} is
\[
 \nu=(\beta,\ r,\ i,\ A,\ H,\ q^{\rm open},\ \mathit{cause}),\qquad
 \mathit{cause}\in\{\tagg{GENESIS},\ \tagg{REGULAR},\
 \tagg{FORCE}(c_f,q_f)\}.
\]
Its session identifier is $\operatorname{sid}(\nu)=(\beta,r)$; $i$ is
the selected region; $A$ is the inherited anchor;
$H=h(A)+\TB_r$ is the terminal DATA height; and $q^{\rm open}$ is the
lower bound on the Espresso header that the session's first DATA may
reference. It is not a selection time. The cause says how the
session was opened; for a Force, $c_f$ identifies the executing atomic
update and $q_f$ is its activation block.  We write $\beta(\nu)$,
$r(\nu)$, $i(\nu)$,
$A(\nu)$, $H(\nu)$, $q^{\rm open}(\nu)$ and
$\mathit{cause}(\nu)$ for these fields.  The descriptor fixes the session context $C_\nu$
of \cref{sec:receipt-kernel}: $\nu$ supplies the pair, region and anchor,
and $\mathit{cfg}_r$ and $\TB_r$ come from $\tagg{AUTH}_{r,0}$.  The L1
bound $\lambda^{\rm open}_{\beta,r}$ is not part of $\nu$; only regional
validation uses it.

Whenever the global finality rule appends a non-genesis DATA block
$B$ to $G$, it binds the immutable \defterm{session record}
\[
 \sigma_G(B):=\nu
\]
using the complete descriptor under which $B$ was accepted. This is
global-chain metadata, not a separately submitted or signed L1 or
Espresso record.  Every DATA produced by the same fixed session binds
the same descriptor value.  A selected session that contributes no DATA
to $G$ leaves no session record in $G$.

\begin{definition}[\defterm{Global chain}]
\label{def:global-chain}
A global chain is a sequence
\[
 G=A_0\cdot B_1\cdots B_n
\]
together with the map $B\mapsto\sigma_G(B)$ on its non-genesis blocks,
such that:
\begin{enumerate}[label=(\roman*),itemsep=2pt]
\item \textbf{Continuity.} Every $B_j$ has $B_{j-1}$ as its parent and
the next DATA height, where $B_0=A_0$.
\item \textbf{Ownership.} If $\sigma_G(B_j)=\nu_j$, then $B_j$ is
produced by $\operatorname{sid}(\nu_j)$, has height at most $H(\nu_j)$,
and is regionally final under $C_{\nu_j}$
(\cref{def:regional-finality}).
\item \textbf{Session runs.} Partition $B_1,\ldots,B_n$ into maximal
nonempty consecutive runs on which $\sigma_G$ is constant.  The first
run, if any, has turn index $0$ and anchor $A_0$.  If a run with
descriptor $\nu'$ follows one with descriptor $\nu$, then
$A(\nu')$ is the last block of the earlier run and
$r(\nu')=r(\nu)+1$.  If
$\mathit{cause}(\nu')=\tagg{REGULAR}$, the earlier run ends at
$H(\nu)$ and $\beta(\nu')=\beta(\nu)$.  If
$\mathit{cause}(\nu')=\tagg{FORCE}(c_f,q_f)$, then
$\beta(\nu')>\beta(\nu)$.  A run that ends below its terminal height can
only be followed by a Force-opened run.
\end{enumerate}
\end{definition}
No empty run appears in $G$.  A selected session with no accepted DATA
is represented by $\tagg{NEXT}(\nu)$ only while it owns the next block;
if a Force replaces it, it leaves no block or session record in $G$.
Repeated Forces before the next accepted DATA may therefore make the
branch version jump by more than one between two adjacent runs.

The rule of \cref{sec:exports} fixes which global chain each verifier
holds at each Espresso height. In the running example, each DATA in the
listed range binds the displayed record:
\begin{center}\small
\begin{tabular}{@{}ll@{}}\toprule
DATA in $G$ & Bound session record\\\midrule
$101,\ldots,108$ & $(7,20,\text{A},100,110,q_{\rm A},\tagg{REGULAR})$\\
$109',110',\ldots$ & $(8,21,\text{B},108,108+\TB_{21},q_f,\tagg{FORCE}(c_f,q_f))$\\\bottomrule
\end{tabular}
\end{center}
A's run stops at 108 because the Force closed version 7 before 109
was accepted. B's regular session $(7,21)$ and its DATA 111--115 do not
appear, because $G$ never passed through them. The production branch
$R$ may run ahead of $G$; the global chain records only what was
accepted.

\begin{definition}[\defterm{Global finality}]
\label{def:global-finality}
A DATA block $B$ is \defterm{globally final at Espresso height $q$} when
$B\in G_q$, the global chain output by the rule of \cref{sec:exports} at
$q$. Observer $u$ treats $B$ as globally final once it has authenticated
the finalized Espresso history through $J_q$ and the complete L1 control
history that applies to it, and has applied the rule through a step that
appends $B$.
\end{definition}

The rule only appends, so $G_{q-1}$ is a prefix of $G_q$ for every $q$
(\cref{prop:global-prefix}). No later Force removes a globally final
block $B$ or changes the record $\sigma_G(B)$ bound to it, and no extra waiting period
follows. Espresso inclusion alone, or a receipt alone, does not make a
block globally final: an old-version block may be finalized in Espresso
as raw bytes and still never enter $G$.

This is finality of the input prefix. Espresso neither executes the
native transactions nor proves their resulting state; the regional
validity contract and the execution-verification interface are
responsible for that. Finality of a state claim or withdrawal on L1 is a
later, separate stage. Publication lag never permits skipping globally
final input.

\subsection{The global finality rule}
\label{sec:exports}
The global finality rule computes the global chain and the selection
information for its next DATA. It is a deterministic function of its
two input streams, so every verifier that reads the same history
computes the same output.

Between heights, the verifier retains the minimal \defterm{replay state}
\[
 \mathcal X_q=(\mathrm{Out}_q,S_{\beta_q},\kappa_q),
\]
where $\beta_q$ is read from $\mathsf{Tail}_q$, $S_{\beta_q}$ is that
branch's immutable complete map, and $\kappa_q$ is the processed control
cursor of \cref{sec:l1}. These are maintained state, not additional
per-height input streams. In particular, no historical sequence of
status maps is kept for successor selection. The public output of the
protocol remains exactly $\mathrm{Out}_q=(G_q,\mathsf{Tail}_q)$.

\paragraph{Input.}
\label{sec:logical-carrier}
For each newly finalized Espresso block $J_q$, Protocol~3 consumes
exactly two logical input streams:
\begin{enumerate}[label=(I\arabic*),itemsep=2pt]
\item \textbf{Espresso-finalized regional results.}  The authenticated
namespace of $J_q$ adds a set $M_q$ of regional objects.  Complete
regional results consist of DATA bodies, FINAL signatures and any
covering path needed to establish regional finality.  The rule does not
depend on how EXPORT records group these objects.  It uses:
 \begin{itemize}[itemsep=1pt]
 \item block bodies
 $B=(\beta,r,k,h,\mathit{type},x,\Hash(B_{\rm parent}),\rho,\xi)$ of type
 DATA or DRAIN, identified by their hash;
 \item FINAL signatures: a signature of a member on the hash of a block
 $D$, in the signed context of $D$'s session.
 \end{itemize}
Write $\mathcal O_q=M_1\cup\cdots\cup M_q$ for everything published
through $J_q$. A DATA block is an append candidate only when the objects
in $\mathcal O_q$ establish $\mathrm{Fin}_q(\nu,B)$ under the current
next-block descriptor.
\item \textbf{The activated update batch.}  From the complete finalized
L1 control log, let $\Lambda_q$ be the ordered valid updates whose
activation thresholds are first crossed at $J_q$
(\cref{def:control-effect}).  The first entry targeting the branch
version at the start of the prelude, if any, atomically creates the next
version and its fixed map; every other entry is a whole-update no-op.
\end{enumerate}
The genesis checkpoint supplies the initial state rather than a third
per-block input: the chain identifier, region order $\mathcal L_R$,
anchor $A_0$, first replay position $q_{\rm init}$, initial map $S_0$
and control prefix. A later trusted checkpoint supplies an earlier
$\mathcal X$ instead. Espresso headers, namespace-completeness
proofs, $\tagg{AUTH}_{r,0}$ and L1 authentication proofs are validation
dependencies for the two inputs above, not additional logical input
streams.

Thus the per-height transition consumes only $\mathcal O_q$ and
$\Lambda_q$, together with its already persisted $\mathcal X_{q-1}$.
It never rereads the full update log to choose a regular successor.

Proposals, QCs, recovery authorities
$\tagg{AUTH}_{r,g}$ with $g>0$ and native sidecars may also appear in
$M_q$; the rule ignores them. The outer fields of an EXPORT, such as its
expected base, interval and version, are routing hints and are ignored
too. A malformed or irrelevant object grants nothing and blocks nothing.

The rule uses one test derived from the input. For a complete session
descriptor $\nu$ and a candidate DATA block $B$, write
$\mathrm{Fin}_q(\nu,B)$ when
$\tagg{CheckFINAL}(C_\nu,B,D,\Sigma,\pi)$ returns \emph{true} for
some block $D$ and path $\pi$ whose bodies all lie in $\mathcal O_q$, with
$\Sigma$ the set of FINAL signatures on $D$ in $\mathcal O_q$. The
signatures may come from different EXPORT records and Espresso blocks;
only their total weight counts.

\paragraph{Output.}
For every Espresso height $q\ge q_{\rm init}$ the rule outputs
\[
 \boxed{\mathrm{Out}_q=(G_q,\mathsf{Tail}_q)},\qquad
 \mathsf{Tail}_q\in\{\tagg{NEXT}(\nu),\ \tagg{PENDING}(\omega)\}.
\]
$G_q$ is the global chain with its block-bound records
(\cref{def:global-chain}).  The tail $\mathsf{Tail}_q$ says who is responsible for
the next DATA after $\operatorname{last}(G_q)$:
\begin{itemize}[itemsep=2pt]
\item $\mathsf{Tail}_q=\tagg{NEXT}(\nu)$ means that the fully specified session
$\operatorname{sid}(\nu)$ is the \defterm{export session}: it owns the
next DATA height and its descriptor supplies the region, anchor,
terminal height, Espresso reference floor and cause. Its accepted count is
\[
 a(G_q,\nu)=
 \bigl|\{B\in G_q:\operatorname{sid}(\sigma_{G_q}(B))
                 =\operatorname{sid}(\nu)\}\bigr|.
 \]
\item $\mathsf{Tail}_q=\tagg{PENDING}(\omega)$, where
\[
 \omega=(\mathit{cause},\beta,r',i_{\rm src},A,q^{\rm open}),
\]
means that the next session identifier $(\beta,r')$, its anchor
$A=\operatorname{last}(G_q)$ and its opening reference floor are fixed,
but the immutable map $S_\beta$ has no eligible target region after
$i_{\rm src}$.
\end{itemize}
The current branch version is $\beta(\nu)$ in the \tagg{NEXT} case and
the displayed $\beta$ in the \tagg{PENDING} case.  Exactly one tail
variant is present at every completed replay subposition.

\begin{protocol}[H]{alg:global-rule}{Global finality rule: update replay state and output $\mathrm{Out}_q$}
Initialize $\kappa$ to the control cursor in the genesis checkpoint. Use
$S_0$ and $\tagg{AUTH}_{0,0}$ to compute
$i_0=\tagg{NextLive}(\bot,S_0;0)$.  If $i_0\ne\bot$, set
\[
 \nu_0=(0,0,i_0,A_0,h(A_0)+\TB_0,q_{\rm init},\tagg{GENESIS}),
 \qquad \mathrm{Out}=(A_0,\tagg{NEXT}(\nu_0));
\]
otherwise set
\[
 \mathrm{Out}=(A_0,
 \tagg{PENDING}(\tagg{GENESIS},0,0,\bot,A_0,q_{\rm init})).
\]
For each finalized block $J_q$ in order, start from $\mathcal X_{q-1}$,
ingest its two inputs---the activated batch $\Lambda_q$ and the cumulative
Espresso-finalized object set $\mathcal O_q$ (after adding $M_q$)---and
run the two steps below. The public result is $\mathrm{Out}_q$; the
updated map and control cursor complete $\mathcal X_q$. Write
$\operatorname{last}(G)$ for the current tip.
\begin{enumerate}[label=\textbf{\arabic*)},itemsep=3pt]
\item \textbf{Apply an atomic branch update.}
Let $\beta_0$ be the branch version at the start of the prelude.  If
$\Lambda_q$ contains no update targeting $\beta_0$, every due entry is
stale, $\mathrm{Out}$ and $S_{\beta_0}$ are unchanged, and $\kappa$
advances past the batch. Otherwise let
\[
 U^*=(\mathit{chain},c^*,\theta_a^*,\beta_0,S^*)
\]
be the first matching update.  Close $\beta_0$, set
$\beta':=\beta_0+1$ and fix $S_{\beta'}:=S^*$.  Determine the target
index and source region as follows:
  \begin{itemize}[itemsep=1pt]
  \item if $\mathsf{Tail}=\tagg{NEXT}(\nu)$, write
  $\operatorname{sid}(\nu)=(\beta_0,r)$, let $i_{\rm src}=i(\nu)$ and
  $a=a(G,\nu)$, and set $r'=r$ if $a=0$ and $r'=r+1$ otherwise;
  \item if $\mathsf{Tail}=\tagg{PENDING}
  (\mathit{cause},\beta_0,r',i_{\rm src},A,q^{\rm open})$, keep its $r'$ and
  $i_{\rm src}$.
  \end{itemize}
Let $A=\operatorname{last}(G)$ and
$i'=\tagg{NextLive}(i_{\rm src},S_{\beta'};r')$.  If $i'=\bot$, set
\[
 \mathsf{Tail}:=\tagg{PENDING}
 (\tagg{FORCE}(c^*,q),\beta',r',i_{\rm src},A,q).
\]
Otherwise set
\[
 \nu'=(\beta',r',i',A,h(A)+\TB_{r'},q,
       \tagg{FORCE}(c^*,q)),
 \qquad \mathsf{Tail}:=\tagg{NEXT}(\nu').
\]
All other entries of $\Lambda_q$ are stale whole-update no-ops.  This
step advances $\kappa$ past every entry exactly once and never changes
$G$ or an existing $\sigma_G(B)$.

\item \textbf{Append matching DATA.}
While $\mathsf{Tail}=\tagg{NEXT}(\nu)$ and $\mathcal O_q$ contains a
DATA block $B$ satisfying, with
$\operatorname{sid}(\nu)=(\beta,r)$,
  \begin{enumerate}[label=(\alph*),itemsep=1pt]
  \item $B$ is signed in this chain's domain by session $(\beta,r)$;
  \item $\mathrm{parent}(B)=\Hash(\operatorname{last}(G))$ and
  $h(B)=h(\operatorname{last}(G))+1\le H(\nu)$;
  \item $\mathrm{Fin}_q(\nu,B)$ holds; and
  \item $\xi(B)$ names a finalized block $J_{q'}$ of the log, with its
  authenticated hash, and $q^{\rm open}(\nu)\le q'<q$,
  \end{enumerate}
append $B$ to $G$ and bind $\sigma_G(B):=\nu$.  If several blocks
qualify, take the one with the least hash.  If $h(B)<H(\nu)$, leave
$\mathsf{Tail}=\tagg{NEXT}(\nu)$ and continue.  If $h(B)=H(\nu)$, write
$\operatorname{sid}(\nu)=(\beta,r)$ and set
\[
 q^{\rm open}_{\rm next}=
 \max\{\operatorname{height}(\xi(B)),q^{\rm open}(\nu)+1\}.
\]
Let $i'=\tagg{NextLive}(i(\nu),S_\beta;r+1)$. If $i'=\bot$, set
\[
 \mathsf{Tail}:=\tagg{PENDING}
 (\tagg{REGULAR},\beta,r+1,i(\nu),B,q^{\rm open}_{\rm next}).
\]
Otherwise set
\[
 \nu'=(\beta,r+1,i',B,h(B)+\TB_{r+1},
       q^{\rm open}_{\rm next},\tagg{REGULAR}),
 \qquad \mathsf{Tail}:=\tagg{NEXT}(\nu').
\]
Then continue looking for another DATA block.
\end{enumerate}
\end{protocol}

\Cref{alg:global-rule} is the public core of \cref{alg:force} and
\cref{alg:external}. Step~1 consumes the control log and changes only the
branch version, its fixed map and the next-block selection. Step~2
appends exactly the regionally finalized DATA that extends the current
tip and matches that selection. \Cref{sec:force,sec:observers} explain
the waiting rules around these operations.

\paragraph{Why each check.}
\begin{itemize}[itemsep=2pt]
\item \textbf{Export session only} (a). Only the session that owes the
next height may extend $G$. A closed version has no export session, so
after a Force its DATA never enters $G$. A regular successor that runs
ahead becomes the export session only after its predecessor's terminal
DATA is accepted.
\item \textbf{Continuity} (b). $G$ has no gap and no fork. DRAIN never
enters $G$, because it is not DATA and does not advance the height.
\item \textbf{Regional finality} (c). By
\cref{prop:final-kernel-soundness}, a correct replica of the session
committed $B$ after fully validating it. The rule needs no proposal or
recovery history.
\item \textbf{Selection reference} (d). A reference cannot justify itself
through the block that carries it. Because $q^{\rm open}\ge q_f$ for a
session opened by a Force at $J_{q_f}$, no new-version DATA is accepted
in $J_{q_f}$ itself.  A regular successor may already have published
output under its deterministic opening floor; that output can follow the
terminal DATA in the same or a later EXPORT once the terminal enters
$G$.
\end{itemize}

\paragraph{Determinism and accumulation.}
Each step reads only the current replay state and the two inputs, so both
$\mathcal X_q$ and its public component $\mathrm{Out}_q$ are functions of
the input through $J_q$. Under regional
safety at most one block qualifies at each iteration of Step~2, since two regionally final
DATA of one session at one height are identical; the hash tie-break only
makes the rule total. The set $\mathcal O_q$ only grows. An object stays
usable after the block that carried it: if DATA 109 is published before
108 is accepted, Step~2 accepts 109 right after 108, in the same or a
later block, without republication. A wrong hint or an invalid neighbor
in the same EXPORT record removes nothing.

For example, immediately before the Force in the running example,
\[
 G=100\cdot101\cdots108,\qquad
 \mathsf{Tail}=\tagg{NEXT}(\nu_A),\qquad
 \operatorname{sid}(\nu_A)=(7,20),\quad i(\nu_A)=\text{A},\quad
 a(G,\nu_A)=8.
\]
Suppose the executing update has sequence number $c_f$ and activation
block $J_{q_f}$. Step~1 leaves $G$ and its records unchanged and
constructs
\[
 \omega_F=(\tagg{FORCE}(c_f,q_f),8,21,\text{A},108,q_f).
\]
Because B is eligible under the newly installed map $S_8$, Step~1 sets
the tail directly to $\tagg{NEXT}(\nu_B)$, where
\[
 \nu_B=(8,21,\text{B},108,108+\TB_{21},q_f,
        \tagg{FORCE}(c_f,q_f)).
\]
B's new DATA $109'$ carries a reference at least $q_f$, so Step~2 can
accept it from $J_{q_f+1}$ on and bind
$\sigma_G(109')=\nu_B$. A's old 109 is still in $\mathcal O_q$ but
fails (a). If no region were eligible under $S_8$, Step~1 would instead
keep $\tagg{PENDING}(\omega_F)$ until another successful atomic update
created a later version with a new map.

\paragraph{Waiting and partial progress.}
\label{sec:subpositions}
\label{sec:terminal-commit}
The rule is stated per Espresso block, but its steps are atomic, and an
observer may stop between them. Step~1 needs the applicable atomic update
and the target authority $\tagg{AUTH}_{r',0}$; the terminal case of Step~2
needs the already fixed $S_\beta$ and $\tagg{AUTH}_{r+1,0}$. If required
evidence is not yet retrievable, the observer waits in that operation;
missing evidence never counts as \emph{no eligible region}. With complete
evidence and no eligible region, the boundary stays pending and can be
changed only by another successful atomic update. Step~2 without a
matching DATA is not a wait either. When Step~2 appends terminal DATA, it
persists that append before selecting the successor; a wait for successor
authority never withholds the DATA.

The chain held after any step within $J_q$ is a prefix of $G_q$ and may
be exposed at once. An implementation persists the block and the last
completed step, its \defterm{replay subposition}, so that a restart
neither repeats nor skips a step. A duplicate delivery of a source block
resumes at that subposition.

\paragraph{What the transport must supply.}
\label{sec:transport-scope}
The rule can accept only what $\mathcal O_q$ contains. The Publish Layer
must therefore put the objects of a complete regional result into
Espresso: the DATA bodies, the FINAL signatures and any covering path.
The delivery bound $T_{\rm pub}$ of \cref{ass:services} applies to
complete results; a hash, a READY message or a receipt without its
payload is not one. The encoding of EXPORT records is the Publish Layer's
choice, provided every reader recovers the same $M_q$ from the namespace
content. An encoding that hides required objects, or a private cache,
does not supply them.

\subsection{What block references do not tell}
\label{sec:late-receipt}
A block's L1 reference and Espresso reference are fixed when it
is proposed. A later recovery or DRAIN may commit the block long
afterwards. Neither reference shows global finality or authorizes a
Force. Regional timers use the real-time clock of \cref{ass:network}.
Controls take effect at the Espresso threshold of an explicit L1 record.
For example, a block that references an Espresso block with time 100 may
be committed much later after a recovery; it was not committed at time
100.

\section{Cross-region rotation}
\label{sec:xprotocols}
This section describes what each regional node does across sessions:
how it learns about handoffs and Forces, how it prepares and enters a new
session, and how it shuts down an old one. It also fixes the interface
through which the multisignature committee issues controls. The public
rules that execute those controls are in \cref{sec:force,sec:observers};
this section consumes their results.

\paragraph{What B does at the same time.}
Return to the running example. A's last DATA is 110. B holds A's
terminal receipt and produces in $(7,21)$ from 111, while $G$ is still
108. L1 now announces a future Force that targets branch version 7. B
keeps one program and two execution states:
\begin{center}\small
\begin{tabular}{@{}lll@{}}\toprule
Object & What it is now & Continues after\\\midrule
program $(7,21)$ & Running the regional BFT & 110 (fixed)\\
baseline state & $\mathit{st}(\operatorname{last}(G))$ & 108, then follows $G$\\\bottomrule
\end{tabular}
\end{center}
No program exists yet for the announced Force. If $G$ grows to 109
before the Force takes effect, the baseline state follows it to 109.
When the Force takes effect, its anchor is $\operatorname{last}(G)$; if
B is selected, the Session Manager creates program $(8,21)$ with that
anchor, starts it from the baseline state, and stops $(7,21)$. The
baseline state never grants the right to produce.
\Cref{sec:runtime-races} works through this case in detail.

\subsection{Cross-region architecture}
\label{sec:runtime-modules}
Every correct regional replica reads the same public inputs and runs a
Session Manager and its session programs.  Regional nodes, observers
and certifiers directly apply the public replay rules of
\cref{alg:global-rule,alg:force,alg:external}; replay is not a separate
module or service.  \Cref{fig:runtime-modules} shows the resulting data
flow. The Multisig Controller and the Publish Layer run outside the
node.

\paragraph{Components.}
\begin{itemize}
\item The \defterm{Multisig Controller} watches verified progress and
operational reports and appends control updates to the L1 control log
(\cref{sec:runtime-controls}).
\item The \defterm{Publish Layer} carries complete regional results to
Espresso (\cref{sec:delivery}).
\item The \defterm{Session Manager} tracks regional, global and L1
progress derived by the node's replay, keeps the executable state at
$\operatorname{last}(G)$, and starts and stops the node's session
programs (\cref{alg:rotation}). It also keeps the evidence and
publication duties of programs that have stopped (\cref{app:storage}).
\item A \defterm{session program} exists for each fixed pair
$(\beta,r)$ in which the node takes part. It runs only the regional BFT
of that session: the bootstrap of \cref{alg:bootstrap}, then
\cref{alg:steady,alg:recovery}.
\end{itemize}
The execution layer supplies payloads and executable state to the
programs. It takes no part in rotation decisions.

\paragraph{Data flow.}
\Cref{fig:runtime-modules} follows one round of the loop.
\begin{enumerate}[label=(\arabic*)]
\item The Multisig Controller appends control updates to L1. The Publish
Layer submits complete regional results to Espresso.
\item Each regional node, observer and certifier replays both histories
block by block. At the start of each Espresso block it applies the due
atomic updates; then it accepts DATA into $G$ and derives boundaries.
The same rules derive the branch version and its immutable map as in
\cref{sec:runtime-controls}. Verified progress goes back to the
controller.
\item The Session Manager uses the node's replayed state, keeps its
baseline at $\operatorname{last}(G)$, assigns new sessions and records
closures.
\item An Assigned program enters, together with the other members of
its session, once its fixed handoff and exact anchor state are present
(\cref{sec:ready}). The programs of a closed version stop together
(\cref{sec:retirement}) and hand their duties to the Manager.
\item A running program produces complete regional results and pushes
them to the Publish Layer. The Session Manager retries the publication
and sends terminal results directly to the selected successor region.
This closes the loop.
\end{enumerate}

\paragraph{Authority rules.}
Three rules keep these roles apart.
\begin{enumerate}[label=(A\arabic*)]
\item Only an L1 control update that takes effect changes eligibility or
the branch version (\cref{prop:explicit-force}).
\item Only the public replay rules change $G$. A program's local phase
neither grants nor revokes permission to extend $G$.
\item Only the Manager's Enter call starts a regional engine, and only
for a fixed handoff. Baseline execution and preparation grant no vote.
\end{enumerate}
A monitoring report, a local timeout or a missing output has no control
effect. An idle region keeps following L1 and Espresso and keeps its
baseline state current, so it can take over at any time.

\begin{figure}[p]\centering
% Cross-region components (Figure: fig:runtime-modules).
% Protocol numbers come from \cref, so they follow the document.
\resizebox{\linewidth}{!}{%
\begin{tikzpicture}[
 x=1cm,y=-1cm,font=\small,
 box/.style={draw=rrBlue!65,rounded corners=3pt,fill=rrBlue!3,
   align=center,inner sep=6pt,text width=2.9cm,minimum height=2.0cm},
 shared/.style={box,draw=rrTeal,fill=rrTeal!4,text width=15.85cm,
   minimum height=1.6cm},
 prog/.style={box,text width=4.38cm,minimum height=2.5cm},
 flow/.style={-{Stealth[length=2.2mm]},semithick},
 lab/.style={fill=white,inner sep=2pt,align=center,font=\footnotesize}
]
\path[use as bounding box] (-1.1,0) rectangle (17.2,15.0);
\node[font=\large\bfseries] at (8.15,.4) {Cross-region data flow on one node};

% Row 1: the two outside services and their logs.
\node[box] (ctrl) at (1.65,2.0) {\textbf{Multisig Controller}\\
 Watch progress\\and reports;\\issue control\\updates.};
\node[box] (l1) at (6.35,2.0) {\textbf{L1}\\
 Authorities and\\the control log.};
\node[box] (esp) at (9.95,2.0) {\textbf{Espresso}\\
 Finalized log with\\timestamps and\\namespace data.};
\node[box] (pub) at (14.65,2.0) {\textbf{Publish Layer}\\
 Leader first,\\then backup senders.};
\draw[flow] (ctrl.east)--node[lab,above=3pt]{control\\updates}(l1.west);
\draw[flow] (pub.west)--node[lab,above=3pt]{EXPORT\\records}(esp.east);

% Row 2: public replay is ordinary node logic, not a separate module.
\node[shared] (manager) at (8.15,5.8) {\textbf{Session Manager:
 \Cref{alg:rotation}}\\
 Directly replays finalized L1 and Espresso histories; observers and
 certifiers apply the same public rules
 (\Cref{alg:global-rule,alg:force,alg:external}).\\
 Views: regional R, global G, control L. Baseline: executable state at
 $\operatorname{last}(G)$.\\
 Assigns new sessions; starts and stops each session in sync with its peers
 (ReadyBundle, close bundle);\\
 retries publication; sends handoff material to the successor region.};
\draw[flow] (l1.south)--(l1.south |- manager.north);
\draw[flow] (esp.south)--(esp.south |- manager.north);
\draw[flow] (manager.west)--(-.25,5.8)--(-.25,2.0)--(ctrl.west);
\node[lab,rotate=90] at (-.25,3.9) {verified progress};

% Row 3: local programs, one per fixed session.
\node[prog,draw=black!45,fill=black!3] (archive) at (2.4,9.5)
 {\textbf{Closing /}\\ \textbf{Terminated} $(\beta,r)$\\
 No fresh actions; timers drain; duties pass to the Manager; tombstone
 kept.};
\node[prog] (new) at (7.9,9.5)
 {\textbf{Assigned $(\beta',r')$}\\
 Handoff and anchor fixed. Starts when Ready holds; all members within
 $\Delta$.};
\node[prog,draw=rrTeal,fill=rrTeal!5] (old) at (13.4,9.5)
 {\textbf{Running $(\beta,r)$}\\
 Regional BFT only. Frozen and recovery stay inside the engine.};
\draw[semithick] (manager.south)--node[lab,right=2pt]{commands,
 events}(8.15,7.75);
\draw[flow] (8.15,7.75)-|(archive.north);
\draw[flow] (8.15,7.75)-|(new.north);
\draw[flow] (8.15,7.75)-|(old.north);

% Row 4: execution layer and the regional engine.
\node[box,draw=black!45,fill=black!3,text width=4.38cm,minimum height=2.3cm]
 (el) at (2.4,13.2)
 {\textbf{Execution layer}\\
 Payloads, executable state and request queues.};
\node[box,draw=rrTeal,fill=rrTeal!4,text width=10.38cm,minimum height=2.3cm]
 (gate) at (10.9,13.2)
 {\textbf{Regional engine: \Cref{alg:bootstrap,alg:steady,alg:recovery}}\\
 STATUS, $2\Delta$ wait, NEWVIEW,\\ ANCHOR QC, then DATA/DRAIN.\\
 Messages only with same-session peers in the region. L1 and Espresso
 references and authorities come from the Manager.};
\draw[flow] (new.south)--node[lab,right=2pt]{Enter}(new.south |- gate.north);
\draw[semithick] (old.south)--node[lab,right=2pt]{its engine}(old.south |- gate.north);
\draw[flow] (el.east)--node[lab,above=2pt]{payloads,\\state}(gate.west);

% Results: one-way push from the engine; the Manager retries.
\draw[flow] (gate.east)--(16.8,13.2)--(16.8,2.0)--(pub.east);
\draw[semithick] (manager.east)--(16.8,5.8);
\node[lab,rotate=90] at (16.8,9.5) {results: one-way push};
\node[lab,rotate=90] at (16.8,4.0) {and Manager retries};
\end{tikzpicture}%
}

\caption{Cross-region data flow. Only L1 control updates change
eligibility or the branch version, only replay changes $G$, and only the
Manager's Enter call starts an engine. The same replay runs at regional
nodes, observers and certifiers.}
\label{fig:runtime-modules}
\end{figure}
\FloatBarrier

\subsection{Runtime use of atomic branch updates}
The atomic branch-update interface, its validity rules and its activation
semantics are defined before global finality in \cref{sec:l1}.  Runtime
components consume only the resulting immutable branch map $S_\beta$,
the executing update and its activation block.  A stale update has no
runtime effect at all.  When an update executes, Protocol~3 preserves
$G$, closes the target version, installs the new map and fixes the new
tail; the Session Manager separately retires local programs.

The committee's policy determines when to issue a branch update.
Because every map change closes the current branch, implementations
should batch operational changes and normally keep at most one
unactivated update in flight.  A stale update's intended map change is
lost and must be resubmitted against the current branch version.
\Cref{app:policy} gives a reference policy.

\subsection{The Session Manager}
\label{sec:session-manager}
The Session Manager decides which sessions the node runs, from which
block each one starts, and when each one stops. It runs no consensus and
makes no public decision. Everything it acts on comes from three views
built from verifiable sources, and its actions follow the fixed rules of
\cref{alg:rotation}.

\paragraph{Three views.}
A view is only a name for a group of objects that the Manager reads.
$\mathcal V_G$ and $\mathcal V_L$ add nothing to the public replay
output; only $\mathcal V_R$ collects objects not defined elsewhere. Given an Espresso block $q$,
\begin{itemize}
\item \textbf{Global view} $\mathcal V_G$. the output
$\mathrm{Out}_q=(G_q,\mathsf{Tail}_q)$ of \cref{alg:global-rule} at Espresso
block $q$: the global chain and the tail that assigns or leaves pending
the session responsible for its next DATA.
\item \textbf{Control view} $\mathcal V_L$. the finalized
L1 state at the same position: the registered authorities
$\tagg{AUTH}_{r,g}$ (\cref{sec:l1}), the control log $\Lambda$, the
current branch version $\beta$ and its immutable status map $S_\beta$
(\cref{def:control-effect}).
\item \textbf{Regional view} $\mathcal V_R$. \emph{Covers:} the
regionally finalized DATA blocks with its complete
receipt (\cref{def:receipt}), grouped by session. 
\end{itemize}

Each node derives $\mathcal V_G$ and $\mathcal V_L$ by applying
\cref{alg:global-rule,alg:force,alg:external}, so every correct node that
reaches the same replay position holds the same two views.

A branch version $\beta'$ is closed in public replay once the prelude
of some block executes a Force on it (\cref{def:control-effect}). A node
learns this from its own replay or, possibly earlier, from a verified
close bundle of another member (\cref{sec:retirement}). Either way it
records a durable \defterm{closed mark} for $\beta'$, and we say that
$\beta'$ is \defterm{marked closed} at the node. The retirement gate and
the entry checks test this mark, not the node's replay position.

\paragraph{Local data.}
Besides the three views, the Manager keeps:
\begin{itemize}
\item two exact execution checkpoints. Write $\mathit{st}(B)$ for the
executable state after block $B$. The \defterm{baseline state}
$\mathit{st}(\operatorname{last}(G))$ is kept by every operational node.
In the region currently expected to take over by regular handoff, the
\defterm{successor state} is $\mathit{st}(\hat B)$, where $\hat B$ is the
latest regionally finalized DATA of the producing session. The baseline
state is determined by public $G$ and is common to nodes at the same
replay position. The successor state may be further ahead and may differ
temporarily across nodes, because direct regional results can arrive at
different times.
\item the closed marks, each with the closure evidence of
\cref{sec:retirement};
\item the \defterm{program table}: for every fixed session in which the
node is a member, the program's state (\cref{sec:session-lifecycle}), its
fixed handoff, a durable entry flag, its entry time $e_v$ and its close
time $c_v$. Records are never deleted; the record of a Terminated
program is a tombstone;
\item the \defterm{duty table}: a durable to-do list for work that
outlives a session program.  It records which close bundles still need
to be sent, which complete results still need publication retries, and
which session data or evidence the Manager must continue to serve, even
if the program has stopped or was never created;
\item the evidence store of \cref{app:storage}, which keeps the exact
bytes of every public object the node may need again.
\end{itemize}

\begin{figure}[htbp]\centering
% The two execution checkpoints maintained by the Session Manager.
\begin{tikzpicture}[
 x=1cm,y=1cm,font=\small,
 block/.style={draw,rounded corners=2pt,minimum width=1.05cm,
   minimum height=.72cm,align=center,inner sep=2pt},
 global/.style={block,draw=rrBlue!75,fill=rrBlue!7},
 regional/.style={block,draw=rrTeal!75,fill=rrTeal!7},
 checkpoint/.style={draw,rounded corners=3pt,align=center,inner sep=4pt,
   minimum width=2.8cm,minimum height=.72cm},
 session/.style={draw,rounded corners=3pt,align=center,inner sep=5pt,
   minimum width=2.65cm,minimum height=.85cm},
 flow/.style={-{Stealth[length=2.2mm]},semithick},
 note/.style={align=center,font=\footnotesize}
]
\node[note,anchor=east] at (.45,3.05) {A's DATA};
\node[note] at (1.25,3.05) {$\cdots$};
\node[global] (b108) at (2.25,3.05) {108};
\node[regional] (b109) at (4.25,3.05) {109};
\node[regional] (b110) at (6.25,3.05) {110};
\draw[semithick] (1.55,3.05)--(b108.west);
\draw[semithick] (b108.east)--(b109.west);
\draw[semithick] (b109.east)--(b110.west);

\node[note,rrBlue] at (2.25,3.85)
  {in $G$ through 108};
\node[note,rrTeal] at (5.25,3.85)
  {regional only: 109--110};

\node[checkpoint,draw=rrBlue!75,fill=rrBlue!7] (sg) at (2.25,1.75)
  {baseline: $\mathit{st}(108)$};
\node[checkpoint,draw=rrTeal!75,fill=rrTeal!7] (st) at (6.25,1.75)
  {successor: $\mathit{st}(110)$};
\draw[flow,rrBlue] (b108.south)--(sg.north);
\draw[flow,rrTeal] (b110.south)--(st.north);

\node[session,draw=rrBlue!75,fill=rrBlue!7] (forced) at (2.25,.15)
  {Force successor $(8,21)$\\first new DATA: $109'$};
\node[session,draw=rrTeal!75,fill=rrTeal!7] (regular) at (6.25,.15)
  {regular successor $(7,21)$\\first DATA: 111};
\draw[flow,rrBlue] (sg.south)--node[note,right]{Force keeps $G$}(forced.north);
\draw[flow,rrTeal] (st.south)--node[note,right]{terminal handoff}(regular.north);
\end{tikzpicture}

\caption{The two Session Manager checkpoints in the running example.}
\label{fig:manager-baselines}
\end{figure}

\begin{protocol}[H]{alg:rotation}{Session Manager rules}
Each correct regional replica runs these rules. It handles one event at
a time; retrieval and execution may run in the background. Below, $q$
is the node's current replay position.
\begin{enumerate}[label=\textbf{\arabic*)},itemsep=3pt]
\item \textbf{Baseline.} Whenever $\mathcal V_G$ appends a DATA block,
execute it, so that the baseline state stays
$\mathit{st}(\operatorname{last}(G))$. If the node's region is the one
that $\tagg{NextLive}$ selects after the producing region under that
session's fixed $S_\beta$ and target authority,
also execute the producing session's regionally finalized DATA from
$\mathcal V_R$ into the successor state. Baseline work signs nothing,
holds no lock and grants no entry.
\item \textbf{Assign.} Create a session program in state Assigned when
one of the two cases below selects the node's region. Each case fixes
the session id, its anchor and its Espresso lower bound.
 \begin{enumerate}[label=(\alph*),itemsep=2pt]
 \item \emph{Force handoff.} $\mathcal V_G$ has
 $\mathsf{Tail}=\tagg{NEXT}(\nu)$ with
 $\nu=(\beta,r',i,A,H,q_f,\tagg{FORCE}(c_f,q_f))$ and $i$ equal to the
 node's region (the next-block selection of \cref{alg:global-rule}). The session id is
 $\operatorname{sid}(\nu)=(\beta,r')$: $\beta$ is the branch version
 opened by the Force at $J_{q_f}$, and $r'$ is the target index fixed by
 Step~1 from the former $\tagg{NEXT}$ tail's accepted count, or kept from
 the former $\tagg{PENDING}$ tail (\cref{sec:retry-slot}). The anchor is
 $A=\operatorname{last}(G)$ at the start of $J_{q_f}$, and
 $q^{\rm open}_{\beta,r'}=q_f$.
 \item \emph{Regular handoff.} $\mathcal V_R$ holds a complete receipt
 for the terminal DATA $T_{\beta,r}$ of a fixed session $(\beta,r)$ in
 region $i_{\beta,r}$, $\beta$ is not marked closed at the node, and the
 fixed branch map selects the node's region:
 $\tagg{NextLive}(i_{\beta,r},S_\beta;r+1)$ returns that region. The
 session id is
 $(\beta,r+1)$, the anchor is $T_{\beta,r}$, and
 \[
 q^{\rm open}_{\beta,r+1}=
 \max\{\operatorname{height}(\xi(T_{\beta,r})),
        q^{\rm open}_{\beta,r}+1\}.
 \]
 This case does not wait for $T_{\beta,r}$ to enter $G$ or for a header
 at the new reference floor. That floor constrains the successor's first
 DATA; it does not delay assignment or entry.
 \end{enumerate}
 In both cases the authority is $\tagg{AUTH}_{r',0}$ of the target turn
 index $r'$, and $\lambda^{\rm open}$ is the larger of the anchor's L1
 reference height and that authority's publication height. Verify the
 handoff with \tagg{VerifyHandoffHistory} (\cref{sec:verify-handoff})
 and record it in the program table. Each pair is assigned at most once;
 a different handoff for an assigned pair is rejected. The genesis
 session is assigned when $\mathcal V_G$ has
 $\mathsf{Tail}=\tagg{NEXT}(\nu)$ with cause \tagg{GENESIS} and
 $i(\nu)$ equal to the node's region. A
 verified ReadyBundle for a pair the node has not yet assigned assigns
 it in the same way and then enters at once (\cref{sec:ready}).
\item \textbf{Start.} When an Assigned program is Ready
(\cref{def:ready}), start it by the synchronous start of
\cref{sec:ready}. Process every valid ReadyBundle under the same rules.
\item \textbf{Stop.} When a branch version closes, run the synchronous
stop of \cref{sec:retirement} for every known fixed session of that
version, whatever its state, and deliver Close to its local program if
one exists.
\item \textbf{Serve.} Give every Running program the L1 and Espresso
inputs it needs (\cref{sec:session-lifecycle}). Offer every complete
regional result to the Publish Layer until it is in Espresso, as a
backup to the program's own push. Send each terminal result directly to
the selected successor region. Keep the data, publication and relay
duties of programs that have stopped.
\item \textbf{Persist.} Record every state change, entry, closure and
duty durably before exposing it or dispatching a signature. A repeated
event resumes the recorded action; it never starts a second one.
\end{enumerate}
\end{protocol}

Three invariants follow from these rules.
\begin{itemize}
\item \emph{Replay only.} $\beta$, $G$ and the current immutable map
change only through replay; no Manager rule increments $\beta$, chooses
a cutoff or edits $G$.
\item \emph{Fixed is final.} A fixed handoff, its pair and its anchor
never change, and every member enters each pair at most once.
\item \emph{Idempotence.} Duplicates, retries and restarts repeat no
logical action. After a crash the Manager restores its records,
closures and duties before it dispatches any callback or signature.
\end{itemize}

\Cref{tab:manager-inputs} lists the events the rules react to. The
Manager's outputs are calls to its programs, ReadyBundles and close
bundles to the other members of a session, results to the Publish Layer
and terminal results to the successor region.

\begin{table}[htbp]\centering\small
\caption{Session Manager events.}
\label{tab:manager-inputs}
\begin{tabularx}{\linewidth}{@{}>{\raggedright\arraybackslash}p{.24\linewidth}>{\raggedright\arraybackslash}p{.24\linewidth}Y@{}}\toprule
Event & From & Handling\\\midrule
Replay step & Finalized L1 and Espresso history & Update $\mathcal V_G$ and $\mathcal V_L$;
then Baseline, Assign~(a) and Stop.\\
Regional result & Own programs, other regions, Espresso & Validate the
DATA, its complete receipt and its ancestry; update $\mathcal V_R$; then
Baseline, Assign~(b) for a terminal result, and Serve.\\
ReadyBundle & A member of the same session & Assign if the pair is new; then Start (\cref{sec:ready}).\\
Close bundle, or any packet of a closed session & A member of the same
session & Stop (\cref{sec:retirement}).\\
Program event & Local program & Update the program table: Entered,
Result, Completed or Drained.\\
Callback & Retrieval or execution & Check the named target and anchor;
update only the baseline state or the matching program.\\\bottomrule
\end{tabularx}
\end{table}

\paragraph{Example: assigning sessions.}
Return to the running example (\cref{fig:manager-baselines}). The
regions are A, B, C and D, in this order. A's session $(7,20)$ has
anchor 100 and quota 10, so $T_{7,20}=110$. It opened at
$q^{\rm open}_{7,20}=480$, and DATA 110 carries
$\xi(110)=(500,\Hash(J_{500}))$. The fixed map $S_7$ sets all four
regions ON.
\begin{itemize}
\item \emph{Regular handoff.} As soon as a node of B holds the complete
receipt for 110 in $\mathcal V_R$, it computes
$q^{\rm open}_{7,21}=\max\{500,481\}=500$ and
$\tagg{NextLive}(\text{A},S_7;21)=\text{B}$. Rule~2(b) assigns
$(7,21)$ with anchor 110, authority $\tagg{AUTH}_{21,0}$ and
$q^{\rm open}_{7,21}=500$, while $G$ is still 108. The nodes of C and D
compute the same selection and assign nothing. Two members of B that
receive the receipt at different times also assign at different times;
the synchronous start of \cref{sec:ready} still makes them enter within
$\Delta$ of each other.
\item \emph{Force handoff.} Update~41 of \cref{sec:runtime-controls}
takes effect at $J_{501}$ while $G=108$ and the export session is
$(7,20)$ with accepted count $a=8$. Replay closes version 7, finds
$\tagg{NextLive}(\text{A},S_8;21)=\text{B}$ because update 41 installs
$S_8$ with A OFF, and
sets $\mathsf{Tail}=\tagg{NEXT}(\nu_B)$ with
$\nu_B=(8,21,\text{B},108,118,501,\tagg{FORCE}(41,501))$. Rule~2(a) at B assigns
$(8,21)$ with anchor 108 and $q^{\rm open}_{8,21}=501$, and Stop closes
$(7,21)$. Had $(7,20)$ added no DATA to $G$ ($a=0$, $G=100$), the
session id would be $(8,20)$ with anchor 100. Had every region been OFF
at $J_{501}$, replay would keep a pending boundary with target index 21,
and no node would assign anything until another successful atomic update
created a later version with an eligible region.
\end{itemize}

\paragraph{Why two checkpoints instead of Force candidates.}
A Force always keeps $\operatorname{last}(G)$ at the start of its
activating block (\cref{alg:global-rule}). A node that keeps its baseline state current
therefore already holds the state at every possible Force anchor,
whichever region and target index the Force selects, and no rule needs
to predict Force candidates. A regular handoff continues after the
predecessor's terminal DATA, which the successor state reaches as soon as that DATA is
regionally final.  In the running example of
\cref{fig:manager-baselines}, the baseline state is $\mathit{st}(108)$
and the successor state is $\mathit{st}(110)$.  A Force uses the former;
the regular successor uses the latter.  Keeping the baseline state is ordinary chain execution rather than
preparation for a guessed Force: each region must be able to take over
at the public prefix at any time.  \Cref{sec:runtime-races} works through
the resulting races.

\subsection{The session program}
\label{sec:session-lifecycle}
A \defterm{session program} runs one fixed session $(\beta,r)$ at one
member of its region. It runs the regional engine of
\cref{sec:protocols}: the bootstrap of \cref{alg:bootstrap}, then
\cref{alg:steady,alg:recovery}. It holds only that engine's state: the
view $(g,t)$, locks, votes, QCs, timers, the committed log and the
executable state. Programs of different sessions share nothing, even
when they reuse a turn index.

\paragraph{Communication.}
A program talks to three parties only.
\begin{center}\small
\begin{tabularx}{\linewidth}{@{}>{\raggedright\arraybackslash}p{.3\linewidth}>{\raggedright\arraybackslash}p{.12\linewidth}Y@{}}\toprule
Party & Direction & Content\\\midrule
Programs of the same $(\beta,r)$ in the region & Both & Regional BFT
messages: STATUS, NEWVIEW, proposals, votes, FINAL, BLAME, quit proofs
and receipts.\\
The local Session Manager & Both & Calls and inputs down, events up
(below).\\
The Publish Layer & Out only & The session's own complete regional
results.\\\bottomrule
\end{tabularx}
\end{center}
A program never reads L1 or Espresso and never talks to another session
or region. The ReadyBundle and the close bundle travel between Managers,
because the receiving program may not exist yet, or any more.

\paragraph{L1 and Espresso as inputs.}
Regional BFT still depends on public history. Every block carries an L1
reference $\rho$ and an Espresso reference $\xi$
(\cref{sec:blocks}), and recovery entry needs the finalized recovery
authority $\tagg{AUTH}_{r,g}$ (\cref{sec:protocols}). The Manager
supplies all of them from $\mathcal V_L$ and $\mathcal V_G$:
\begin{itemize}
\item \emph{at Enter}: the fixed handoff, the anchor $A_{\beta,r,0,0}$
with its commitment evidence, the executable state after it,
$\tagg{AUTH}_{r,0}$, the quota $\TB_r$ and the lower bounds
$q^{\rm open}_{\beta,r}$ and $\lambda^{\rm open}_{\beta,r}$;
\item \emph{references}: a current finalized L1 reference with its
registry proof and a finalized Espresso header at or above the session's
reference floor, for the leader's new blocks, and the authenticated history
needed to check the references in a peer's block;
\item \emph{authorities}: each recovery authority $\tagg{AUTH}_{r,g}$
with $g>0$, in order, once it is finalized. A change of branch version
does not reset this sequence.
\end{itemize}
These inputs are facts with proofs; none of them is an instruction. A
missing input makes the program wait, as a missing dependency does in
\cref{sec:admission}.

\paragraph{Results and receipts.}
\label{sec:retirement-receipt-relay}
When a program first holds a complete regional result of its own
session, it pushes the result to the Publish Layer and reports it to the
Manager. The push is one way: the program neither waits for an answer
nor retries; the Manager's Serve rule does that. $T_{\rm pub}$ starts
when the complete result is available.

The first time a correct member holds a complete valid receipt of its
own session, it multicasts the receipt and all its validation
dependencies to the session's fixed membership; a terminal receipt goes
out no later than the local Completion event. Deduplication by session
and DATA hash suppresses repeated encodings or signer sets, not this
first complete relay. Recipients apply an own-session terminal receipt
under Completion even while closing, independently of its generation and
before a simultaneous timer. An Assigned program caches it; after Exit
the Manager archives it and keeps the relay duty. The relay forwards
existing signatures; it adds no vote or certificate. The regional
complete-message bound applies to the actual relay, not to first forming
a receipt.

\paragraph{States.}
A program is in one of four states (\cref{fig:program-states}). The
table relates them to the engine phases of \cref{sec:admission}.
\begin{center}\small
\begin{tabularx}{\linewidth}{@{}>{\raggedright\arraybackslash}p{.15\linewidth}Y>{\raggedright\arraybackslash}p{.25\linewidth}@{}}\toprule
State & Meaning & Engine phase\\\midrule
Assigned & Its handoff is fixed and it has not entered. It holds no vote,
lock or timer. & None yet\\
Running & Entered; production or recovery is under way. & Active or
Frozen\\
Closing & Produces nothing new: Completion has run, or its version is
closed. Running timers drain and duties pass to the Manager. &
Completed, or Active or Frozen behind the closed gate\\
Terminated & Exited. A tombstone stops late messages from restarting
it. & None\\\bottomrule
\end{tabularx}
\end{center}
\begin{itemize}
\item \emph{Assign} creates an Assigned program (rule~2 of \cref{alg:rotation}).
\item \emph{Enter} moves Assigned to Running, at every correct member
within $\Delta$ of the first (\cref{sec:ready}).
\item \emph{Completion} (\cref{alg:steady}) or \emph{Close}
(\cref{sec:retirement}) moves Running to Closing.
\item \emph{Exit} moves Closing to Terminated (\cref{sec:retirement}).
\item \emph{Cancel} moves Assigned to Terminated when its version closes.
\end{itemize}
Quit and recovery move the engine between Active and Frozen inside
Running; the Manager sees no change of state. A closed program keeps its
engine phase until a genuine Quit, Completion or Exit: Close is neither
a quit proof nor a terminal receipt. Every transition is persisted
before its effects are exposed.

\begin{figure}[htbp]\centering
% Session program lifecycle (Figure: fig:program-states).
\begin{tikzpicture}[
 x=1cm,y=-1cm,font=\small,
 st/.style={draw=rrBlue!65,rounded corners=3pt,fill=rrBlue!3,
   align=center,inner sep=5pt,text width=2.4cm,minimum height=1.25cm},
 flow/.style={-{Stealth[length=2.2mm]},semithick},
 lab/.style={fill=white,inner sep=2pt,align=center,font=\footnotesize}
]
\node[st] (asg) at (0,2.2) {\textbf{Assigned}\\ handoff fixed};
\node[st,draw=rrTeal,fill=rrTeal!5] (run) at (5.6,2.2)
 {\textbf{Running}\\ engine Active\\or Frozen};
\node[st,draw=rrTeal,fill=rrTeal!5] (cls) at (11.2,2.2)
 {\textbf{Closing}\\ no fresh actions};
\node[st,draw=black!45,fill=black!3] (ter) at (11.2,5.0)
 {\textbf{Terminated}\\ tombstone};
\draw[flow] (0,0.4)--node[lab,right=2pt]{Assign (\cref{alg:rotation})}(asg.north);
\draw[flow] ([xshift=-5mm]run.north) .. controls +(-.4,-1.3) and +(.4,-1.3)
 .. node[lab,above=1pt]{Quit, recovery (\cref{alg:steady,alg:recovery})}
 ([xshift=5mm]run.north);
\draw[flow] (asg.east)--node[lab,above=3pt]{Enter, in sync\\(\cref{sec:ready})}(run.west);
\draw[flow] (run.east)--node[lab,above=3pt]{Completion or\\Close (\cref{sec:retirement})}(cls.west);
\draw[flow] (cls.south)--node[lab,right=2pt]{Exit (\cref{sec:retirement})}(ter.north);
\draw[flow,dashed] (asg.south)--(0,5.0)--node[lab,above=2pt]{Cancel: version
 closed (\cref{sec:retirement})}(ter.west);
\end{tikzpicture}

\caption{Session program states. A program exists only once its handoff
is fixed; the engine's Active and Frozen phases both lie inside
Running.}
\label{fig:program-states}
\end{figure}

\paragraph{Interface with the Manager.}
\begin{center}\small
\begin{tabularx}{\linewidth}{@{}>{\raggedright\arraybackslash}p{.2\linewidth}>{\raggedright\arraybackslash}p{.27\linewidth}Y@{}}\toprule
Direction & Call or event & Effect\\\midrule
Manager to program & Enter & Install the fixed start and enter
$(\beta,r,0,0)$ (\cref{alg:bootstrap}).\\
& References, Authority & The inputs listed above.\\
& Close & Shut the fresh-action gate (\cref{sec:retirement}).\\
& Release & Confirm that the duty table holds every duty; allow Exit.\\
Program to Manager & Entered, Result, Completed, Drained & Entry time; a
complete result; Completion; all timers done.\\\bottomrule
\end{tabularx}
\end{center}
Enter and Close are local calls in the node's single event order, not
network messages. Each takes effect before the program's next action,
and $e_v$ and $c_v$ are recorded at that moment.

\paragraph{What the program assumes.}
The program relies on two guarantees that only the Managers of its
region can give.
\begin{itemize}
\item \emph{Synchronous start}: if the first correct member enters at
$e$, every correct member that has not closed the session enters by
$e+\Delta$ (\cref{prop:sync-start}).
\item \emph{Synchronous stop}: if the first correct member closes the
session at $c$, every correct member closes it by $c+\Delta$
(\cref{prop:retirement-drain}).
\end{itemize}
Part~I's vote-window argument needs the first: a correct member that
votes at $s$ has entered, so every correct member can process the
proposal by $s+\Delta$. The drain bound needs the second.
\Cref{sec:ready,sec:retirement} show how bundle relays provide both.

\subsection{Fixing a handoff}
\label{sec:witness}
A handoff record, or \defterm{witness}, answers three questions that the
successor and every verifier must agree on: why the handoff happened,
who takes over, and which block the new session continues after. It is
not an extra round of consensus. Once fixed, a handoff never changes:
its target pair, region, anchor, source context, branch map, authority and
reference bounds stay as recorded. Changing eligibility requires an
atomic update that closes the old branch; it does not rewrite a regular
handoff already fixed inside that branch.

\paragraph{Choosing the next region: NextLive.}
$\tagg{NextLive}(i,S_\beta;r')$ searches the other regions after $i$ in
the fixed cyclic order $\mathcal L_R$ that are ON in $S_\beta$ and have
positive total weight under $\tagg{AUTH}_{r',0}$. It
returns $i$ itself only if no other region is eligible and $i$ still is.
Otherwise it returns $\bot$. A synthetic source $\bot$ searches from the
beginning of the order.

REGULAR handoffs and atomic branch updates both use NextLive, but with
different branch maps: a regular handoff uses the unchanged $S_\beta$,
whereas a successful update uses the newly installed $S_{\beta+1}$. A
Force does not implicitly disable its source region. A physical
region may be selected again, but its new pair has fresh locks, votes and
timers.


\paragraph{The target index of a Force.}
\label{sec:retry-slot}
If $\mathsf{Tail}=\tagg{NEXT}(\nu)$, write
$\operatorname{sid}(\nu)=(\beta,r)$, $i=i(\nu)$ and
$a=a(G,\nu)$.  This is the number of that session's DATA accepted into
$G$; it is not a regional height or a number of raw publications. The
target index is
\[
 r'=
 \begin{cases}
 r,&a=0,\\
 r+1,&a>0.
 \end{cases}
\]
In words: if the export session has put nothing into $G$, the Force
retries the same authority index; otherwise it moves to the next one.
This rule chooses only the authority index. It does not decide whether a
Force is allowed, and the new region may be a different physical region
even when $a=0$. Both cases change the branch version to $\beta+1$ and
start their new engine at $g=t=0$. The retained anchor is exactly $G$ at
the start of the activating block, with its original producing pair,
height, hash and commitment evidence.

In the running example, $G=108$ and the next DATA 109 belongs to A's
$(7,20)$, which has already put 101--108 into $G$. So $r'=21$, even though
B's regular $(7,21)$ already exists and B has produced up to 115. The
export session is chosen by what $G$ needs next, not by the region that
runs furthest ahead.

An unassigned boundary is
$\mathsf{Tail}=\tagg{PENDING}(\omega)$, where
$\omega=(\mathit{cause},\beta,r',i_{\rm src},A,q^{\rm open})$. It stores its source region $i_{\rm src}$ and its
target index $r'$. A later Force changes the branch version and the
cause but keeps that target index; it does not recompute it. A regular
boundary stores $i_{\rm src}=i(\nu)$ and $r'=r(\nu)+1$; genesis stores
$i_{\rm src}=\bot$ and $r'=0$. This handles repeated Forces
before any new duty is assigned, without creating empty sessions.

Regular completion opens $(\beta,r+1)$ from exactly $T_{\beta,r}$. A
zero-output retry may reuse the number $r$, even at another physical
region, but never rebinds $\tagg{AUTH}_{r,g}$. For a fixed region and
$(r,g,t)$, the recorded leader does not depend on $\beta$: the branch
version separates message and engine identities, not the schedule. The
fixed snapshot and quota stay those of the record. No old Quit, vote,
lock or timer becomes part of a new engine's state.


\paragraph{Fixed-map selection and the reference floor.}
All nodes select from the map attached to the target branch version, not
from the map or wall-clock view they happen to hold when a receipt
arrives. For a REGULAR boundary from $(\beta,r)$, compute
\[
 i'=\tagg{NextLive}(i_{\beta,r},S_\beta;r+1),\qquad
 q^{\rm open}_{\beta,r+1}
 =\max\{\operatorname{height}(\xi(T_{\beta,r})),
         q^{\rm open}_{\beta,r}+1\}.
\]
For a successful atomic update at $J_{q_f}$, replay first installs
$S_{\beta+1}$, selects under that map and sets the new session's opening
floor to $q_f$. Missing source or authority evidence causes a wait; it
never turns an unknown candidate into an ineligible one.

The opening floor is not a selection height. A REGULAR successor is
fixed as soon as the direct terminal result and its dependencies are
verified, without waiting for the terminal DATA to enter $G$ and without
waiting for a later Espresso header. If the numerical floor names a
future height, the successor may still enter; its first DATA proposal
waits until it can carry a finalized reference at or above that floor.
A speculative regular handoff does not itself append DATA to $G$.

An EXPORT can span several sessions when every successor DATA already
carries a reference below the containing block and all preceding DATA
and receipts are present. If the required reference is not yet available,
that successor DATA is simply not valid yet; this does not delay global
acceptance of the predecessor's terminal DATA.


\paragraph{Handoff packet.}
A handoff packet contains the fields of \cref{tab:handoff-packet}.
References may name authenticated cached objects or objects still to be
fetched; a reference alone does not establish Ready.
\begin{table}[htbp]\centering\small
\caption{Handoff packet fields.}
\label{tab:handoff-packet}
\begin{tabularx}{\linewidth}{@{}>{\raggedright\arraybackslash}p{.26\linewidth}Y@{}}\toprule
Field & Type and meaning\\\midrule
Domain and kind & Chain/protocol domain; REGULAR or MSIG; genesis uses its trusted initial boundary.\\
Source context & Predecessor pair $(\beta_m,r_m)$ and source region
$i_{\rm src}$, including the accepted count when assigned and the stored target
index when unassigned.\\
Target & Target pair $(\beta',r')$, the reason for its index, and the
selected region $i_{\beta',r'}$.\\
Anchor & Producing pair, DATA height, hash, full DATA or reference,
and original direct or covering commitment evidence.\\
Authority & Finalized $\tagg{AUTH}_{r',0}$, its fixed configuration,
quota, and publication-height proof.\\
Selection and bounds & The immutable target map $S_{\beta'}$, the
NextLive result, $q^{\rm open}_{\beta',r'}$ and
$\lambda^{\rm open}_{\beta',r'}$.\\
Cause and evidence & Mode-specific cause; trusted or validated replay
checkpoint, complete required L1/Espresso history and regional proofs.\\\bottomrule
\end{tabularx}
\end{table}

For the Force in the running example, the packet says: the cause is a
valid MSIG effective at a given Espresso block; it closes version 7 and
opens version 8; the export session was A's $(7,20)$, owing DATA 109; the
target is B's $(8,21)$, selected under the newly installed $S_8$; the anchor is
DATA 108 with its exact hash, data and receipt, still produced by
$(7,20)$; and the authority is the finalized $\tagg{AUTH}_{21,0}$. A
regular packet has the same shape, but its cause is the terminal receipt
of the previous session, and its anchor is that terminal DATA.

Two checks use the packet. \tagg{VerifyHandoffHistory} asks whether the
handoff is valid \emph{as history}. \tagg{CanEnter} asks whether a node
may enter \emph{now}. For example, after a Force closes version 7, DATA
108 of $(7,20)$ is still valid history and can anchor $(8,21)$, but that
does not let anyone re-enter $(7,20)$.


\paragraph{Handoff verification.}
\label{sec:verify-handoff}
The pure procedure $\tagg{VerifyHandoffHistory}(P)$ checks a packet $P$
against an authorized replay checkpoint and the required source history.
It returns the historical fixed fields, \emph{invalid},
or \emph{needs evidence}. It authenticates what the boundary established.
It does not require the branch version to be open now, and it grants no
entry permission. Neither procedure changes public replay.
\begin{enumerate}[label=(\roman*),itemsep=2pt]
\item Check the domain, field types, and source and target names. Replay
authenticated compatible history from the checkpoint through all positions needed by
the cause and the incoming boundaries. Verify
the complete applicable controls, not only the named control. Consult
validated caches when available. A historical receipt from a producing
session may use \tagg{CheckFINAL} for intrinsic input validity; this check
does not fetch that producer's optional recovery-admission path. It still
authenticates each fixed initial snapshot and the exact DATA ancestry. The
target's executable state and its own regional admission are verified
separately.
\item For REGULAR, with predecessor session $(\beta_m,r_m)$, validate
its fixed boundary, initial authority, exact quota and terminal DATA
$h(T_{\beta_m,r_m})=h(A_{\beta_m,r_m,0,0})+\TB_{r_m}$. Validate its terminal
FINAL kernel, actual ancestry and signed $\xi$, including the separate
historical source and boundary checks. Reconstruct the immutable map
$S_{\beta_m}$. Require $\beta'=\beta_m$ and
$r'=r_m+1$. The anchor must equal the terminal DATA. Validate the
predecessor's own boundary recursively in the same historical mode;
current closure does not erase its original evidence.
\item For MSIG, verify its authorized finalized L1 record, control
sequence and first-crossing header $J_{q(\theta_a)}$. Replay the full
prelude and check that this is the first due instruction whose target
matches the branch version at the start of the block. A stale or already
consumed instruction cannot form another boundary. No proof of missing
blocks is required. Take the captured source context, the old and new
branch versions, the retained prefix and the boundary from this unique
transition. Require $\beta'=\beta_0+1$ for the version $\beta_0$ at the start of the block, the target index of
\cref{sec:retry-slot}, and anchor equality with $G$ at the start of the
block, including original provenance. Verify the zero or positive
accepted count directly; do not infer it from the anchor's producer. For
an unassigned boundary, verify its stored target index and any Forces in
between. A later Force removes entry permission without invalidating the
historical evidence of the earlier boundary.
\item Recompute the target with NextLive under the designated target
authority and the immutable map of the target branch. For REGULAR, use
$S_{\beta_m}$ and set
$q^{\rm open}_{\beta',r'}=\max\{\operatorname{height}(\xi(T_{\beta_m,r_m})),
q^{\rm open}_{\beta_m,r_m}+1\}$. For MSIG, use the executing update's
complete $S_{\rm new}$ and set $q^{\rm open}_{\beta',r'}=q(\theta_a)$.
Check the target region and immutable target authority, and set
$\lambda^{\rm open}_{\beta',r'}$ to the maximum of the anchor's L1
reference height and the initial authority's publication height. Validate
the anchor's original direct or
covering commitment evidence under \cref{sec:coverage,sec:receipt-kernel},
together with its complete block and required ancestry. Do not relabel
it with the target pair.
\item Compare the derived fixed fields with $P$ and with any existing
fixed record for its target pair. Equal duplicates share that record;
different fixed fields for the same pair are rejected. A packet cannot
replace an installed handoff. This step records historical facts, not a
claim that the target may still enter.
\end{enumerate}


\paragraph{Baseline state and races.}
\label{sec:runtime-races}
Return to the running example. A's session $(7,20)$ ends at 110, B
produces in $(7,21)$ from 111, and $G=108$. Every operational node keeps
its baseline state at $\mathit{st}(108)$. B's nodes, as A's expected
regular successor, also kept the successor state up to $\mathit{st}(110)$.
When A's terminal receipt arrives, Assign creates $(7,21)$ with anchor 110,
and the successor state already matches. Later a Force keeps 108 and
selects B for $(8,21)$: Assign creates that program with anchor 108, and
the baseline state already matches. Had $G$ reached 112 first, with D
selected for $(8,22)$, D's baseline state would be $\mathit{st}(112)$ in
the same way. No
region prepares for a predicted Force, and a pending update needs no
action before it takes effect. If a node's execution lags its replay,
Ready waits until the baseline state reaches the anchor; \cref{ass:ready-service}
bounds that wait.

Races between the two kinds of handoff resolve as follows. A terminal
receipt received before the Force can fix and start a regular successor
in the old branch version; once the node verifies the Force, Stop closes
that program. If the Force is verified first, a late receipt is kept as
regional evidence and cannot fix a handoff in the closed version. Nodes
may briefly hold different programs, but the same replay position, with
complete evidence, fixes the same Force boundary everywhere. Every state
computation keeps the anchor it was computed for: state after 110
cannot make a program anchored at 108 or 112 Ready, and duplicate
receipts or callbacks cannot fix or enter a pair twice. A stale atomic
update changes nothing: it neither closes the new branch version nor
installs its proposed map.

\subsection{Synchronous start}
\label{sec:ready}
An Assigned program starts when two things hold: the node has
everything it needs to continue after the anchor, and the session may
still start now. All correct members of the session then start within
$\Delta$ of one another, because the first to start sends the others a
bundle that lets each of them start in a single step.

\paragraph{Permission to start now.}
\label{sec:can-enter}
$\tagg{CanEnter}(P)$ checks a packet $P$ whose historical check has
passed against the Session Manager's current verified public state. It
requires that:
\begin{itemize}
\item the target branch version is current and not marked closed at
the node;
\item the packet is the fixed boundary that still applies, and this
replica's region is the selected one;
\item the local record has no earlier initial entry, no delivered Close
and no Terminated tombstone.
\end{itemize}
A closed mark blocks entry even before a complete per-session $c_v$ can
be set; context resolution and the closure relay still run before the
refusal. Ready also checks the exact anchor state and the dependencies
below. A duplicate packet for an already entered record is serviced
without another Enter.

All initial entry paths use both checks. Receipt, ancestry and
retained-anchor validation use historical validation only, and global
EXPORT acceptance separately checks current append permission. A cache
miss goes through the same historical checks as a cache hit; no earlier
cached verdict is assumed.

A Force without an eligible successor still has a fixed cause, anchor,
source context and new branch version. Replay keeps that
pending boundary. It has no complete selected packet, and no program can
become Ready on that branch version. Only a later successful atomic
update can supersede it, install a new map and select a target while
keeping the stored target index and predecessor context. The synthetic genesis boundary is verified
against the initial checkpoint and the selection rule of \cref{sec:l1};
it claims no REGULAR or MSIG cause.


\begin{definition}[\defterm{Ready}]
\label{def:ready}
An Assigned program for $(\beta,r)$ is \defterm{Ready} exactly when:
\begin{enumerate}[label=(\roman*)]
\item historical handoff verification fixes this target pair and selects
the replica's region, or the authenticated genesis rule does so; and
the separate current-entry check $\tagg{CanEnter}$ also passes;
\item its exact anchor's original producing pair, height, hash, complete
DATA and commitment evidence have been validated;
\item the native execution state needed to continue after this anchor
has been recovered and verified against that same anchor;
\item the matching finalized authority, fixed quota, immutable branch
map, separate reference lower bounds and all evidence from which those
bounds are derived are available. If the numerical Espresso floor is
above the terminal reference, no header at that floor is required for
entry; the first DATA will require one;
\item the Session Manager's processed public history has not closed or
superseded the handoff; and
\item the durable record has never performed initial Enter and has no
Terminated tombstone.
\end{enumerate}
\end{definition}
In the running example, the new B must: confirm that it is selected and
version 8 is open; hold the exact DATA 108, with its producer $(7,20)$,
hash and receipt; have the execution state after 108 (state after 110, or
after another branch, does not qualify); hold $\tagg{AUTH}_{21,0}$, the
quota and the opening references; and confirm that no later Force has
replaced the handoff and that it has never entered $(8,21)$ before.

With the baseline of \cref{alg:rotation}, condition (iii) usually holds
as soon as the handoff is fixed. A Force anchor is
$\operatorname{last}(G)$, which the baseline state already reaches; a
regular anchor is the predecessor's terminal DATA, which the successor
state already reaches. At
Enter the program takes its own copy of that state, and the baseline
keeps following $G$. State is checked through the native validation
interface: state after 110 cannot make a program anchored at 112 Ready.
An announced update that has not taken effect fixes nothing and starts
nothing.

\paragraph{How long preparation takes.}
The preparation bound concerns one question: once public replay has
assigned the next export session to a region, how long until that region
is ready?

\begin{assumption}[Bounded preparation for a selected live region]
\label{ass:ready-service}
Let $J_o$ be the canonical block in which replay opens an assigned
export session, and let $t_0=t_E(o)$ be its finalization time.
Suppose the selected region stays operational under its fixed snapshot:
its correct members stay online, the regional complete-message contract
holds, and the declared input, load and state-service conditions hold.
Suppose also that the selected handoff stays applicable throughout
preparation. Then by $t_0+T_{\rm ready}$, each correct target member has
either already performed initial Enter, or obtained and validated the
exact anchor DATA, commitment evidence, matching executable state,
finalized authority and public context needed for its guarded initial
Enter.
\end{assumption}
The budget includes observing the opening block, retrieval, validation
and state installation; its clock does not wait until authority or data
has arrived. A regular successor that already entered through direct
delivery meets the obligation earlier. A Force that supersedes the
handoff ends the obligation for that target; not every Force leaves time
for a full preparation interval.

The bound applies to the declared, continuously served configuration and
state interface, not to arbitrary offline catch-up. It does not, by
itself, bound entry skew among members, require new output, or restrict
the next Force. $T_{\rm ready}$ is measured in real time, while controls
use Espresso timestamps and the advance-announcement rule, so no new
two-sided Espresso clock bound is involved. The paragraph below says where
$t_0$ lies for each kind of handoff.


On a node whose execution keeps up with its replay, the baseline makes
this wait short for a Force: the state at the anchor is already
present. The assumption does not rely on that.


\paragraph{Public assignment and the preparation clock.}
Public replay assigns the next export session independently of local
program phases. The \emph{public opening position} of that assignment is:
\begin{itemize}
\item for a normal successor that is selectable at once, the position of
the EXPORT that delivers the terminal DATA;
\item for a forced successor that is selected at once, the triggering
block;
\end{itemize}
A late proof of an earlier selection does not move this historical
position to the local retrieval time.

The real finalization time of the public opening block is the origin
$t_0$ of \cref{ass:ready-service}. A direct regular successor may already
be Running before its export assignment opens. An unassigned boundary has
no selected target, so no preparation obligation applies to it; later
Espresso blocks alone cannot assign one. There is
no public delivery deadline: download completion, Ready, restart and
silence never generate a cross-region Force.


\paragraph{A complete regional ReadyBundle.}
\label{sec:ready-bundle}
The members of a selected region pass the entry material among themselves
in a ReadyBundle. It makes the transferred state explicit. It is not a
READY certificate and does not change the public opening of an export
duty.

For target session $(\beta,r)$, let $A=A_{\beta,r,0,0}$ and
$\Delta=\Delta_{i_{\beta,r}}$. A \tagg{READYBUNDLE} is an authenticated
regional message with the typed fields of the table below. Every required
object is present as bytes, once, in its normalized dependency set. A
URL, a state root, an unresolved hash or a promise to fetch later is not
a substitute.
\begin{table}[htbp]\centering\small
\begin{tabularx}{\linewidth}{@{}>{\raggedright\arraybackslash}p{.25\linewidth}Y@{}}\toprule
Field & Required value\\\midrule
Domain and target & Chain and protocol domain, encoding/native-state
format identifier, target $(\beta,r)$, selected region and sender identity.\\
Fixed handoff & The complete handoff of \cref{sec:witness}, including
its cause, predecessor/source context, selection and fixed target.\\
Exact anchor & $A$'s original producing pair, height, hash, full DATA,
original direct or covering commitment evidence and required ancestry.\\
Executable state $\mathit{st}(A)$ & The full, already constructed application-state
object from which execution after $A$ can continue. It includes all state,
code and continuation metadata required by the native execution interface;
if partial execution is supported, its exact continuation position is
included. It carries no regional locks, votes or entry timers.\\
State evidence & Native validation evidence connecting this exact
$\mathit{st}(A)$ to this exact $A$, with all proof inputs and dependencies, under the
sound application-validation interface. A sender's assertion or a bare
state hash is not enough.\\
Authority and bounds & Finalized $\tagg{AUTH}_{r,0}$, fixed configuration
and quota, $Q^0_{\beta,r}$, $q^{\rm open}_{\beta,r}$ and
$\lambda^{\rm open}_{\beta,r}$, and all authority evidence needed to
validate the inherited history.\\
Public evidence & An authorized replay checkpoint, with its trust and
authorization dependencies, and the complete finalized L1
control/authority and Espresso header/namespace histories needed to
verify the handoff, references and authorization context. It contains the
finalized header named by the predecessor's terminal reference (or the
Force activation header) and its evidence. If the numerical opening floor
is higher, no header at that floor is required for entry; the successor
obtains one before proposing its first DATA.\\
Dependency set & Full bytes for the transitive validation dependencies
of the fields above. A receiver need not fetch an object, replay
execution to rebuild $\mathit{st}(A)$, or obtain a new authority before processing
the bundle.\\\bottomrule
\end{tabularx}
\end{table}

The native interface determines how the state evidence verifies $\mathit{st}(A)$.
If the block format does not commit to a state root directly, the bundle
supplies the execution evidence and authenticated starting state that the
interface requires; this document does not invent such a commitment. A
serialized snapshot that still needs restoration or conversion is not yet
the directly executable state required here. Recipients validate the
packet themselves; they do not trust the sender to authorize a handoff or
vouch for state.

\paragraph{Sending and receiving a ReadyBundle.}
Each correct selected member builds its complete bundle before its first
Enter for the pair. Gathering histories, fetching across
regions, rebuilding state and forming the bundle all happen before that
event, and the model gives them no fixed time bound. At the event, the
member installs its verified state, executes initial Enter once, and
sends the complete bundle to every member of the fixed selected-region
snapshot. Since every correct member does this at its own first event,
the first correct entrant sends without needing to know that it is first.
Building and dispatching the bundle are not postponed past the recorded
entry time.

The receiver processes the bundle through its serialized Session Manager:
\begin{enumerate}[label=(\roman*),itemsep=2pt]
\item Authenticate the domain, sender and target, and inspect supplied
evidence under the evidence-first rule. Validate the historical handoff,
the exact anchor, all required authorities, references and native state
evidence. Incomplete or invalid material does not authorize entry. A
compatible existing cache can avoid duplicate validation, but the sender
never relies on an unknown receiver cache to omit required bytes.
\item Apply $\tagg{CanEnter}$ against the receiver's processed public
state. A closed or superseded target, a conflicting fixed handoff, a
delivered Close or a Terminated tombstone prevents initial Enter. The packet cannot
replace a fixed anchor or override a newer verified closure. If the
target is closed, run the Session Manager's phase-independent closure
relay using the packet's verified historical snapshot before refusing
entry. This also applies at a Terminated record and after a restart; no
program is recreated. If the same pair has already entered, keep
any useful evidence without resetting its state or entry time.
\item For an unentered valid target, take the state after the anchor:
the node's own baseline state if the native interface identifies it as
the state after the same anchor, and otherwise the received $\mathit{st}(A)$ after
validating it. Check every remaining condition of \cref{def:ready}, and
then execute the Enter step of \cref{alg:bootstrap} in this same
message-processing event. The durable entry flag, the attached state and
the entry time are recorded atomically before signatures are dispatched.
Relay the complete bundle as above. No Ready queue, state-fetch task or
post-Ready installation job lies between this event and Enter.
\end{enumerate}

To \emph{adopt} the state means that the validated received state object
becomes the new program's executable application state. The program works
on that object; it does not first rebuild a different database from it.
This is a contract on representation and on the message handler, not a
claim that database import, disk copying or runtime initialization take
zero time. No receiver can pass Ready while installation work remains.


\paragraph{Synchronous start.}
\begin{proposition}[Synchronous start]
\label{prop:sync-start}
Let $e$ be the first correct initial Enter time for session
$(\beta,r)$. Suppose the fixed handoff stays activatable, no closure or
other valid entry prohibition intervenes, and validation, adoption and
the atomic Enter transition together form the message-processing step
covered by \cref{ass:network}. Then
\[
 e_v(\beta,r,0,0)\le e+\Delta
 \qquad\text{for every correct selected member }v.
\]
\end{proposition}
\begin{proof}
The send rule places a complete correct-to-correct bundle send at $e$.
The first correct sender supplies the full executable object and all
validation dependencies. By the delivery-and-processing bound, each
correct recipient finishes that handler by $e+\Delta$. Its last guarded
step is Enter, unless it has entered earlier. No further
state-restoration premise is needed.
\end{proof}
A recipient whose own baseline state already matches the anchor uses
it, so in the normal case the bundle's state bytes serve only members
that lag.


This conclusion depends on the abstract-message reading above. The
network premise covers delivery and processing of complete validation
dependencies; it does not bound an external native import or storage
operation merely because the protocol calls it ``installation''. If such
work remains after bounded message validation, let $u_v\ge0$ be its
actual residual delay. Then only
\[
 e_v(\beta,r,0,0)\le e+\Delta+u_v
\]
holds, when entry eventually completes and stays authorized; the network
assumption gives no uniform bound on $u_v$, and calling that work one
event does not make $u_v=0$. In such an implementation, the $e+\Delta$
claim and the bootstrap bounds derived from it do not hold. A concrete
native adapter must show that it meets this contract; sound state proofs
alone do not establish a timing property. \Cref{ass:ready-service} bounds
total preparation from public assignment, but it does not prove $u_v=0$
or replace this handler obligation. The network bound also does not
bound construction of the first complete packet or historical catch-up.

A ReadyBundle only transports existing state and evidence within a
region. It adds no quorum of READY signatures, no Espresso transaction,
no publisher guarantee and no public activation certificate, and it does
not execute a Force. Public assignment and the current-entry check remain
those of \cref{alg:force,alg:external,sec:ready}.


\begin{protocol}[H]{alg:bootstrap}{Enter and bootstrap the selected session}
Replica $v\in V_{i_{\beta,r}}$ runs these steps for one fixed target
$(\beta,r)$. Let $\Delta=\Delta_{i_{\beta,r}}$.
\begin{enumerate}[label=\textbf{\arabic*)},itemsep=3pt]
\item \textbf{Enter.} When the program is Ready, atomically recheck
$\tagg{CanEnter}$ and the closed marks, then check and set the durable
initial-entry flag. Install $A_{\beta,r,0,0}$, $\tagg{AUTH}_{r,0}$,
$\TB_r$ and both opening reference bounds. Set the entry lock to
$Q^0_{\beta,r}$. Enter $(\beta,r,0,0)$ unfrozen, move the program from
Assigned to Running, record $e_v(\beta,r,0,0)$ once, and start the
progress checks of \cref{alg:recovery}. Send the complete ReadyBundle of
\cref{sec:ready-bundle} to the fixed membership. Send STATUS with this
lock and the fixed handoff to the initial leader, and process buffered
proofs for this session. Persist entry before dispatching signatures.
\item \textbf{Declare the fixed anchor.} The initial leader waits
$2\Delta$ from its own entry. If the view is still current and unfrozen
and the retirement contract permits the action, it sends NEWVIEW choosing
exactly $A_{\beta,r,0,0}$, justified by $Q^0_{\beta,r}$ and the fixed
handoff. The leader does not choose the anchor; the handoff has fixed it.
\item \textbf{Certify and produce.} Members run the ANCHOR-vote step of
\cref{alg:recovery}. On the ANCHOR QC, the leader starts
\cref{alg:steady} at DATA height $h(A_{\beta,r,0,0})+1$ and physical
position 1. Local failures use the same recovery rules. All these steps
occur while the program is Running, unless a Quit, Completion or Close
intervenes.
\end{enumerate}
\end{protocol}


For example, with anchor 108, the first DATA is 109. ANCHOR voting itself
starts no commit timer.

Let $e=\min_v e_v(\beta,r,0,0)$ over correct members. By
\cref{prop:sync-start}, all correct members enter by $e+\Delta$ if the
handoff stays activatable. With a correct initial leader, no Quit and no
Close:
\begin{center}\small
\begin{tabular}{@{}ll@{}}\toprule
Event & Bound\\\midrule
all correct members have entered & $e+\Delta$\\
NEWVIEW is sent & $e+3\Delta$\\
its ANCHOR QC reaches the leader & $e+5\Delta$\\
the first DATA reaches all correct members & $e+6\Delta$\\\bottomrule
\end{tabular}
\end{center}
These bounds start at the first actual Enter, not at public
assignment. The last one says the DATA has arrived, not that it
is committed or globally accepted. If native installation is not part of
bounded message processing, these bounds have not been established; the
extra delay is not silently absorbed into $\Delta$. \Cref{ass:ready-service}
alone does not prove this entry-skew bound.


\subsection{Synchronous stop}
\label{sec:retirement}
\label{sec:runtime-close}
Closing a branch version publicly and stopping a local program are
separate events. Public closure is immediate: once a Force takes effect
at its Espresso block, \cref{alg:force} stops the closed version from
appending anything new to $G$. Local stop takes time: a program that
learns of the closure stops new proposals, votes, BLAME and term or
recovery changes, but its running commit timers, FINAL messages,
receipts and evidence are processed as before. A new session need not
wait for an old one to stop. The rules below make sure that closure is
never silent: every correct member of an old session learns of it within
$\Delta$ of the first, running timers finish, and no evidence is lost.

Fix a session $(\beta,r)$ and write $\Delta=\Delta_{i_{\beta,r}}$,
using its immutable voting snapshot. Closure is a duty of the Session
Manager, whether or not a program exists or has entered. The Manager
stores the closed mark for $\beta$ and its complete closure evidence
durably, not just a hash or a retrieval promise, and keeps those bytes
available to the receive handler below. This branch-wide mark
immediately prohibits every old-$\beta$ initial entry and fresh
engine action. It is separate from the program states and from the
per-session close event defined below.

For each known fixed session in which $v$ is a member, the Manager keeps
its record in every state, and also when no program exists. The record
includes the fixed boundary, source region, immutable authority and
membership, whether Close has been performed, the first close time $c_v$
and pending relays. A program closed while Running keeps its engine
phase while it drains; only a genuine Quit or terminal Completion
changes that phase. Requests in a discarded suffix stay available to the
execution layer for reproposal.

\paragraph{Closure handling.}
The following handler runs at the Session Manager, in every program
state and when no program exists, not only at an entered engine.
\begin{enumerate}[label=(\roman*),itemsep=2pt]
\item \emph{Record public closure.} Verify complete evidence that the
canonical prelude of an authenticated finalized Espresso block closed
$\beta$. Retain its triggering header, applicable complete control and
namespace history from a validated checkpoint, and closure and
retained-prefix context. Check the explicit MSIG rule; a pending
MSIG, a proposed OFF flag, an unfinalized header or a bare timestamp is insufficient.
Atomically persist the closed mark for $\beta$, prohibit fresh actions for
that version and notify all its local control records, regardless of
phase. This consumes an established closure; it does not execute Force
or select an anchor a second time.
\item \emph{Disseminate for every known snapshot.} For each known fixed
session $(\beta,r)$ of this region, build its \defterm{complete close
bundle}. It carries two parts.
 \begin{itemize}
 \item The \emph{closure evidence} for $\beta$: the update $U\in\Lambda$
 that executed the Force, with $\Sigma_U$ and its L1 finality proof; the
 authenticated header of the activating block $J_q$; and the control
 history and activating headers from a validated checkpoint through
 $U$, which show that the branch version at the start of $J_q$ was
 $\beta$ and that $U$ is the first update of $\Lambda_q$ whose target is
 $\beta$.
 \item The \emph{snapshot context}: the session's immutable authority,
 membership and historical admission context.
 \end{itemize}
From these bytes alone another member can validate both the closure and
the exact recipient snapshot, without a prior program, live-entry
permission, an unprovided cache or its own replay reaching $J_q$. On first complete
processing, perform Close atomically: mark the record closed, record
$c_v$ as the current time, deliver Close to the local program if one
exists, and multicast the bundle to the fixed membership. This includes
Assigned, Running, Closing and Terminated records. A recipient validates
the bundle, persists the closed mark and performs the same idempotent operation;
no program or engine entry is created. Repeated bundles keep the first $c_v$.
\item \emph{Keep unresolved dissemination pending.} If the fixed snapshot or
its dependencies are not yet known, retain the closed mark with its
closure evidence and a pending relay obligation. The entry gate is already shut,
but $c_v$ is not assigned and no $4\Delta$ clock is invented.
Whenever any old-pair ReadyBundle, proposal, authority, history or other packet
provides the missing authenticated snapshot context, the Session Manager
completes and multicasts the close bundle in that receive
handler, and replies with it to the packet's sender. Do this before
ordinary BFT processing, rejecting a live-entry request, discarding a
duplicate, or returning a Terminated/tombstone result. Historical
context is validated without requiring that the old version remain
open. Complete proposals from correct voters carry or resolve this
context within the complete-message bound. This handler is
mandatory even if the old program was never created, its ReadyBundle
is refused, or a delayed callback arrives after exit or restart.
\item \emph{Gate new work.} The closed mark prohibits every old-version
DATA/DRAIN proposal, NEWVIEW, DATA/DRAIN/ANCHOR vote, new commit timer,
progress check producing BLAME, term entry and generation entry.
The checks cover initial Enter, delayed authority/recovery callbacks,
terminal-QC rotation and buffered-message processing. Neither a new
STATUS nor recreation of a local program bypasses the persistent gate.
Already dispatched messages are not withdrawn. New-session preparation
and entry use their own pair and need not wait for retirement.
\item \emph{Drain running work.} Keep running commit timers and
their deadlines; a closed mark alone does not cancel a
timer or release its evidence-processing owner. Continue complete
validation, comparison with retained signed material, and evidence-first
processing. A genuine current-generation quit proof still runs Quit,
broadcasts its proof, cancels its timers and preserves commitments;
it cannot reopen the gate. A surviving timer commits its block and
ancestors. For each DATA thereby committed, its FINAL from \cref{alg:steady} goes
both to the fixed regional membership and to the publication collector,
preserving \cref{def:receipt}. A partial FINAL does not invoke the
complete-result guarantee of \cref{ass:services}(i).
\item \emph{Keep evidence duties.} Continue collecting old votes, QCs,
FINAL messages and receipts. Multicast complete own-session receipts
as specified above, and offer every complete result to the Publish
Layer. A valid terminal receipt still runs genuine Completion,
possibly moving the engine to Completed and canceling timers; a
drain deadline that is already set still applies. Neither receipt arrival nor
any later update reopens the closed branch. Keep data, control-relay and publication
responsibilities until their durable transfer.
\end{enumerate}
All closure relays forward authenticated source evidence and
signatures that already exist, not new BFT votes or certificates. The Session Manager
continues transporting them after the program exits. A known snapshot
is never forgotten merely because its program exits. Recovery from a
crash restores closed marks, session close times, pending relays and owners
before dispatching callbacks or signatures. Messages, including complete
closure, quit and terminal-receipt evidence, are processed before a
simultaneous timer, entry or retirement event. No receipt, lock, vote,
committed block or provenance is erased by this handler.


\paragraph{Exit.}
For a program that is Running when it first learns closure, first
require its complete per-session Close event, at time $c_v$.
An earlier closed mark without a complete-close event is not a substitute
for $c_v$ and never authorizes early release of timer/evidence processing.
Exit may mark it Terminated only at a time $t\ge c_v+4\Delta$ when:
\begin{enumerate}[label=(\roman*),itemsep=2pt]
\item every commit timer that was running has executed or was canceled by
genuine Quit or Completion, and no queued action can pass a fresh-action
gate;
\item persistent storage owns the exit record, the closed mark, the per-session close records,
committed log, required signed evidence and validation dependencies,
receipts, anchor data and recoverable state, without releasing the last
required copy;
\item the Manager's duty table holds pending closure-context discovery,
complete-close relays and replies, complete-receipt multicasts,
complete-result offers, publication retries and data-service duties, and each locally complete result has been offered to the
Publish Layer; and
\item the dispatcher routes later messages for $(\beta,r)$ to the
Manager. It must run the closure handler even for an
unrecognized old pair, and may collect or relay existing signatures,
validate late receipts, serve data and continue publication, but cannot restart the
engine, execute its BFT timers or issue fresh BFT signatures for it.
\end{enumerate}
The ownership transfer may be local to the same correct replica; it
requires no additional agreement or certificate. The Manager must already hold the required evidence and receive route before
the program releases them; an unresolved retrieval task is not such
a handoff. Only after these
guards hold may the program and its timer machinery
be released. A late terminal receipt records archival evidence and does
not move a Terminated program back to another state.

A cancelled Assigned program has no regional timers to drain. It may exit without waiting $4\Delta$ only after
its durable exit record and all duties above are in the Manager's
duty table. In particular, learning Force before its first ReadyBundle never
permits silent cancellation: known snapshots trigger immediate
Close; unknown snapshots leave pending relays that the refused
ReadyBundle or a later proposal must resolve. A program that reached
Closing through Completion has already canceled its timers and may exit
after the same responsibility transfer; it likewise keeps all
control-relay duties. Termination never deletes a commitment,
removes DATA from $G$, or authorizes stake withdrawal.


\begin{proposition}[No silent closure in a vote window]
\label{prop:retirement-observable-close}
Let a correct member vote for and completely forward a valid proposal
of $(\beta,r)$ at time $s$. If a correct member's matching vote or entry
is suppressed by its closed mark during processing of
that proposal by $s+\Delta$, a correct member processes and multicasts
a complete close bundle for this session by $s+\Delta$. Thus, writing
$c=\min_{v\ {\rm correct}}c_v$ over all phase-independent complete-close
events (with $c=+\infty$ if none occurs), such suppression implies
$c\le s+\Delta$.
\end{proposition}

\begin{proposition}[Synchronous stop]
\label{prop:retirement-drain}
Under \cref{ass:network}, let $c=\min_{v\ {\rm correct}}c_v$ be the first
complete per-session Close time among this snapshot's correct
members, in any program state or with no program. Then every correct member has closed by $c+\Delta$.
Every surviving honest commit timer expires by $c+3\Delta$, and the
regional FINAL messages generated by those timer events reach every
correct member by $c+4\Delta$. Consequently $c_v+4\Delta$ is a
conservative local drain threshold, not a deadline for finishing all
retirement duties.
\end{proposition}

The bound starts at complete per-session closure processing and its
actual regional relay. It does not start at an earlier closed mark
whose session snapshot was unknown, a control signature, the triggering Espresso timestamp, or the
start of unbounded historical catch-up. It does not add a new source
visibility assumption. Three local $\Delta$ intervals cover the last
honest timer event but not its FINAL delivery: a member can vote just
before $c+\Delta$, commit just before $c+3\Delta$, and deliver FINAL
only just before $c+4\Delta$.

Four $\Delta$ intervals do not bound subsequent receipt relays,
Byzantine message release, a storage copy, $T_{\rm pub}$, $T_E$, or
successful EXPORT acceptance. Their outstanding duties pass to the
Manager's duty table; no bounded storage size or transfer latency is
claimed. Elapsed grace alone therefore never authorizes Exit.

\begin{proposition}[Conditional regional safety during retirement]
\label{prop:retirement-safety}
Assume the stated regional model, the regional validation
and DATA/DRAIN soundness premises, and the complete-message ReadyBundle
installation contract of \cref{sec:ready}. With the phase-independent closure handler of the Session Manager,
complete-receipt multicast and retirement rules above, correct regional committed logs
remain prefix-comparable and irrevocable. The external prefix argument
of \cref{sec:safety-argument} is preserved. This does not assert a
protecting QC for every late surviving timer, old-session progress,
or complete-receipt formation.
\end{proposition}


The proofs of the three propositions are in
\cref{app:retirement-proofs}.

\section{The Force handler}
\label{sec:force}
\Cref{alg:force} is the only cross-region Force transition. It executes
explicit L1 instructions. \Cref{alg:rotation} consumes the verified
closure and stops local programs separately (\cref{sec:retirement}). Silence, a health report
or a local timeout never runs this transition.

\subsection{What an instruction says}
\label{sec:force-stages}
An MSIG is an atomic branch update
$U=(\mathit{chain},c,\theta_a,\beta_{\rm target},S_{\rm new})$
(\cref{def:control-update}). It names the branch version to close and the
complete map to install on the next version. These two effects are
inseparable: if the target is current, both occur; if the target is
stale, neither occurs. The update names no retained height. Where the new
branch starts is decided when the instruction executes: if the committee
saw $G=108$ when it signed but $G=112$ when the update takes effect, 112
is kept. The committee submits the instruction under
\cref{ass:committee,ass:services}. L1 finalization makes it known; its
effect begins at $J_{q(\theta_a)}$. Until then, regular handoffs and
valid old-version output continue. A regular handoff changes $r$, not
$\beta$.

\paragraph{Why the target version is checked.}
Suppose instruction X says ``at time 120, close branch version 7''. If
in the meantime A only handed off to B by a regular handoff, the version
is still 7, so X executes at 120 and closes all of version 7, including
B's regular successor. But if another Force has already moved the
version to 8, X does not execute: its target is 7, and it never turns
into ``close whatever version is current''. So a late or duplicate
instruction cannot cut off a newer branch version: its target is stale
(\cref{def:control-effect}).

\paragraph{The export session.}
The \defterm{export session} is responsible for the next DATA after $G$.
Its source context, fixed quota and accepted count decide the target
index of a Force (\cref{sec:retry-slot}). It has no cross-region failure
timer. Its program may still be Assigned while this public responsibility
exists, and a later regular program may already be producing.

The newly installed map changes eligibility for the new branch only. An
update whose map leaves the source region ON can select that region again
when it is the only eligible one. If no region is eligible, the new
boundary keeps its anchor and target index without opening a regional
engine. Since the map is immutable within that version, only another
successful atomic update can supersede this unassigned boundary and try
a different complete map.

\begin{protocol}{alg:force}{Apply an atomic branch update}
Before the payload of authenticated finalized block $J_q$, let
$\theta=\theta_q$ and record the branch version $\beta_0$ at the start
of the block.
\begin{enumerate}[label=\textbf{\arabic*)},itemsep=3pt]
\item \textbf{Authenticate the due batch.}
From the complete applicable L1 control log, collect the valid due
updates $\Lambda_q$ not yet processed, in sequence order
(\cref{def:control-effect}). Repeated source delivery or equal timestamps
never processes an update twice.

\item \textbf{Choose the executing update.}
If no entry of $\Lambda_q$ targets $\beta_0$, mark every due entry
processed, leave all protocol state unchanged and return to ordered
payload processing. Otherwise let
$U^*=(\mathit{chain},c^*,\theta_a^*,\beta_0,S^*)$ be the first matching
entry. It is the sole executing update in this prelude; every other due
entry is a stale whole-update no-op. A version created in this prelude is
never a fresh matching target. No output, Ready, elapsed-time or health
condition is checked.

\item \textbf{Close, install the map and keep $G$.}
If $\mathsf{Tail}=\tagg{NEXT}(\nu)$, record
$\operatorname{sid}(\nu)$, $i(\nu)$ and $a(G,\nu)$; if
$\mathsf{Tail}=\tagg{PENDING}(\omega)$, read the stored source context
and target index from $\omega$.  Let $i_{\rm src}=i(\nu)$ in the first
case and the stored source region in the second. Keep all of $G$ as it
stands at the start of $J_q$, with its original block-bound records. Keep
$U^*$ and the header of $J_q$ as closure evidence. Close $\beta_0$, set
$\beta'=\beta_0+1$ exactly once and install $S_{\beta'}:=S^*$. Fix the
target index $r'$ by \cref{sec:retry-slot}, let
$A=\operatorname{last}(G)$ and compute
$i'=\tagg{NextLive}(i_{\rm src},S_{\beta'};r')$. If the required target
authority is missing, pause this operation without interpreting the
candidate set. With complete evidence, set
\[
 \mathsf{Tail}:=
 \begin{cases}
 \tagg{PENDING}(\tagg{FORCE}(c^*,q),\beta',r',i_{\rm src},A,q),
   & i'=\bot,\\
 \tagg{NEXT}(\beta',r',i',A,h(A)+\TB_{r'},q,
              \tagg{FORCE}(c^*,q)),& i'\ne\bot.
 \end{cases}
\]
Archive the closed version's unconfirmed suffix and regular successors.
The new engine, when entered, starts with fresh core state and $g=t=0$.
No new-version DATA is accepted in $J_q$ itself.

\item \textbf{Persist once.}
Persist the processed control cursor, the new version and map, the closed
mark, $G$ and the new tail as one replay transition. Mark all other due
entries processed as stale no-ops. Then return to ordered payload
processing.
\end{enumerate}
\end{protocol}

The cutoff and retained anchor are fixed even if target evidence arrives
later. Resume the same canonical operation; never rerun a persisted
update. A successful update can have a fixed closure while target
selection still waits for its authority evidence. In contrast, later
Espresso blocks alone never revisit a PENDING boundary: changing its map
requires another successful atomic update.

Every prelude executes at most one Force. The activation thresholds of
two Forces are at least $T_{\rm update}$ apart
(\cref{def:control-valid}); no DATA, Ready or other output is required
between them.
Each update must still satisfy its own authorization,
advance-announcement and target-version conditions.

\subsection{Atomic maps, stale updates and the cutoff}
If an update's target is stale, the entire update is a no-op. It does not
install $S_{\rm new}$, close the current version, select a region or
choose an anchor. If the map change is still desired, the committee must
submit a fresh update targeting the current version. This atomic rule is
why a map change cannot retroactively revise a regular handoff within an
open version.

All control effects come before the payload of the triggering block.
DATA accepted in earlier blocks stays in $G$. An old-version result in
the triggering block cannot cancel a Force by arriving with late
progress. It may be kept as regional evidence but cannot append its
suffix. A newer version starts only from the fixed retained anchor.

For example, take threshold $d=120$ and Espresso blocks with timestamps
119, 122 and 125. The block at 119 accepts valid DATA of the current
version. The block at 122 is the first to reach 120, so the Force runs
first and old-version DATA in that block can no longer be accepted;
nor can any later old-version DATA. New-version DATA can be accepted only
from the block after 122.

A proof of closure contains the actual MSIG and the complete applicable
source context. A bare timestamp, an unavailable RPC response, an
unfinalized control or a monitoring alarm is not such a proof. After
learning of the closure, regional Session Managers apply
\cref{sec:retirement}; the cutoff does not wait for their local cleanup.

\section{The external observer}
\label{sec:observers}
To \emph{replay} means to start from a verified state and apply
\cref{alg:global-rule} to the finalized Espresso blocks in order, which
yields the same control results and the same $G$. Any correct verifier
that processes, say, one due update targeting version 7 and installing a
complete map with B ON, followed by an old-version DATA 109, reaches the
same map, the same closure and the same rejection of 109. There is no
hidden private decision, and the
result does not depend on who sent a handoff message to whom first.

An observer keeps the production branch, the replay state
$\mathcal X=(\mathrm{Out},S_\beta,\kappa)$ of \cref{alg:global-rule}, the
object set $\mathcal O_q$ and its source-verification cursors. Its public
output is still $\mathrm{Out}=(G,\mathsf{Tail})$. It also persists its
replay subposition. It stores no per-block
relay clocks and no committee-protected prefix. The same deterministic
logic can run inside a certifier.

\begin{protocol}{alg:external}{Verify Espresso output and advance the prefix}
\begin{enumerate}[label=\textbf{\arabic*)},itemsep=3pt]
\item \textbf{Authenticate the next source input.} Obtain the next
finalized Espresso block, its complete namespace, and the complete
finalized L1 control context that applies, linked to the trusted prior
prefix. Resume a durable subposition if this block is partly executed.
Obtain every authority $\tagg{AUTH}_{r',0}$ needed by an atomic update or
a terminal regular selection in this block.
Regional objects that are not yet published are not fetched: the rule
reads only $\mathcal O_q$ (\cref{sec:record-completeness}), and optional
proposal, recovery-authority and native sidecars do not delay it. Do not
require the completion of a later successor before appending the
current DATA. Missing genuine source inputs cause a wait, never an
empty-history interpretation.

\item \textbf{Run the prelude once.} Execute \cref{alg:force}, which
contains Step~1 of \cref{alg:global-rule}, before any DATA is appended,
and persist its stage and historical snapshots.
Deliver a verified closure and its source proof to the Session Manager.
Correct regional nodes store and propagate the version-wide gate for
known snapshots, including Assigned and Terminated programs; unresolved
contexts keep durable relay obligations under \cref{sec:retirement}.

\item \textbf{Apply the global finality rule.} Add the block's regional
objects to $\mathcal O_q$ (\cref{sec:logical-carrier}) and apply
Step~2 of \cref{alg:global-rule} while a matching DATA exists.
Durably expose each appended DATA, its count and the subposition at
once. Invalid or incomplete objects do not block another DATA or
reverse an append. No duplicate DATA is counted again.

\item \textbf{Select after a terminal commit.} The terminal case of
Step~2 constructs the REGULAR boundary, computes its reference floor and
applies NextLive under the unchanged $S_\beta$. With complete evidence
and an eligible region it installs $\tagg{NEXT}(\nu)$. A genuine evidence wait persists
here without withholding the terminal DATA that is already accepted;
resume exactly once at this subposition. With complete evidence but no
candidate, it installs $\tagg{PENDING}(\omega)$ and source replay
continues.

\item \textbf{Finish or resume.} When no step applies, the block is
done; advance to the next block unless a step is waiting for evidence.
Keep the block's snapshots taken in the prelude unchanged. Persist every
completed step; later source or boundary processing cannot roll it back.
Expose provisional production separately. A Force archives only the
unconfirmed suffix and keeps $G$.
\end{enumerate}
\end{protocol}

Missing evidence and having no candidate are different. If B's authority
has not been downloaded, the node does not know whether B is eligible and
must fetch the evidence. Only when all authorities and eligibility are
known and no region qualifies does the boundary stay pending; it changes
only if a later successful atomic update creates a new branch and map. A
download failure is never recorded as ``no candidate'',
and late evidence never moves a control's effective position to the
present. A missing regional object is different again: it is not a wait
at all. Step~2 simply finds no candidate, the block finishes, and the
object is used in whichever later block publishes it.

The outer fields of an EXPORT, such as its expected base, are not inputs
to the rule. Step~2 derives the real base from $G$ and checks the signed
parent link. A certifier attestation binds the complete actual previous
and next replay checkpoints: $\mathcal X=(\mathrm{Out},S_\beta,\kappa)$,
the replay subposition, pending controls and the source position. It can attest an
appended prefix before a later successor selection finishes; it cannot
replace a pending boundary with an invented empty candidate set. The certifier
checks that issued controls were executed correctly. It does not judge,
on the committee's behalf, whether a region was missing blocks for too
long.

\subsection{Execution-layer requests and stored data}
\label{sec:requests}
The execution layer supplies admissible payloads and manages client
queues, ordering policy, persistence and resubmission. The consensus
layer may accept empty DATA under its normal quota rules. It does not
promise that an individual request survives a discarded unconfirmed
suffix or is eventually included. For example, if transaction $x$ is in
an unconfirmed 109 that a Force discards, the user or builder may
resubmit $x$. If $x$ is already in $G$, no Force can remove it.

This does not weaken the storage of block data and proofs. Keep the DATA,
receipts and exact state evidence needed for recovery, handoff and
checkpoint verification. Before Espresso finalization, keep complete
results available to the Publish Layer and direct readers. Afterward,
authenticated history and the node-level store provide the required
objects. Global acceptance never authorizes deleting the only needed
copy. The accepted $G$ binds every publication path.

\paragraph{Faithful publication.}
The proposer or batcher publishes the next complete interval after
the accepted DA publication prefix, even if a Force has changed the current region. The
interval must belong to $G$, keep order and multiplicity, and match the DA
payload and native execution evidence. An unavailable accepted payload
causes waiting, not an invented replacement. The native adapter is
specified separately: a Force cannot move back from a partial-EVM anchor
just because the next signer cannot continue it. It must obtain a legal
continuation or wait.

\section{Examples and guarantees}
\label{sec:examples}
\subsection{Regular rotation and a late DRAIN}
\label{sec:regular-example}
Let $h(A_{\beta,r,0,0})=100$, $\TB_r=10$ and $K=4$. The DATA terms produce
101--104, 105--108 and 109--110. A terminal QC below 110 moves to the next
term. A receipt for DATA 110 opens the next region, which starts DATA 111
after Ready and bootstrap. Step by step:
\begin{enumerate}
\item A commits 110 under its regional BFT rules, if necessary through a
DRAIN during recovery.
\item A's members collect the complete regional receipt for 110 and send
it to B.
\item B verifies that it is selected and obtains the data of 110, the
execution state after 110 and the target authority.
\item B enters $(7,21,0,0)$, runs the ANCHOR bootstrap and produces 111.
\item A's result is published to Espresso in parallel; $G$ may reach 110
later.
\end{enumerate}
Step 2 needs the receipt. A QC for 110 alone is not enough for B to take
over.

\begin{figure}[htbp]\centering
% Regular handoff: a terminal receipt transfers the turn from A to B.
% Representative arrows only; time runs downward and is not to scale.
\begin{tikzpicture}[
  x=1cm,y=-1cm,font=\fontsize{9}{10.5}\selectfont,
  party/.style={draw,rounded corners=2pt,fill=white,align=center,
    inner sep=3pt,minimum width=2.2cm},
  lane/.style={draw=black!30,densely dashed},
  msg/.style={-{Stealth[length=1.7mm]},semithick,draw=rrBlue},
  direct/.style={msg,densely dashed},
  tag/.style={fill=white,inner sep=1.5pt,font=\fontsize{8}{9}\selectfont},
  note/.style={anchor=north west,align=flush left,text width=5.35cm,
    inner sep=0pt,font=\fontsize{8.5}{10}\selectfont},
  phase/.style={anchor=west,fill=white,inner sep=2pt,
    font=\fontsize{9}{10.5}\selectfont\bfseries},
  dot/.style={circle,fill=black,inner sep=1.1pt},
  timer/.style={line width=2.6pt,draw=rrTeal!75},
  wait/.style={line width=2.6pt,draw=rrBlue!70!black,
    dash pattern=on 2.5pt off 1.5pt},
  entry/.style={draw=rrBlue,line width=1.2pt,dotted}
]
  \def\xA{1.2}\def\xR{3.7}\def\xL{6.2}\def\xB{8.7}\def\xn{9.65}
  \path[use as bounding box] (-.1,-.65) rectangle (15.05,15.75);
  \node[party] at (\xA,0) {A: leader /\\receipt holder};
  \node[party] at (\xR,0) {A: other\\replicas};
  \node[party] at (\xL,0) {B: initial leader\\session $(\beta,r+1)$};
  \node[party] at (\xB,0) {B: other\\replicas};
  \foreach \x in {\xA,\xR,\xL,\xB}
    \draw[lane] (\x,.55)--(\x,15.35);

  \node[phase] at (0,.85) {1.\ A commits terminal DATA in session $(\beta,r)$};
  \draw[msg] (\xA,1.40)--node[tag,above,sloped] {$T_{\beta,r}=110$}(\xR,1.70);
  \draw[msg] (\xR,2.15)--node[tag,above,sloped] {DATA votes}(\xA,2.45);
  \foreach \x/\s/\e in {\xA/1.40/3.80,\xR/1.70/4.10} {
    \draw[timer] (\x+.20,\s)--(\x+.20,\e);
    \node[dot] at (\x,\s) {};
    \draw[thick] (\x-.12,\e)--(\x+.45,\e);
  }
  \node[note] at (\xn,1.35) {Replicas validate and vote for the last DATA
    block. Each starts its own $2\Delta_{\rm A}$ commit timer.};
  \node[tag,anchor=east] at (\xA-.15,2.90) {QC$(110)$};
  \node[note] at (\xn,2.80) {The terminal QC does not hand the turn to B.
    The quota permits no DATA 111 in A's session.};
  \draw[msg] (\xR,4.25)--node[tag,above,sloped] {FINAL$(110)$}(\xA,4.55);
  \node[dot] at (\xA,4.55) {};
  \node[note] at (\xn,3.95) {Timers expire; replicas commit 110 and send
    FINAL. Matching signatures of weight $>F_{\rm A}$ form
    $\FC_{\beta,r}(T_{\beta,r})$. Receipt holders complete the old engine.};

  \node[phase] at (0,5.55) {2.\ Direct transfer of the regular boundary};
  \node[note] at (\xn,5.50) {Verify the session and quota. Select B from
    the fixed map $S_\beta$ under $\tagg{AUTH}_{r+1,0}$. The handoff keeps
    $\beta$ and changes $r$ to $r+1$.};
  \draw[direct] (\xA,6.35)--node[tag,pos=.55,above,sloped]
    {REGULAR witness + bundle}(\xL,7.35);
  \draw[direct] (\xA,6.35)--(\xB,7.55);
  \node[note] at (\xn,6.90) {Direct transfer may precede Espresso
    publication. It waits for neither global acceptance nor a future
    reference header.};

  \node[phase] at (0,8.05) {3.\ B opens $(\beta,r+1)$ and starts DATA};
  \draw[entry] (\xL-.6,8.45)--(\xL,8.45)--(\xB,8.65)--(\xB+.6,8.65);
  \foreach \x/\y in {\xL/8.45,\xB/8.65} \node[dot] at (\x,\y) {};
  \draw[msg] (\xL,8.75)--node[tag,above,sloped] {Ready relay}(\xB,8.95);
  \draw[msg] (\xB,9.25)--node[tag,above,sloped] {STATUS}(\xL,9.50);
  \draw[wait] (\xL+.20,8.45)--(\xL+.20,10.85);
  \node[note] at (\xn,8.40) {Ready requires the boundary, anchor 110,
    history and authority. Enter view $(\beta,r+1,0,0)$; relay the bundle
    and send STATUS.};
  \node[note] at (\xn,9.95) {The initial leader waits $2\Delta_{\rm B}$
    from its own entry. Progress checks already run.};
  \draw[msg] (\xL,10.90)--node[tag,above,sloped] {NEWVIEW$(110)$}(\xB,11.20);
  \draw[msg] (\xB,11.65)--node[tag,above,sloped] {ANCHOR votes}(\xL,11.95);
  \node[dot] at (\xL,11.95) {};
  \node[note] at (\xn,11.05) {NEWVIEW must choose the inherited anchor.
    Weight $>F_{\rm B}$ gives an ANCHOR QC; these votes start no timer.};
  \draw[msg] (\xL,12.55)--node[tag,above,sloped] {DATA 111}(\xB,12.85);
  \foreach \x/\s/\e in {\xL/12.55/14.95,\xB/12.85/15.25} {
    \draw[timer] (\x+.20,\s)--(\x+.20,\e);
    \node[dot] at (\x,\s) {};
    \draw[thick] (\x-.12,\e)--(\x+.45,\e);
  }
  \draw[msg] (\xB,13.35)--node[tag,above,sloped] {DATA votes}(\xL,13.65);
  \node[note] at (\xn,12.55) {The first fresh block has DATA height 111
    and physical position 1. It waits for a finalized Espresso reference
    at or above $q^{\rm open}_{\beta,r+1}$, then its votes start new
    commit timers.};
  \node[note] at (\xn,14.60) {If uninterrupted, the new timers commit 111;
    replicas send FINAL for B's session.};
\end{tikzpicture}

\caption{Regular cross-region rotation keeps the DATA, FINAL and ANCHOR
sequence. The immutable branch map $S_\beta$ fixes the successor, while
the terminal Espresso reference fixes only the next reference floor.
Regional production and Espresso publication proceed in parallel. Neither
the terminal receipt nor the new region waits for global confirmation.}
\label{fig:regular-cross-region}
\end{figure}
\FloatBarrier

Now suppose 110 has a QC but a Quit cancels its timer. Recovery opens at
height 110 and proposes a DRAIN, which supplies a new timer without using
DATA quota. When the DRAIN commits the terminal DATA, replicas send FINAL
for 110, not for the DRAIN. The regular successor still inherits 110. The
old L1 and Espresso references of 110 say nothing about when it was
committed or exported.

Suppose $G$ is only at 100 while A has committed 110 and B has committed
111. The Publish Layer can export 101--111 with the required A and B
receipts and the regular witness. Once those DATA are accepted from
finalized Espresso history, the whole interval is global. A record that
declares only 111 can still supply the predecessors first if it also
carries their DATA and receipts. It cannot skip ancestors that are truly
missing.

\subsection{A multisignature Force and the retained anchor}
\label{sec:force-example}
This is the running example of \cref{sec:overview}. The region order is
A, B, C. A's session $(7,20)$ produces DATA 101--110. Its terminal receipt
lets B start $(7,21)$ before all of A's DATA is globally accepted. B has
produced through 115, while $G$ ends at 108. The export session is still
A's: the next required DATA is 109.

The committee submits one atomic update with activation threshold 120
that satisfies advance announcement. It targets branch version 7 and
supplies a complete new map that sets A OFF and keeps B, C and D ON.
Before the first finalized Espresso block that
reaches 120, valid old output may still extend $G$.

If $G$ still ends at 108 at that block, replay checks the target and
atomically closes version 7 and installs the new map with A OFF before
the payload. It keeps 108. Old
A:109--110 and B:111--115 can no longer be appended. Their receipts stay
as historical evidence of regional commitment.

\begin{center}\small
\begin{verbatim}
Production before the Force:
  A: 101 ... 108 109 110
  B:                     111 ... 115
Kept in G at the Force:
     101 ... 108
New B after the Force:
                new 109, new 110, ...
\end{verbatim}
\end{center}

If B is eligible, the forced target is $(8,21)$, because A has already
contributed accepted DATA. The new engine enters $(8,21,0,0)$ from 108 and
may propose DATA 109 after Ready and the ANCHOR bootstrap. It does not
wait for Espresso results to be registered on L1. New-version DATA is
accepted only in a later Espresso block.

If an older update had already moved version 7 to 8 before 120, this
update would be stale in full: it would neither close version 8 nor
install its map. If instead A had completed its export and B had become
responsible through regular rotation, the version would still be 7: the
instruction would execute, keep the newer $G$, and select after the
actual export region B.

With quota ten, old B $(7,21)$ starts after 110 and ends at 120. New B
$(8,21)$ starts after 108 and ends at 118. They use the same immutable
authority index but different engines and anchors.

\begin{figure}[p]\centering
% Atomic MSIG branch update and the existing regional bootstrap.
% beta is the numeric branch version, and (beta,r) identifies a session.
% Vertical spacing illustrates event order, not elapsed time.
\begin{tikzpicture}[
  x=1cm,y=-1cm,font=\fontsize{8.3}{9.8}\selectfont,
  party/.style={draw,rounded corners=2pt,fill=white,align=center,
    inner sep=3pt,text width=1.22cm,minimum height=.66cm},
  lane/.style={draw=black!28,densely dashed},
  msg/.style={-{Stealth[length=1.6mm]},semithick,draw=rrBlue},
  read/.style={msg,densely dashed},
  force/.style={msg,draw=rrRed},
  tag/.style={fill=white,inner sep=1.5pt,font=\fontsize{7.7}{9}\selectfont},
  note/.style={anchor=north west,align=flush left,text width=5.55cm,
    inner sep=0pt,font=\fontsize{8.3}{9.8}\selectfont},
  phase/.style={anchor=west,fill=white,inner sep=2pt,
    font=\fontsize{9}{10.5}\selectfont\bfseries},
  dot/.style={circle,fill=black,inner sep=1.1pt},
  wait/.style={line width=2.5pt,draw=rrTeal!75,
    dash pattern=on 2.5pt off 1.5pt},
  rulebox/.style={draw=rrRed!55,rounded corners=2pt,fill=rrRed!3,
    align=center,text width=7.9cm,inner sep=5pt}
]
  \def\xC{.75}\def\xI{2.3}\def\xE{3.85}
  \def\xN{5.4}\def\xL{6.95}\def\xB{8.5}\def\xn{9.45}
  \path[use as bounding box] (-.05,-.55) rectangle (15.05,16.05);
  \node[party] at (\xC,0) {Multisig};
  \node[party] at (\xI,0) {L1};
  \node[party] at (\xE,0) {Espresso};
  \node[party] at (\xN,0) {Node /\\certifier};
  \node[party] at (\xL,0) {B: initial\\leader};
  \node[party] at (\xB,0) {B: other\\replicas};
  \foreach \x in {\xC,\xI,\xE,\xN,\xL,\xB}
    \draw[lane] (\x,.52)--(\x,15.65);

  \node[phase] at (0,.88) {1.\ Announce an atomic branch update on L1};
  \draw[force] (\xC,1.55)--(\xI,1.85);
  \node[tag,anchor=south] at (1.52,1.48) {submit $U$ at $s$};
  \node[note] at (\xn,1.30) {$U$ names target version $\beta$, future
    time $\theta_a$, and a complete map $S_{\rm new}$. If $\beta$ is
    stale, the entire update is a no-op.};
  \node[dot] at (\xI,2.65) {};
  \draw[read,draw=rrRed] (\xI,2.65)--node[tag,above,sloped]
    {finalized update + complete history}(\xN,3.30);
  \node[dot] at (\xN,3.30) {};
  \node[note] at (\xn,2.85) {By $s+T_{\rm L1}$, all correct participants
    can verify the finalized update. $T_{\rm L1}$ includes dissemination
    and verification. This precedes the Espresso crossing.};
  \node[tag,align=center] at (5.35,4.15)
    {Before activation: version $\beta$ and $S_\beta$ stay fixed;\\
     normal regional handoffs may continue.};
  \node[phase] at (0,5.25) {2.\ Finalize the Espresso boundary; apply one update};
  \draw[force] (\xE,5.95)--(\xN,6.25);
  \node[tag,anchor=south] at (4.62,5.90) {$J_{q^*}$ finalized};
  \node[note] at (\xn,5.75) {$q^*=\min\{q:\theta_q\ge\theta_a\}$ for
    this update. Interpret the complete history through $J_{q^*-1}$
    first.};
  \node[rulebox] at (4.65,7.30) {At the start of $J_{q^*}$:\\[2pt]
    first due update targeting current $\beta$ executes;\\
    close $\beta$ and install its full $S_{\rm new}$ atomically;\\
    every other due update is a whole-update no-op;\\
    only then process EXPORT records.};
  \node[note] at (\xn,7.05) {Health detection belongs to multisig policy.
    A report or missing DATA cannot execute Force. No Ready or output
    condition gates a valid matching update.};
  \node[dot,fill=rrRed] at (\xN,8.55) {};
  \node[tag,anchor=east,text=rrRed] at (\xN-.12,8.55)
    {anchor $:=G$ before payload};
  \node[note] at (\xn,8.70) {Close version $\beta$ for appends; keep $G$.
    Advance to $\beta+1$. Use $S_{\beta+1}$ and NextLive to select B,
    or keep a PENDING tail.};
  \draw[read] (\xN,9.60)--node[tag,pos=.55,above,sloped]
    {complete boundary}(\xL,10.00);
  \draw[read] (\xN,9.60)--(\xB,10.15);
  \node[note] at (\xn,10.15) {If the old export session has accepted
    count $a$: keep $r$ when $a=0$, otherwise use $r+1$.
    New-version DATA can append only in a later Espresso block.};

  \node[phase] at (0,11.10) {3.\ B uses the existing ANCHOR bootstrap};
  \node[dot] at (\xL,11.65) {};
  \node[dot] at (\xB,11.85) {};
  \draw[msg] (\xB,12.20)--node[tag,above,sloped] {STATUS}(\xL,12.45);
  \draw[wait] (\xL+.18,11.65)--(\xL+.18,13.15);
  \node[tag,anchor=east] at (\xL-.12,12.85) {wait $2\Delta_{\rm B}$};
  \node[note] at (\xn,11.75) {B verifies the boundary, anchor and state.
    At Ready, enter $(\beta+1,r',0,0)$. The leader waits
    $2\Delta_{\rm B}$, then sends NEWVIEW.};
  \draw[msg] (\xL,13.20)--node[tag,above,sloped] {NEWVIEW}(\xB,13.45);
  \draw[msg] (\xB,14.00)--node[tag,above,sloped] {ANCHOR votes}(\xL,14.25);
  \node[dot] at (\xL,14.25) {};
  \node[note] at (\xn,13.65) {An ANCHOR QC with weight $>F_{\rm B}$
    permits DATA after $G$. ANCHOR votes start no commit timer.
    DATA and DRAIN retain their rules.};
  \draw[msg] (\xL,15.00)--node[tag,above,sloped]
    {DATA at $h(G)+1$}(\xB,15.30);
  \node[note] at (\xn,15.20) {The forced handoff preserves $G$;
    it does not wait for $G$ to be registered on L1.};
\end{tikzpicture}

\caption{Explicit Force: L1 supplies the instruction, Espresso time fixes
its cutoff, and the next region inherits exactly the $G$ at the start of
the activating block. There is no automatic trigger and no second
Espresso-to-L1 wait.}
\label{fig:forced-cross-region}
\end{figure}
\FloatBarrier

\subsection{Concrete view updates}
\label{sec:force-view-examples}
The following examples assume complete authority and an eligible target.
Every Force has an explicit valid MSIG cause.

\paragraph{A has accepted output.}
A is in $(7,20,2,3)$, producing 101--110, while $G=108$. Its export
session is 20 with a positive accepted count. Selecting B gives
\[
 (7,20,2,3)\xrightarrow{\mathrm{Force}}(8,21,0,0).
\]
B continues after 108. Only the new engine's $g$ and $t$ start at zero;
the retained DATA keeps its original provenance.

\paragraph{Regular rotation has already reached B.}
A first hands off
\[
 (7,20,0,2)\xrightarrow{\mathrm{regular}}(7,21,0,0).
\]
B reaches $(7,21,0,1)$ and contributes DATA through 113 to $G$. The next
DATA, 114, belongs to B's session 21. A Force selecting C gives
\[
 (7,21,0,1)\xrightarrow{\mathrm{Force}}(8,22,0,0).
\]
C continues with DATA 114. A scheduled MSIG targeting version 7 still
matches, because regular rotation did not change $\beta$. Even if the
committee saw only $G=108$ when it signed, it cannot make the chain go
back to 108.

\paragraph{Speculative production is ahead of the export session.}
A's $(7,20)$ has regional DATA 101--110; B's $(7,21)$ has 111--120; C runs
$(7,22,1,2)$. If $G=108$, the next required DATA is still A's 109. A Force
selecting B opens $(8,21,0,0)$, not $(8,23,0,0)$. It closes the entire old
version, including C's speculative program. Using a lower turn index than
the latest speculative one removes no globally accepted DATA.

\paragraph{A stale instruction changes no view field.}
One MSIG targets version 7 at threshold 120. An older update has already
moved the version to 8 at an earlier position. Suppose the new
engine is now in $(8,21,1,2)$. Processing the stale instruction leaves
\[
 (8,21,1,2)\xrightarrow{\mathrm{stale\ MSIG}}(8,21,1,2).
\]
Its proposed map does not take effect. To change the map or force again,
the committee needs another valid atomic update that targets version 8.

\subsection{Zero output, an offline region and explicit exclusion}
\label{sec:retry-example}
A's $(7,20)$ opens after DATA 100 but gets no DATA accepted. The committee
observes the failure and submits one update that sets A OFF and carries
the complete new map while targeting version 7. At its activation, if
$G$ is still 100, the target index stays 20. NextLive can choose B, which
enters $(8,20,0,0)$ with a fresh engine. Reusing the number 20 does not
mean selecting A again.

If B also adds no DATA, another explicit MSIG can choose C in
$(9,20,0,0)$. Neither B's Ready nor a first DATA is required before that
instruction. If A's DATA 101 had been accepted before the first Force,
the new index would be 21 with anchor 101.

The monitoring report itself has no effect.
\begin{center}\small
\begin{tabularx}{\linewidth}{@{}>{\raggedright\arraybackslash}p{.45\linewidth}Y@{}}\toprule
Situation & What the protocol does\\\midrule
A misses blocks for a long time, no MSIG & $\beta$ does not change\\
matching update whose new map leaves only A ON & close the old version; NextLive may select A again\\
matching update whose new map has A OFF and B ON & close the old version and hand off to B\\
matching update whose new map sets every region OFF & keep $G$ and the target index in a PENDING tail\\
later update targeting that pending version with B ON & create another version and hand off to B\\\bottomrule
\end{tabularx}
\end{center}

\subsection{Repeated Force and an unassigned boundary}
\label{sec:repeated-force}
A valid MSIG can close a newly assigned version while its target is
still Assigned, or after it starts but before any DATA. Forces need no DATA
in between; their thresholds need only be $T_{\rm update}$ apart. Each
new instruction still needs a matching
target and a valid advance announcement.

If closing session $(7,20)$, which has accepted output, finds no eligible
region, replay stores target index 21 and anchor $G$ in version 8. It
opens no regional engine. A later valid MSIG can move the version to 9
while keeping target index 21 and the same $G$. If that update's complete
map sets B ON, B enters $(9,21,0,0)$, not $(9,22,0,0)$: a session that
never started does not count as a completed turn. If no successful update
ever supplies an eligible map, the wait can last forever.

\subsection{Strict cutoff, equal timestamps and concurrent preparation}
\label{sec:cutoff-example}
An MSIG targets version 7 at threshold 300. B may start regular session
$(7,21)$ from A's terminal receipt before that threshold. Meanwhile its
baseline state follows $G$, so it already holds the state at any Force
anchor. The baseline grants no entry.

At the first finalized block with timestamp at least 300, the prelude
checks the still-current version. A match closes every session of
version 7 and keeps the exact preceding $G$. Old DATA in that block's
payload cannot revoke the control. A program created for that
Force starts from the state at that exact anchor, never from state
computed for another one.

If Espresso timestamps jump across the threshold, the first crossing
block is still the activation block. A later block with the same
timestamp does not repeat an already processed instruction. A distinct
valid instruction can execute at a later position if its own target
matches. If the jump also crosses the threshold of the next update, both
updates take effect in that prelude, in sequence order. A single prelude
never chains several Forces.

\subsection{Conditional safety and lifecycle invariants}
\label{sec:safety-argument}
The two propositions below cover public replay and the local programs.
They are safety statements only; progress is discussed in
\cref{sec:liveness}. They use the model and verification interfaces of
\cref{sec:model}. In plain words: everyone who reads the same complete
history accepts the same chain; DATA in $G$ is never removed by a Force;
and no program enters twice or mixes state with another session.

\begin{proposition}[Common interpretation and preservation of global DATA]
\label{prop:global-prefix}
Let two correct observers start from the same authenticated replay
checkpoint and apply \cref{alg:global-rule} under the stated history,
announcement and verification premises. Their global DATA sequences are
prefix-comparable: one is a prefix of the other. Replaying to the same canonical subposition with all applicable source
evidence and L1 controls verified yields the same output
$\mathrm{Out}=(G,\mathsf{Tail})$ of \cref{sec:exports}. At a completed Espresso prefix
this gives the corresponding completed-prefix result as a special case. It does not require identical
future-control queues or source-reading checkpoints. Every accepted DATA
block keeps its content, position, producing identity and bound
session record forever. Once a version is closed, no
later EXPORT adds new DATA from that version to $G$.
\end{proposition}
\begin{proof}
Induct over the common finalized Espresso sequence and, within a block,
over the steps of \cref{alg:global-rule}.
The initial checkpoints agree. Complete applicable L1 controls,
immutable authorities, Espresso headers and namespace bytes authenticate
the same next input. Advance announcement supplies controls before their
crossing; missing actual evidence pauses at the same operation. Atomic
update order, the first matching target, target indices and fixed-map
NextLive selection are deterministic. They give the same prelude and
immutable snapshots. Future controls not yet due change neither result.

Both observers read the same object set $\mathcal O_q$ from the same
authenticated namespace content. A private cache changes retrieval work,
not $\mathcal O_q$. The matching conditions of Step~2 and its hash tie-break are
functions of $\mathcal O_q$ and the current output, so both observers
append the same block. By \cref{prop:final-kernel-soundness},
\tagg{CheckFINAL} establishes intrinsic input validity from a genuine
correct signer's regional commit. The independent boundary, ancestry,
source-order and cutoff checks establish current append permission;
regional safety identifies the same next DATA. Optional admission or
native sidecars cannot veto that kernel. Executable state, native output
and DA correspondence remain subject to their separate sound interfaces. Each DATA transaction appends the
same next block and accounting, and persists the same subposition.
A terminal transaction constructs the same boundary and applies the same
fixed-map selector. It either selects the same eligible successor, returns
the same pending tail, or pauses at the same missing evidence. Missing evidence
cannot undo the already committed transaction or alter its timestamp.
These rules give agreement by induction over all committed subpositions.
Two observers that have reached different subpositions hold ordered
states on this append-only DATA sequence, including when one is paused
inside a source record. Extra future-control queues and private validation
caches are not part of the compared replay result.

Inspect preservation at each step. Control ingestion removes no DATA. A
successful update keeps all of $G$ at the start of the block and changes
only the branch version, its map and the tail for future appends. A
pending boundary also preserves $G$. Each accepted DATA transaction only extends it. Retained DATA keeps its
original producing pair and hash. Closed versions are never reopened,
and every later DATA candidate checks the current version. Thus
neither an old regular successor nor a late receipt can append a closed
suffix. Local program phases and retirement do not mutate $G$, so they
add no exception to this induction.
\end{proof}

\begin{proposition}[Fixed entry and isolation of local programs]
\label{prop:program-isolation}
Under the regional-validity and source-verification premises, each
opened pair $(\beta,r)$ has one fixed handoff. At a correct replica,
initial Enter occurs at most once for that pair, uses exactly that
handoff's anchor and state, and never uses an obsolete task callback
as authority. Actions for one pair do not alter another pair's locks,
votes, timers or entry time. Baseline execution creates none of
these BFT actions.
\end{proposition}
\begin{proof}
Proceed from the initial checkpoint through successive handoffs.
For a regular handoff, the predecessor's fixed boundary and quota fix
its terminal height; regional validity fixes the terminal DATA and its
signed Espresso reference. The immutable $S_\beta$ and target authority
fix the region, while the terminal reference fixes the next opening
floor. For Force, common replay fixes the old and new versions, source
context, exact retained $G$ and newly installed map; NextLive under that
map fixes the target. Repeated Forces without assignment
change $\beta$ but keep the stored predecessor index. They cannot
reuse an already opened pair. Handoff verification rejects conflicting
fixed fields instead of replacing an existing record.

Baseline state follows $G$ over time but is not a program. Only Assign
creates a program, and only after verification installs the actual
handoff.
Material callbacks keep the requested pair and anchor identity;
the gate checks the exact fixed state and separate current-activation predicate.
The serialized Enter call tests and sets the durable entry flag
before dispatching signatures. Repetition, restart and late callbacks
therefore cannot repeat Enter or reset its entry time. Each core
message is authenticated and routed by the pair it signs; local tasks
use the same pair and anchor context. An Assigned program does not run
the core.
Closing and Terminated records prevent a late event from reopening an
old program. These checks establish the stated isolation by induction
over the local event sequence.
\end{proof}

This is local authorization relative to a node's processed public
prefix. It does not claim that all nodes learn a finalized Force at the
same real time. A lagging node can briefly run an old program, but
\cref{prop:global-prefix} still excludes its closed-version output.
\Cref{prop:retirement-safety} covers regional behavior during
retirement. Both propositions rely on Part~I and on the DATA/DRAIN and
application-validation interfaces rather than re-proving them.
\subsection{Explicit control and conditional preparation}
\label{sec:liveness}
The guarantees in this subsection do not require every region to receive
service, and they do not make the consensus layer responsible for
request persistence.

\begin{proposition}[No unannounced cross-region Force]
\label{prop:explicit-force}
Starting from an authenticated checkpoint, every later increment of
$\beta$ has a unique effective MSIG cause in the verified L1 control
history. Missing DATA, a local health report or a Ready event does not
increment $\beta$; there is no separate status-update transition.
\end{proposition}
\begin{proof}
Step~3 of \cref{alg:force} is the only transition that increments
$\beta$. It is reached only for the first update due at the block whose
target matches the branch version recorded at its start. Control identities
and persisted subpositions prevent a second application. Other
transitions select under an already fixed map, append DATA or update
local programs; none changes $\beta$ or installs a map. Induction over committed
transitions proves the claim.
\end{proof}

\begin{proposition}[Preparation under a continuing assignment]
\label{prop:ready-bound}
Let a public assignment open in block $J_o$, with finalization time
$t_0$. If the selected region and its state service satisfy
\cref{ass:ready-service}, the fixed handoff stays applicable, and no
Force supersedes it during preparation, then correct target members can
complete Enter by $t_0+T_{\rm ready}$. Entry uses the retained
anchor of that assignment and happens at most once per member and pair.
Entry does not require any DATA production.
\end{proposition}
\begin{proof}
A member that already entered satisfies the claim. For a member that has
not entered, the service assumption supplies every material condition of
\cref{def:ready} by the stated time, and continuing applicability
supplies the current-entry check. \Cref{alg:rotation,alg:bootstrap} then
perform the guarded entry in the same complete event.
\Cref{prop:program-isolation} gives anchor binding and at-most-once
entry. The gate does not look at future DATA output. A valid Force before
completion would remove the continuing-assignment premise; it would not
violate any waiting rule.
\end{proof}

\subsection{Conditions for progress}
\label{sec:progress}
Being Ready only means that a region can start. Global progress further
needs regional consensus, a FINAL receipt, publication and global
acceptance. The following conditions are separate from prefix safety.
\begin{itemize}
\item Regional service keeps Part~I's fair leader schedules, eventual
exclusion of exposed equivocators, actual entry into correct turns, and
uninterrupted commit opportunities. These are conditions inside a region,
not fairness among regions.
\item Espresso keeps finalizing retrievable history with growing
timestamps. Complete results reach the Publish Layer, and the declared
publication and finalization services hold.
\item The required authority, exact anchor state and acceptance evidence
are supplied. A selected live region satisfies \cref{ass:ready-service}.
The next required DATA is published in a form and position that pass the
current-branch acceptance checks.
\item If the current region cannot serve, recovery additionally depends
on the multisignature service actually submitting a suitable atomic
update whose complete map keeps or enables an eligible working region,
and eventually allowing
enough uninterrupted execution for regional commitment and delivery.
\end{itemize}
The last item is a liveness condition on the control service, not a
condition for accepting an MSIG. The protocol specifies no detection or
response bound, the only interval it enforces between Forces is the
Espresso-time gap $T_{\rm update}$, and $T_{\rm L1}$ starts only at actual
submission. So an honest service can still fail to
give a useful execution opportunity unless its policy provides one.

A wholly offline region is not an operational region for these service
claims. The regional safety and lifetime key-compromise bounds still
govern every accepted certificate. Request persistence, resubmission and
eventual inclusion belong to the execution layer and the client
interface. If Espresso fails, global confirmation stops; switching the
finality service is outside this protocol.

\subsection{What is needed for global progress}
\label{sec:global-progress}
Under the Part~I conditions above, a session that keeps running can form
a complete result. The bounded preparation assumption alone does not
provide that uninterrupted regional opportunity.

If a complete result is available at time $s$, the declared services give
submission by $s+T_{\rm pub}$ and Espresso visibility by
$s+T_{\rm pub}+T_E$. This is transport of the result. To advance $G$, the
result must supply the next required DATA under its fixed boundary, while
its producing branch version is still open at the canonical acceptance
position, and its references and predecessors must pass the checks of
\cref{sec:exports}. Because the rule ignores outer hints, reads the
cumulative $\mathcal O_q$ and appends one DATA at a time, wrong hints or
a later pending successor do not erase an otherwise valid next DATA. They never permit a gap or an append from a
closed version.

With these regional, transport, acceptance and control-service
conditions, replay of the suitable result advances $G$. The service
bounds give no wall-clock processing deadline for every arbitrary
observer. No unconditional global-growth theorem follows from repeated
Ready events, published raw records or committee honesty alone.

\paragraph{Offline recovery requires control.}
When the current region cannot continue, progress needs the committee to
submit a suitable atomic update whose complete map disables a failed
candidate when necessary and keeps or enables a live eligible successor.
If no such update is
submitted, nothing else forces a switch. If all regions are OFF, the
boundary can stay unassigned forever. If $T_{\rm update}$ is shorter than
a useful interval, MSIGs can keep closing every one, and no regional
output need ever become global. These executions
are allowed unless the governance service provides more liveness
conditions. The committee can thus stall the chain, but under the
safety assumptions it cannot use a Force to roll $G$ back.

\paragraph{Failure and safety are separate.}
An offline region does not meet the operational premises of the
preparation and regional-service claims. A violated lifetime Byzantine
weight bound is a different failure, of the certificate model. Neither
case lets the multisignature committee roll back $G$. No result here
promises fair service among regions, automatic ON, one DATA between
Forces, or eventual inclusion of a client request.

\section{Implementing the Publish Layer}
\label{sec:publish-implementation}
The preceding protocols use a delivery-service assumption. The following
leader-first design explains how existing participants can implement it.
It adds no consensus round and no exclusive publisher role.
The implementation must meet $T_{\rm pub}$; the sketch alone is not a
proof of that bound.

\subsection{Leader first, with permissionless backup}
The proposing regional leader is the preferred sender. Once it has the
complete finalized result, it builds and submits an EXPORT immediately.
Other replicas continue collecting FINAL messages, keeping payloads and
making complete results available to the service. Commitment and
Completion do not wait for a publication acknowledgement.

Another replica or an external data verifier may submit the same result
without the leader's permission. An implementation may give the leader a
short first opportunity, but the backup delay, data retrieval and retries
must fit within $T_{\rm pub}$ from complete-result availability. A private
announcement that the leader has sent something is not enough to suppress
backup indefinitely. The service checks submission and Espresso evidence.

The fallback sender uses the original DATA and receipt. It cannot replace
the DATA, rewrite its session pair or create FINAL on behalf of the region.
The service guarantees delivery through an available path; no particular
sender is guaranteed to win the submission race.

\subsection{Fee sponsorship}
A backup sender need not hold tokens or pay from its own balance. It may
use a funded relay, a prepaid submission account, or another sponsored
route that can submit to Espresso. The route carries the verified result;
it does not acquire authority to change its content.

Someone still supplies the fee resources. To justify the publication
bound, a complete result covered by the service assumption must have access to a funded,
working route even when the original leader is unavailable. A sender can
try another route if one refuses or fails. The concrete sponsorship and
admission mechanism is an implementation choice; this protocol assumes
no Espresso application smart contract for it.

\subsection{Duplicates, retries and acceptance}
Different senders may carry the same complete result in different EXPORTs.
The signed producing pair and block hash identify each DATA; changing a
receipt encoding, signer set or outer hint does not create a new output.
Accepted ancestors are reused as evidence and receive no second count.
A compatible next DATA can be extracted even when the sender's expected
base is stale or its result is listed only among supporting ancestry.

The Session Manager tracks the original EXPORT record and its per-DATA
acceptance. A complete next result already published with a bad hint is
interpreted at that source position under the same rule; a corrected
EXPORT is unnecessary.
If a genuine gap or an absent required object prevents acceptance, the
objects already published stay in $\mathcal O_q$; once the missing
objects appear in a later block, \cref{alg:global-rule} accepts the
result there without republication. No bound on when the missing
objects arrive is inferred. A closed version never regains append
permission.

The two service bounds concern complete-result delivery and visibility.
The conditional progress discussion of \cref{sec:global-progress}
uses those actual recorded bytes. It does not assume a bounded retry
with a corrected expected-base hint. Terminal DATA commits before successor authority
resolution. Preparation for a selected live region has the separate
scope of \cref{ass:ready-service}; arbitrary archive catch-up has no
uniform bound. \Cref{sec:transport-scope} states what the transport
must put into Espresso for the rule to use it.

\subsection{Retention and incentives}
Before finalized Espresso inclusion, the holders of a complete result
keep its payload and evidence available to backup senders. Afterward,
Espresso's retained, authenticated history supplies returning readers and
initial hydration. Prepared regional nodes additionally keep the raw
source bytes under \cref{sec:source-retention}; finalization does not
replace that retained copy by an eventual archive pointer. Direct transfer
can still accelerate regular handoff and state recovery. Retirement and
checkpoint rules preserve required evidence and outgoing duties.

Submission rewards and reimbursement may fund this service. They must not
count duplicate DATA as new progress. A report of delayed delivery
alone does not prove that the original leader, verifier or sponsor
equivocated. Health policy and any objectively verifiable penalties
are separate from the Force acceptance rule.

\Needspace{14\baselineskip}
\begingroup\small
\begin{thebibliography}{9}
\bibitem{part1}
Sync HotStuff RR, Part I: Core Protocol. Design document, 13 October 2026.
Companion source: \href{run:sync_hotstuff_rr_part1_core.tex}{Part I core protocol}.
\bibitem{sync}
I. Abraham, D. Malkhi, K. Nayak, L. Ren, and M. Yin.
\emph{Sync HotStuff: Simple and Practical Synchronous State Machine Replication}.
IEEE Symposium on Security and Privacy, 2020, pp. 106--118.
\href{https://conferences.computer.org/sp/pdfs/sp/2020/349700a106.pdf}{Published paper}.
\bibitem{certification}
MegaETH Specification, PR~8, \emph{Block Certification}.
Commit \texttt{aedc2826663b}.
\href{https://github.com/megaeth-labs/specification/tree/aedc2826663b134ee25d26bea17e70c6d23d321c/specs/block-certification}{Pinned specification}.
Its Espresso and certifier verification interfaces are a reference;
the rotation and explicit-control rules in this document are new.
\end{thebibliography}
\endgroup
\appendix
\crefalias{section}{appendix}
\section{Proof of \texorpdfstring{\Cref{prop:final-kernel-soundness}}{the FINAL kernel proposition}}
\label{app:final-kernel-proof}
We restate \cref{prop:final-kernel-soundness} in the precise form the
proof uses.

\paragraph{Precise statement.}
\textit{Assume the fixed-snapshot fault bound, the accepted-lifetime
key-compromise bound, the cryptographic assumptions, sound validation
interfaces, and the regional full-validation and FINAL-generation rules.
If $C$ is the authentic incoming context of $(\beta,r)$ and
$\tagg{RawFINAL}_C(D,\Sigma)$ holds, then a correct replica committed
this exact DATA $D$ at a timer event, directly or as an ancestor. Before
that event it fully checked its own-session proposal and required input
ancestry, using established historical kernel conclusions only for
inherited anchors. So a genuine valid historical witness exists for $D$.
Consequently the exact DATA has valid historical regional admission and
native input validity under its authenticated anchor, including the
regional $\rho$ and $\xi$ validity conditions. Exact same-session DATA
ancestors on that hash-linked path inherit these commitment and
input-validity conclusions.}

\begin{proof}
By the fixed-snapshot fault bound, a set of distinct signers with weight
greater than $F_i$ includes a correct signer. The lifetime compromise
condition and signature unforgeability make its signature an actual
correct FINAL action on the stated DATA hash and context. This step uses
the raw threshold, not the definition of a valid regional receipt.

A correct replica sends FINAL only after Commit's timer event commits the
named DATA, either as the timed DATA itself or in the ancestry of the
timed DATA/DRAIN block. Vote started that timer only after all
current-session dependencies were present and the proposal and required
ancestry passed full validation. In particular the replica checked its
session's historical authority and admission, parent justifications,
quota and positions, the L1 and Espresso reference rules, and native
input validity. The sound interfaces apply to these actual regional
checks. An inherited anchor from another producer may rely on that
producer's historical kernel conclusion, justified by the induction
below; this does not replace the replica's exact native state or its
own-session validation. These facts give a genuine valid historical
witness for the committed DATA. It need not be the witness that a later
publisher attaches.

A commit triggered by a receipt, including Completion, emits no new
FINAL. So it cannot supply the first correct signature without the
earlier fully validated timer event identified above. Dependencies on
earlier accepted anchors are justified by induction over execution
events from the trusted initial checkpoint. The FINAL signatures used by
an inherited anchor exist before the new Vote event, so their own timer
origins are earlier cases of the induction. The same argument covers an
unaccepted REGULAR anchor without assuming its global acceptance. No step
assumes the validity of the kernel acceptance being proved. Hash binding
identifies the supplied DATA and exact parent-linked bodies with those in
the correct signer's validated history. The timer's ancestor-commit rule
gives the same-session coverage claim.
\end{proof}

\section{A reference control policy}
\label{app:policy}
This appendix sketches one policy that an honest committee can run on
top of the interface of \cref{sec:runtime-controls}. The protocol does
not depend on it, and certifiers do not check it. Other conditions or
thresholds keep every guarantee of the main text.

\paragraph{What the controller observes.}
The controller acts on a verified replay position $q$. There it reads
$\mathrm{Out}_q$, the current immutable map $S_\beta$ and the
control log, including its pending updates. It also knows the complete
regional results it has received, and operational reports. When the
$\mathsf{Tail}_q=\tagg{NEXT}(\nu)$, write
$(\beta,r,i,a)=(\beta(\nu),r(\nu),i(\nu),a(G_q,\nu))$
for the export context. Two times are read on the global clock of
\cref{ass:network}: $t_0$, the finalization time of the block
that opened the export session (as in \cref{ass:ready-service}), and
$t_G$, the finalization time of the block that appended
$\operatorname{last}(G)$. Let $t$ be the controller's current time.

\paragraph{Conditions.}
The policy uses the following predicates, with policy parameters
$\tau_{\rm boot}$, $\tau_{\rm stall}$ and $k_{\rm off}\ge1$.
\begin{itemize}
\item $\tagg{Pub}$: the controller knows a complete regional result for
DATA height $h(\operatorname{last}(G))+1$ of the export session that is
not yet in $\mathcal O_q$.
\item $\tagg{Boot}$: $\mathsf{Tail}_q=\tagg{NEXT}(\nu)$, $a=0$, $\tagg{Pub}$ fails,
and $t-t_0>\tau_{\rm boot}$.
\item $\tagg{Stall}$: $\mathsf{Tail}_q=\tagg{NEXT}(\nu)$, $a>0$, $\tagg{Pub}$ fails,
and $t-t_G>\tau_{\rm stall}$.
\item $\tagg{Fault}(i)$: the controller holds evidence that region $i$
cannot be trusted or cannot serve. The evidence is \emph{checkable}, for
example two different regionally final DATA blocks of one session at one
height, or \emph{attested}, for example an operator or security report.
\item $\tagg{Maint}(i,t_m)$: region $i$ is to stop serving at time $t_m$.
\item $\tagg{Restore}(i)$: region $i$ is OFF, and the controller holds
evidence that it can serve again.
\item $\tagg{Unassigned}$: $\mathsf{Tail}_q=\tagg{PENDING}(\omega)$.
\end{itemize}
$\tagg{PendU}$ holds when the controller has an issued update that has
not yet reached its activation block. $\mathit{strikes}(i)$ counts the
Forces issued under \tagg{BOOT} or \tagg{STALL} whose export region was
$i$, since $i$ last added DATA to $G$.

\paragraph{Updates issued.}
The reference policy is \defterm{single-flight}: while $\tagg{PendU}$
holds, it issues no second update. It accumulates desired operational
changes off chain. After the outstanding entry is canonically processed,
the controller rereads the resulting $\beta$ and $S_\beta$, reevaluates
all conditions, and, if action is still needed, submits one update that
targets that $\beta$ and carries one complete new map. This avoids a
queue in which the first update changes $\beta$ and makes the rest stale.

When no update is in flight, start from $S^\circ=S_\beta$, apply every
condition that currently holds, and batch the result into one complete
$S_{\rm new}$. \Cref{tab:control-conditions} gives the intended map
change. Every issued update also forces a branch change; there is no
status-only row.

\begin{table}[htbp]\centering\small
\caption{Reference policy: conditions and the update issued.}
\label{tab:control-conditions}
\begin{tabularx}{\linewidth}{@{}>{\raggedright\arraybackslash}p{.15\linewidth}>{\raggedright\arraybackslash}p{.2\linewidth}YY@{}}\toprule
Code & Condition & Complete $S_{\rm new}$ & Replay effect\\\midrule
\tagg{FAULT} & $\tagg{Fault}(i)$ & $S^\circ$ with $i$ OFF & Close
$\beta$, retain $G$, and select under the new map\\
\tagg{MAINT} & $\tagg{Maint}(i,t_m)$ & $S^\circ$ with $i$ OFF, with
activation chosen for $t_m$ & Close $\beta$ at activation and select
under the new map\\
\tagg{BOOT} & $\tagg{Boot}$, not $\tagg{PendU}$ & $S^\circ$, with export region $i$
OFF if $\mathit{strikes}(i)+1\ge k_{\rm off}$ & Close
$\beta$; target index $r$; selection after $i$\\
\tagg{STALL} & $\tagg{Stall}$, not $\tagg{PendU}$ & As for \tagg{BOOT} &
Close $\beta$; target index $r+1$; selection after $i$\\
\tagg{PUB} & $\tagg{Pub}$ & No update; ask a backup sender to publish
(\cref{sec:publish-implementation}) & None\\
\tagg{RESTORE} & $\tagg{Restore}(i)$ & $S^\circ$ with $i$ ON & Close
$\beta$ and make $i$ eligible on the new branch\\
\tagg{UNASSIGNED} & $\tagg{Unassigned}$ & $S^\circ$ with a chosen region
$j$ ON & Close the pending version; keep its target index and try $j$ on
the new branch\\
none & None of the above, including an L1 or Espresso outage & No update
& None\\\bottomrule
\end{tabularx}
\end{table}

Under this interface even \tagg{RESTORE} and maintenance of an idle
region close the current branch, because every map change is atomic with
a Force. The policy should therefore batch such changes rather than use
the interface for frequent health telemetry. A \tagg{FAULT} with
checkable evidence forces at once. $\tagg{Boot}$ and
$\tagg{Stall}$ require that $\tagg{Pub}$ fails on purpose: when the next
DATA exists but has not been published, a Force would discard committed
output and gain nothing.

\paragraph{Issuance rules.}
\begin{enumerate}[label=(R\arabic*)]
\item \textbf{Complete map and target.} Every update names the branch
version the controller observed and carries a complete map derived from
that same replay state. A pending update is never rewritten. If it turns
out stale, its entire effect is lost and the controller reevaluates the
conditions against the new current version.
\item \textbf{Single flight.} Normally issue no update while one is
unactivated. After it executes or becomes stale, resubmit any still
desired intent with a fresh sequence, target and complete map.
\item \textbf{Timing.} At actual submission time $s$, choose the
earliest admissible threshold: above $s+T_{\rm L1}+\epsilon_E$
(\cref{sec:assumptions}) and at least $T_{\rm update}$ after the last
issued threshold. The update takes effect in whatever state replay has
reached by then. If submission slips, sign a fresh update.
\item \textbf{Short horizon.} Do not schedule an update more than a
horizon $T_{\rm horizon}$ ahead; keep planned changes off-chain until they are near.
A far-future in-flight update delays any later intervention under the
single-flight policy.
\item \textbf{Keep a candidate.} Do not issue a map that leaves no
eligible region, unless the code is \tagg{FAULT} with checkable
evidence.
\item \textbf{No anchor.} The controller never tries to pick the
retained DATA; replay keeps the actual $G$ at the start of
$J_{q(\theta_a)}$.
\end{enumerate}

\paragraph{Choosing the thresholds.}
\tagg{BOOT} and \tagg{STALL} should never fire against a region that
serves as the model allows. Let $T_{\rm first}$ bound the time from
$t_0+T_{\rm ready}$ until a complete result for the session's first DATA
is available, and $T_{\rm next}$ the time from one complete result of
the session until the next. These bounds come from the regional service
conditions of Part~I; this document does not derive them. Choose
\[
 \tau_{\rm boot}\ge T_{\rm ready}+T_{\rm first}+T_{\rm pub}+T_E,
 \qquad
 \tau_{\rm stall}\ge T_{\rm next}+T_{\rm pub}+T_E.
\]

\begin{proposition}[The policy does not interrupt a serving region]
\label{prop:policy-noninterference}
Let the thresholds satisfy the bounds above, and let the controller keep
up with Espresso history as in \cref{ass:services}(iii). Consider an
assigned export session that no Force closes, whose region satisfies
\cref{ass:ready-service}, whose first complete result is available by
$t_0+T_{\rm ready}+T_{\rm first}$, whose later results each follow the
previous one within $T_{\rm next}$, and whose results pass the
acceptance checks of \cref{sec:exports}. Then $\tagg{Boot}$ and
$\tagg{Stall}$ never hold for this session.
\end{proposition}
\begin{proof}
The first result is available by $t_0+T_{\rm ready}+T_{\rm first}$. By
\cref{ass:services}(i) and (iii), it is in a finalized Espresso block,
and verified by the controller, by
$t_0+T_{\rm ready}+T_{\rm first}+T_{\rm pub}+T_E\le
t_0+\tau_{\rm boot}$. Replay accepts it in that block, because
it extends $G$, passes the acceptance checks and its branch version is
open. So once $t-t_0>\tau_{\rm boot}$, the controller already
sees $a>0$, and $\tagg{Boot}$ fails. For $\tagg{Stall}$, let the last
accepted DATA be appended at $t_G$. Its result was available at some
time $s\le t_G$. The next result is available by $s+T_{\rm next}$ and is
accepted and verified by $t_G+T_{\rm next}+T_{\rm pub}+T_E\le
t_G+\tau_{\rm stall}$. So $t_G$ moves forward before $\tagg{Stall}$ can
hold. When the session's last DATA is accepted, the tail no longer names
this session, and neither predicate applies to it.
\end{proof}

The proposition covers only \tagg{BOOT} and \tagg{STALL}. A \tagg{FAULT}
or \tagg{MAINT} decision, or a threshold set too low, can still
interrupt a serving region. While $G$ is fixed, each \tagg{BOOT} or
\tagg{STALL} Force moves the selection to the next eligible region, and
a serving region is never interrupted by them; this is the step that
recovery needs in \cref{sec:progress}. A full recovery theorem also needs
the authority of each new target index to keep that region eligible; we
do not prove it here.

\section{Proofs for the synchronous stop}
\label{app:retirement-proofs}
This appendix proves the propositions of \cref{sec:retirement}.

\begin{proof}[Proof of \cref{prop:retirement-observable-close}]
The forwarded proposal and its historical authority/admission dependencies
reach the member's Session Manager within $\Delta$, even if it has
no program or only a tombstone. These dependencies identify the fixed
regional snapshot. If the closed mark was already recorded, the receive handler
combines them with its retained closure evidence and sends the complete
bundle before refusing ordinary processing. If closure is learned after
that packet but by the time it suppresses a vote, the snapshot is already
known and closure's own handler performs the multicast immediately.
Either way some $c_v\le s+\Delta$. No new initial entry, executable-state
restoration or remote membership discovery is required for this control
relay. Restart and duplicate paths use the same persistent handler.
\end{proof}

\begin{proof}[Proof of \cref{prop:retirement-drain}]
The first correct relay includes closure and the fixed snapshot's full
validation context. Its delivery and processing are covered by the
regional bound even when no receiving program exists. Thus
$c\le c_v\le c+\Delta$ and every old-version fresh-action gate is shut
by $c+\Delta$. A member may have set its closed mark still earlier; that can only
stop its votes earlier. A correct DATA/DRAIN vote at time
$s=\tau_v(P)$ starts its only timer, with expiry $s+2\Delta$.
The close gate gives $s\le c_v\le c+\Delta$; no callback starts a
replacement timer. Thus its expiry is at most $c+3\Delta$.
Quit and Completion can cancel a timer but cannot postpone it.
The FINAL messages actually generated at surviving expiries are
delivered regionally within another $\Delta$. Finally $c_v\ge c$
implies $c_v+4\Delta\ge c+4\Delta$. Message-first processing covers equality.
\end{proof}

\begin{proof}[Proof of \cref{prop:retirement-safety}]
We verify the prerequisites of the timer/lock argument explicitly;
Part~I's safety theorem is not used to infer an unproved universal-QC property.
Fix a correct direct timer for DATA or DRAIN block $B$ in generation
$g$, voted at $s$, that actually expires at $s+2\Delta$ without
cancellation. All comparisons below concern this same session.
Let $e_0$ be its first correct initial-entry time. By the synchronous
start (\cref{prop:sync-start}), every nonclosed correct member enters by $e_0+\Delta$.
An own-session terminal receipt requires a correct FINAL after a timer,
whose earliest possible expiry is $e_0+2\Delta$. Hence such a receipt
cannot evade Completion by arriving only at an unentered cache of a
correct member that will run the chosen timer.

\emph{Quit and Completion cannot suppress its early vote window.}
If any correct member first quits $g$ at $a\le s+\Delta$, its complete
quit proof reaches the committer by $a+\Delta\le s+2\Delta$ and
cancels the timer. Thus any first correct Quit satisfies
$a>s+\Delta$. If any correct member Completes at $b\le s+\Delta$,
its complete terminal receipt is multicast by $b$ and reaches the
committer by $b+\Delta\le s+2\Delta$. The committer then executes
Completion: it commits the terminal DATA and ancestors and cancels
this timer, including a DRAIN timer. Again this contradicts the chosen
surviving direct timer. No correct Completion can therefore occur at
or before $s+\Delta$. This is an operational cancellation argument;
it does not assume that receipt-triggered commitments are already
known to be prefix-comparable.

\emph{The proposal is processable within its regional window.}
For $g=0$, the self-contained ReadyBundle relay and installation contract
in \cref{sec:ready} supplies entry by $s+\Delta$; this assertion is
conditional on that precise contract, not a bound on state restoration
before any member is Ready. For $g>0$, let $a^-$ be the first correct
Quit of $g-1$, and let $e$ be the first correct entry into $g$.
The recovery rule gives $e\ge a^-+\Delta$ and $s\ge e$.
The complete quit relay reaches every correct member by $a^-+\Delta$;
its own mandatory wait has therefore ended by
$a^-+2\Delta\le e+\Delta\le s+\Delta$. The forwarded complete
proposal supplies the finalized authority and remaining admission
history. Hence a receiver with no Quit, Completion or close preventing
entry can process it in $g$ by $s+\Delta$. Applying the same
complete-proof argument inductively to earlier generations supplies
any predecessor entry needed to process the quit proof. This uses
the recovery waits and complete validation bundles, not a
new simultaneous authority-visibility time.

\emph{Early and late windows.}
Let $c$ be the first complete per-session Close event among
correct Session Managers, in any program state. A prior closed mark or
unentered/Terminated record is not excluded merely by definition of $c$.
Instead, \cref{prop:retirement-observable-close} shows that if it
suppressed a matching vote or entry by $s+\Delta$, then $c\le s+\Delta$.
Consequently, if $s+\Delta<c$, no closure suppresses that vote or entry.
The preceding cancellation arguments also exclude Quit and Completion
through that window. Earlier-term processing and prior indirect
commitment do not suppress that vote either. An incompatible retained
history would expose a genuine quit proof in this window and contradict
timer survival. Every correct member therefore has cast, or now casts,
the matching DATA/DRAIN vote by $s+\Delta$. Their total weight is at least
$W-F>F$, so the votes form $B$'s protecting QC by $s+2\Delta$ at
all members that can still take a later recovery entry. A member that
Completed after the vote window cannot take such an entry; it need
not be assigned an unused recovery lock.

Any correct entry into $g+1$ must obey
\[
 e_v(\beta,r,g+1,0)\ge a_v+\Delta
 \ge a+\Delta>s+2\Delta.
\]
In the early case, every such entrant therefore knows the protecting
QC before its lock snapshot. The highest-QC rule gives the
recovery-lock invariant of Part~I. Closure preserves collecting those already
sent votes and QCs, and may only remove a later entry. In the late
case $s+\Delta\ge c$, the timer expires at or after $c+\Delta$.
The complete snapshot-bearing relay has closed every correct member's
entry gate by then, including previously unentered and archived members; the strict displayed
inequality prohibits every later-generation entry and thus its
correct signatures. A later-generation QC needs a correct signer,
so no recovery can exploit a missing protecting QC in this case.

\emph{Direct timers and receipt-based commits.}
In one generation, complete signed histories either extend one
another or expose an admissible equivocation proof.
A correct voter forwards its complete history. A correct signer of
an incompatible history either processes $B$ first and cannot vote,
or forwards its history before processing $B$, no later than
$s+\Delta$. That history reaches the committer by expiry and exposes
the proof. Genuine Quit, or Completion before expiry, would cancel
the timer. Thus no conflicting same-generation correct direct timer
can survive; the same argument excludes a conflicting QC, including
one assembled later from old votes.

Across generations, a surviving lower-generation direct timer cannot
be in the late case if a correct higher-generation vote exists.
Its early-case QC is known before every later live lock snapshot.
The ranked-QC descent of Part~I now has its actual prerequisites:
correct ANCHOR voters meet those locks; DATA/DRAIN QCs inherit a
lower-ranked parent justification. It excludes a later QC that does
not extend the committed block, and the parent of a later direct
proposal consequently extends it. These are precisely the
protecting-QC, recovery-lock and unique-extension steps of Part~I,
with the explicit vote-window checks above and the
DATA/DRAIN validation rules. Direct commitments are comparable, and
ancestor commitments inherit that property.

Finally a valid receipt contains a correct FINAL signer. That signer
sent FINAL only after a timer committed its named DATA, directly or
as an ancestor; committing from another receipt emits no FINAL.
Consequently every receipt-based commit, including Completion, traces
to a causally earlier direct timer commitment of that DATA or a
descendant. It adds an already protected prefix. This provenance
argument does not take the safety of receipt-based commits as its own
premise. A timer canceled by Completion contributes no additional
direct commitment, particularly no commitment of its unexpired DRAIN.
No retirement step removes a prior commitment.

The global finality rule still preserves all of $G$ at its canonical
Force boundary and rejects further old-version DATA. Retirement and
receipt relay neither select a retained anchor nor authorize an
append, so they do not affect those external transitions.
\end{proof}


The universal protecting-QC lemma from Part~I cannot simply be copied:
its proof requires matching votes from all correct members, whereas a
closed member keeps a proposal without voting. For example, with
three unit-weight members and $F=1$, let correct $v_0$ close at $c=0$
and correct $v_1$ receive closure at $\Delta$. A Byzantine leader gives
a valid proposal to $v_1$ at $0<s<\Delta$. $v_1$ votes, forwards and
later commits; $v_0$ keeps it without voting. If the Byzantine member
withholds its vote and FINAL, only $v_1$'s FINAL exists. Its weight is
not greater than $F$: neither a proposal QC nor a complete receipt need
ever form. Retirement preserves this local commitment and evidence
without calling it Completion. The Publish Layer's complete-result
premise does not repair the missing signature.

\paragraph{Force before the first ReadyBundle.}
Suppose $v_0$ enters at $e_0$ and relays its complete ReadyBundle, while
$v_1$ learns Force before that bundle arrives. $v_1$ immediately stores
its closed mark. If its snapshot is already known it multicasts closure at once;
otherwise the ReadyBundle supplies the historical snapshot by
$e_0+\Delta$ and triggers the same multicast before its entry refusal.
The complete close evidence therefore reaches $v_0$ by $e_0+2\Delta$,
regardless of $v_0$'s later direct Espresso retrieval. Control processing
wins equality with its initial leader wait. If a proposal is the first
packet to reveal the snapshot instead, \cref{prop:retirement-observable-close}
gives the required $c\le s+\Delta$ and the late-window argument blocks
later recovery. Program exit, a stale callback and restart cannot
remove this relay path. This rules out a silently absent member being
used to classify an unprotected timer as an early-window commitment.

\section{Node storage and source relay}
\label{app:storage}
\label{sec:source-step}
\label{sec:source-retention}
A node must keep the exact bytes of old public objects. The object set
$\mathcal O_q$ of \cref{alg:global-rule} only grows, so the rule may
append a DATA whose body or FINAL signatures arrived many blocks
earlier; a handoff or recovery may also refer to old objects. Keeping
only a hash is not enough. This appendix states what each regional member stores, how it starts storing, and how it relays new source data to its region.

\paragraph{The public store.}
Each correct regional member keeps a durable \tagg{PublicStore} shared
by its source reader, replay logic, Session Manager and engines. It stores content, not only cursors and indexes. The store
has the following distinct components.
\begin{table}[htbp]\centering\small
\begin{tabularx}{\linewidth}{@{}>{\raggedright\arraybackslash}p{.24\linewidth}Y@{}}\toprule
Component & Stored content\\\midrule
Raw Espresso source & Every acquired complete relevant namespace,
including rejected or malformed records; headers, finality/linkage
proofs, authenticated absence, and canonical source order.\\
Object set $\mathcal O_q$ & The block bodies and FINAL signatures
published through $J_q$, with their bytes, indexed by hash and session,
and the source block that carried each. An index without the bytes is
not a stored object.\\
L1 source objects & Every obtained finalized authority, control,
header and registry/completeness proof, with its authentication inputs;
pending controls as well as applied controls.\\
Regional witnesses & Full validation bundles actually used by correct
regional actions, committed DATA and receipt dependencies. Possession
does not put a private object into $\mathcal O_q$.\\
Prepared anchors & Exact executable continuation objects and their
validated provenance, separately keyed by anchor, configuration,
region and native format.\\
Replay state & Verified checkpoint,
$\mathcal X=(\mathrm{Out},S_\beta,\kappa)$ of \cref{alg:global-rule}, the replay subposition,
source-availability cursors, and typed pending fetches.
The source-availability cursor is distinct from the replay cursor.\\\bottomrule
\end{tabularx}
\end{table}

At each source-availability event, store the entire authenticated raw
namespace and its evidence before advancing the availability cursor,
releasing the input buffer, or making a relay depend on it. An object
that \cref{alg:global-rule} has not used, or never uses, is stored all
the same. Likewise store each acquired L1
object before advancing its authority-fetch cursor. Regional validation
stores the witness it used, so a later required complete send does not
reconstruct that witness from an archive. Correct FINAL and complete
receipt sends obey the full-bundle rule of \cref{sec:admission}; a
reference is used only when its receiver can resolve it within the
regional delay bound.

A new branch version, session closure, advancement of $G$, checkpoint
compaction and program exit do not delete these objects. Cancellation
removes a consumer, not the store's retention responsibility. Transfer
outstanding sends and ownership durably before releasing resources.
Restart restores the exact store and replay subposition; loss of required
bytes requires hydration before the prepared timing claim can resume.
These are protocol storage rules, not an additional storage service.

\paragraph{Finite initial hydration and pruning.}
\label{sec:source-hydration}
Choose a fixed authenticated source cut $q_0$ when beginning this
retention discipline. Hydrate the raw source interval still legally
referenceable after that cut and all exact objects designated by the
starting checkpoint. Because $\mathcal O_q$ is cumulative, this
includes relevant raw namespaces back to the authorized source origin,
including objects that were never appended. Storing only accepted
ancestry is insufficient. Also prepare the fixed anchor's provenance and the
authority and state dependencies of the assigned handoff.
Store every newly available namespace beyond $q_0$ while hydration
runs. Neither a new branch version nor a later source height restarts the fixed
old task or changes its target to a moving cut.

The interval and each admitted payload are finite. Eventual historical
retrieval gives a finite completion time for that fixed set; the anchor
and finite AUTH prefix have their separate eventual-availability
premises. Taking a finite maximum imposes no bound on arbitrary archive catch-up.
The selected-region preparation service is stated separately in
\cref{ass:ready-service}. It does not include all future
L1 objects that an arbitrary sidecar might name.

A replay checkpoint may compact computation without deleting
referenceable bytes. Deduplication and equivalent lossless storage are
allowed while exact content and source-occurrence mappings remain
available to the modeled handler. A hash or an only-eventually-readable
archive is not a replacement for those bytes. Pruning requires a
protocol proof that no future supported reference, checkpoint,
anchor, pending validation or publication duty can need the object.
Advancing $G$ or terminating its producing engine does not prove this.
\Cref{alg:global-rule} itself never appends an object of a closed
version, or a DATA at or below the height of $\operatorname{last}(G)$
that is not in $G$; other duties may still need such objects, so this
alone does not justify pruning them.
The present rule supplies no finite-storage bound.

\paragraph{What the rule reads from the store.}
\label{sec:record-completeness}
\Cref{alg:global-rule} reads only $\mathcal O_q$ and the L1 inputs of
\cref{sec:logical-carrier}. $\mathcal O_q$ consists of the block bodies
and FINAL signatures carried by the authenticated namespace content of
$J_1,\ldots,J_q$, together with the objects that a protocol-authorized
checkpoint establishes. A private cache, a newer private checkpoint, a
hash or name without its object, or an unresolved sender assertion adds
nothing to $\mathcal O_q$. L1 objects keep their own authentication
interface and need not first occur in Espresso.

For a candidate $B$ in Step~2, the rule needs exactly the body of $B$,
the bodies of a covering path to a root $D$, the FINAL signatures on
$D$, and the descriptor $\nu$ in the tail $\tagg{NEXT}(\nu)$ that
acceptance will bind to
$B$. A publisher-selected PROPOSE,
recovery $\tagg{AUTH}_{r,g}$ path, L1 reference proof sidecar or unbound
native claim is not an input; its optional verification never delays the
rule. The correct signer's full regional checks justify intrinsic input
validity by \cref{prop:final-kernel-soundness}; they do not turn a
private object or a sidecar assertion into a public input. Check typed
identities, actual hashes, the signed domain, the current pair and the
fixed snapshot before verifying FINAL signatures and applying
\tagg{CheckFINAL}. The outer fields of an EXPORT are not checked.

An absent required object or an invalid signature only means that the
candidate does not qualify yet. Other candidates are unaffected, and
since $\mathcal O_q$ only grows, the candidate can qualify in a later
block once the missing object is published. An old object stays in
$\mathcal O_q$ even if nothing was appended when it arrived. It never
grants append permission to a closed version. In contrast, a missing
complete namespace, applicable control prefix or $\tagg{AUTH}_{r',0}$
needed by an atomic update or terminal successor selection pauses the rule at that
operation. It is never zero
weight or an empty candidate set. A cache miss for an object already in
$\mathcal O_q$ requires retrieval and resumption at the same
subposition; the prepared retention invariant below rules out such an
Espresso byte miss in the timed cohort. Optional regional evidence is
stored and routed to its own engine; a complete genuine Quit keeps its
evidence-first handling.

Step~2 appends, persists and exposes the terminal DATA before its
terminal branch applies NextLive for the successor. A wait for successor
authority therefore cannot withhold that DATA. Packet preparation cannot
invert this order.

\paragraph{Immediate regional source relay.}
A \tagg{SOURCESTEP} is an authenticated message to an installed fixed
regional snapshot. It carries its chain/protocol domain and sender,
the source height and full header, complete relevant namespace or
absence, finality/linkage evidence and canonical record order. It
identifies the maintained verified predecessor checkpoint, supplies
applicable complete L1 controls and the available fixed-boundary
inputs, and supplies required exact public bytes not already resolvable
by the recipient under the complete-message contract. The
recipient never trusts a sender-selected new checkpoint.

On acquiring a complete raw source input, each correct member stores
and sends it in that availability event, before waiting for its own
\cref{alg:external} result. Add its regional objects to the store
without first fetching optional witness targets. This relay
need not know which region replay will select. Native preparation does
not delay it. Missing mandatory boundary evidence remains an explicit
wait and is supplied in a later complete dependency event. An initial
raw relay is not falsely described as a complete boundary operation
when that evidence is absent. Equal packets are deduplicated by source
position and commitment without suppressing the first correct relay.

A sender includes full required bytes unless recipient resolution is
already covered by the regional full-bundle contract. Its own cache hit
alone does not justify omitting bytes at a recipient. Acquisition of an
additional exact dependency immediately resumes its saved operation;
there is no independent polling tick or replay job.

\paragraph{Receive, commit and activate.}
The receiver authenticates and stores the input, then processes the
contiguous dependency-complete prefix from its maintained checkpoint.
Out-of-order source inputs wait for their predecessors. For each ready
block it executes \cref{alg:external}, that is, the steps of
\cref{alg:global-rule}. Each step records its change to the output and
its subposition, and exposes appended DATA before any later wait. A
duplicate cannot repeat a completed step, count or boundary. Local
receive times do not replace canonical timestamps.

The same complete event delivers the verified transition to the Session
Manager. When the exact anchor continuation and remaining full Ready
inputs are already prepared, the Session Manager checks the fixed handoff and
the Ready condition and performs initial Enter once. It
attaches the already executable object; it does not import a database
or reconstruct native state here. No old lock, vote or entry time is
inherited. A Force-created version still cannot export in its triggering
block.

On closure, persist the version-wide closed mark and control
routing obligations atomically. Every entry, fresh action and callback
checks that mark by its own branch version. Background scans, archive copies and
retirement completion are not prerequisites to gating the old version
or entering a valid new pair. A complete old-session context still
invokes the complete-close handler and multicast under
\cref{sec:retirement}, in every program state.
Unresolved snapshot context remains pending. The per-session
$c_v$, the $4\Delta$ drain and the evidence duties of \cref{sec:retirement} still apply; recording
the closed mark alone neither starts that drain clock nor permits exit.

\begin{proposition}[Stored public source bytes]
\label{prop:retained-public-bytes}
After fixed-cut hydration, a continuously online history-following
member that executes the retention rules keeps every legally
referenceable acquired Espresso object and its authentication evidence.
If $J_q$ first becomes final at real time $t_E(q)$, the required
Espresso bytes through $q$ are available to that member by $t_E(q)+T_E$, without a
new archive read for a previously acquired object.
\end{proposition}
\begin{proof}
Hydration establishes the claim through the fixed cut. At each later
availability event, insertion of the entire raw namespace precedes the
availability-cursor advance and any use by \cref{alg:global-rule}.
Earlier objects of $\mathcal O_q$ are already stored; new objects reside
in this namespace. An absent
object is not fabricated. No retry, checkpoint or program transition
removes referenceable content. Induction over source positions preserves
the claim. Every predecessor through $q$ became final no later than $t_E(q)$,
so the stated visibility premise supplies its new source bytes and
Espresso evidence by $t_E(q)+T_E$. Older hydrated bytes remain present.
This is a statement about Espresso bytes, not a deadline for new L1 authorities.
\end{proof}

\paragraph{Prepared regional processing.}
\label{sec:source-step-timing}
When source bytes, authority and exact executable continuation are
resident, the complete receive handler validates them, performs the
canonical public transition and checks the guarded initial entry.
No historical fetch or native import is silently charged to the
regional communication bound.

The complete-message contract bounds delivery and processing of its
declared inputs. Initial preparation has the separate scope and origin
of \cref{ass:ready-service}. Returning-node catch-up outside that scope
remains eventual. The source-availability cursor, canonical replay
cursor and local Ready event can differ; none supplies authority to
execute an unannounced Force.

\end{document}
